\documentclass[11pt]{article}
\usepackage[a4paper,margin=26mm]{geometry}
\usepackage[T1]{fontenc}
\usepackage{amsmath,amssymb,amsthm,booktabs,hyperref}
\hypersetup{colorlinks=true,urlcolor=blue,linkcolor=blue}
\newtheorem{theorem}{Theorem}
\newtheorem{lemma}[theorem]{Lemma}
\newtheorem{proposition}[theorem]{Proposition}
\newtheorem{corollary}[theorem]{Corollary}
\newcommand{\TV}{\mathrm{TV}}
\newcommand{\Unif}{\operatorname{Unif}}
\newcommand{\Law}{\mathcal L}
\title{CYCLON convergence:\\proved cases and progress toward a general theorem}
\author{Research working note}
\date{20 October 2026}
\begin{document}
\maketitle

\begin{abstract}
We specify the adopted CYCLON algorithm and assumptions first:
mandatory partner removal, protection of received identities,
minimum-age duplicate merging, and independent uniform activations.
For every fixed shuffle length, arbitrary network size and cache
capacity, and arbitrary finite initial integer ages, we prove
uniform access to a finite bounded-age family. A decreasing
lexicographic color potential supplies the general construction.
Consequently there are finitely many stationary regimes,
continuous-time full-state convergence to their possibly
initial-dependent mixture, and exponential stationary age tails.
For shuffle length one a constructive common atom additionally
proves unique stationarity and uniform local samples conditional
on availability, with an explicit uniform geometric bound.
At capacity one the stationary law is explicitly uniform on
zero-age rooted spanning trees. A second complete theorem
now covers cache capacity two and shuffle length two for
every network size. It combines a four-exchange fresh-copy
gain with an extremal residual-component argument to prove
that all closed classes reach the same canonical state.
For capacity and shuffle length both two, a quadratic-length
controlled path gives the explicit, pessimistic bound
$\exp(O(n^2\log n))\log(1/\varepsilon)$ on mixing.
For ordinary random executions in that case, we additionally
prove an expected cut-contact bound of $7n$ and temporal
connectivity within $O(n^2+n\log(1/\varepsilon))$ global
activations with high probability. On three peers, an exact
positivity projection and integer contraction certificate give
an $O(\log(1/\varepsilon))$ single-sample bound uniformly
over all initial ages. An arbitrary-population polynomial
sampling bound remains open.
A general two-entry theorem now covers every allowed
cache capacity and population size. Its outward-spread
argument combines a bounded count of young descriptors
with a numerical-age argument excluding infinite plateaux.
For fixed shuffle lengths at least three,
a common-regime or sufficient classwise-symmetry argument
is still needed for general conditional uniformity.
For two-entry shuffles we additionally prove uniform contact
across every cut, accessibility of a state reachable from a
prescribed root, and reduction of the common-atom problem to
accessibility of any universally cached identity.
At length three, for every allowed cache capacity and
population size, we prove holder-only uniform first
contact across every cut and show that a
universal zero-age anchor suffices for a common atom;
combining cooperative callbacks and an extremal component
argument now proves its accessibility and full uniform
sampling for $c\ge n-2$. Smaller capacities at length
three, and arbitrary higher fixed lengths, remain open.

For every fixed shuffle length at least three, a specified
uniform-victim completion admits full-cache, zero-age, strongly
connected initial states whose conditional local samples fail
to converge to uniform. Protecting received identities excludes
this construction and preserves connectivity. With that explicit
protection rule and a uniformly positive probability of singleton
shuffles, we prove full convergence for every maximum shuffle
length. This modified-protocol theorem is distinguished from
the remaining fixed-length question.
The negative construction concerns a different replacement
completion and does not apply to the adopted protected algorithm.
Age-dependent initial-lineage erasure and sharper two-entry
age-reset results for other conventions are retained separately.

Earlier results are retained: a capacity-one retention theorem and
counterexamples for particular scheduling and retention conventions.
Each statement specifies its assumptions. In particular, the new
positive theorems use independent uniform activations rather than the
fixed periodic phases in the original paper. A general fixed-length
theorem must therefore specify both the scheduler and the replacement
rule; no universal conclusion across all these completions is claimed.
The adopted target for further fixed-length work is the thesis
convention in Section~\ref{sec:thesis}: protect received identities
and keep the minimum age when duplicates are merged.
\end{abstract}

\section{Algorithm and assumptions for the proof target}
\label{sec:target}

This section specifies the \emph{fixed algorithm} whose general
convergence proof is the current target. It uses received-identity
protection and minimum-age duplicate merging, following
\cite[Figure 2.3, p.~19]{cyclonthesis}. Its transition kernel is
denoted $\widehat P_\ell$. Later sections also retain results for
other conventions; those statements specify their own kernels.

\subsection{System, initialization, and randomness}
\begin{enumerate}
\item Membership is a fixed set $V=\{1,\ldots,n\}$ with $n\ge2$.
All peers follow the algorithm. There are no joins, departures,
crashes, message losses, or external introductions of identities.
\item Every peer has the same cache capacity $c$ and fixed shuffle
length $\ell$, with $1\le\ell\le c\le n-1$. A cache $C_v$
contains at most $c$ descriptors $u^a$, where $u\ne v$ and
$a\in\mathbb Z_{\ge0}$. No cache contains two descriptors for
the same identity. Empty and underfull caches are allowed.
\item Initial ages are arbitrary finite nonnegative integers,
without any common a priori bound. The initial directed overlay
has an edge $v\to u$ when $u\in C_v$; its underlying undirected
graph is connected. Strong connectivity and initially full caches
are not assumed.
\item At each global event choose an initiator independently and
uniformly from $V$, including peers with empty caches. Equivalently,
independent rate-one Poisson clocks at the peers generate this
embedded event chain. Positive results for this scheduler do not
automatically cover fixed periodic phases.
\item An exchange is reliable and atomic. Both buffers are formed
before either cache processes incoming entries. No other exchange
intervenes. Oldest-age ties, extra subsets, reply subsets, and
eligible victim sets are chosen uniformly. Conditional on the
current state and earlier choices in that exchange, the random
choices are independent, including victim choices at its two
endpoints. Fresh choices are used at subsequent events.
\end{enumerate}

\subsection{One activation and exchange}
For a packet or cache $A$, write $\operatorname{ids}(A)$ for
its identity set. Removing an identity removes its descriptor.
An event with initiator $p$ is the following procedure.
\begin{enumerate}
\item If $C_p$ is empty, do nothing and end the event.
Otherwise increment the age of every descriptor in $C_p$ by one.
\item Choose a maximum-age identity $q\in C_p$, uniformly among
ties. Choose a uniform subset $S$ of
$\min(\ell-1,|C_p|-1)$ other descriptors from $C_p\setminus\{q\}$.
The outgoing request packet is
\[
 B_p=\{p^0\}\cup S.
\]
Thus the partner's descriptor is replaced in the packet by a
fresh self descriptor, while the extra descriptors keep their
incremented ages.
\item Peer $q$ forms its reply $B_q$ as a uniform subset of
$\min(\ell,|C_q|)$ descriptors from its pre-exchange cache.
It does not increment its ages when acting as responder.
\item Remove the contacted partner $q$ unconditionally and
apply the merge procedure below at both endpoints, using the
already formed packets:
\[
\begin{split}
 C_p'&=\operatorname{Merge}_p
       (C_p\setminus\{q\},B_q,\operatorname{ids}(S)),\\
 C_q'&=\operatorname{Merge}_q
       (C_q,B_p,\operatorname{ids}(B_q)).
\end{split}
\]
Here the right-hand $C_p$ has the incremented ages from step~1;
the displayed set difference performs the mandatory removal.
All other caches
remain unchanged.
\end{enumerate}

\paragraph{Merge procedure.}
$\operatorname{Merge}_v(A,R,E)$ receives a resident cache $A$,
an incoming packet $R$, and the identity set $E$ of sent entries
eligible for replacement. Compute
\[
 D=\operatorname{ids}(R)\setminus\{v\},\qquad
 N=D\setminus\operatorname{ids}(A),\qquad
 d=\bigl(|N|-(c-|A|)\bigr)_+ .
\]
Ignore the incoming self descriptor, if any. For every identity
in $D\cap\operatorname{ids}(A)$ retain one descriptor with
age equal to the minimum of its resident and incoming ages.
These received identities are protected from eviction.
Choose a uniform $d$-element subset of
$E\cap\operatorname{ids}(A)\setminus D$ as victims, delete
those residents, and insert all descriptors with identities
in $N$, carrying their received ages. Newly inserted entries
cannot be victims, and unsent residents cannot be victims.
Lemma~\ref{lem:receiveprotect} proves that enough eligible
victims always exist, including for underfull exchanges.

The two key rules are therefore explicit: a valid received
identity is present after merging, and a duplicate keeps the
younger age. Sending an extra descriptor alone does not delete
its resident copy; it is deleted only if selected as an eligible
victim. The contacted partner is removed unconditionally.

\subsection{Convergence statement sought}
At a deterministic observation time, an observer $i$ requesting
$b\le c$ peers samples $b$ distinct identities uniformly without
replacement from its cache and orders them uniformly.
The target is positive limiting availability and
\[
 \left\|\Law_x\!\left(S_i^{(b)}(t)\mid |C_i(t)|\ge b\right)
 -\Unif\bigl((V\setminus\{i\})_b\bigr)\right\|_{\TV}
 \longrightarrow0
\]
for every allowed initial state $x$, observer $i$, and request
size $b$. The sample rule is independent of descriptor ages.
Waiting or retrying until a cache can answer is a different
observation rule. Independence of samples across times or peers,
uniformity of the full overlay graph, and a useful mixing-time
bound are additional properties, not assumed consequences.

A sufficient route is to prove a unique stationary law and
full-state convergence, including ages, together with positive
stationary availability. Relabelling symmetry then gives the
local uniform law. These are desired conclusions, not model
assumptions. Theorem~\ref{thm:thesisone} proves the full target
for $\ell=1$, and Theorem~\ref{thm:thesisc2} proves it
for $c=\ell=2$ at every network size.
Theorem~\ref{thm:thesisgeneraltwo} proves it for $\ell=2$
at every allowed cache capacity and population size.
Theorem~\ref{thm:dense3uniform} additionally proves it for
$\ell=3$ at every allowed capacity $c\ge n-2$.
Theorem~\ref{thm:thesisgeneralreset} and
Corollary~\ref{cor:thesisgeneralregimes} prove finite-age
recurrence and continuous-time convergence to finitely many
possible regimes for every fixed $\ell\ge2$; outside
the complete cases just listed, uniqueness and conditional
local uniformity are the remaining general step.
Shuffle length remains fixed throughout
the main target; occasional singleton shuffles are used only
in separately identified variant theorems.

\begin{center}
\begin{tabular}{lll}
\toprule
Shuffle length & Cache capacity & Conditional uniform sampling\\
\midrule
$\ell=1$ & any allowed $c$ & proved\\
$\ell=2$ & any allowed $c$ & proved\\
$\ell=3$ & $c\ge n-2$ & proved\\
$\ell=3$ & $3\le c<n-2$ & open\\
$4\le\ell\le c$ & any allowed $c$ & open\\
\bottomrule
\end{tabular}
\end{center}
All rows use the algorithm and assumptions above. Uniform
finite-age recurrence is proved even in the open sampling cases.

\section{Motivation and comparison with other conventions}

The motivating target is a local total-variation guarantee of the
kind proved for PEERSWAP in~\cite{peerswap}. Its interchange-process
reduction uses a fixed underlying graph and does not automatically
extend to partial cache exchanges.

We consider a static membership $V=\{1,\ldots,n\}$, reliable atomic
exchanges, and caches without self entries or duplicate identities.
Descriptors contain an identity and a nonnegative integer age.
Enhanced CYCLON increments the initiator's cached ages, selects an
oldest descriptor as partner, sends a fresh self descriptor in place
of that partner along with other selected descriptors, and receives
a subset of the partner's cache~\cite{cyclon}.

The update convention matters. We distinguish:
\begin{itemize}
\item \emph{Retention:} discard self and already-known incoming
descriptors, and replace sent cache entries only when a usable incoming
descriptor is inserted. A contacted partner is retained when no
replacement is usable. Discarded duplicates do not refresh ages.
\item \emph{Mandatory removal:} unconditionally remove the contacted
partner, even when the response yields no usable replacement.
\end{itemize}
These conventions differ on reciprocal exchanges. The published
paper's seven steps are followed by the explicit statement that the
contacted peer disappears~\cite[Section 2.2]{cyclon}. We therefore use
mandatory removal as the baseline for the new positive results. The
retention results below concern the separately specified alternative;
their obstructions do not establish failure of mandatory-removal
CYCLON. Underfull buffers, replacement choices, and scheduling are
specified explicitly wherever needed.

Time counts global activation events. Periodic phases and independent
uniform random activations are treated separately. For a cache that
can supply $b$ entries, $S_i^{(b)}$ is a uniformly selected, uniformly
ordered $b$-tuple of distinct cached identities. When a requested
tuple is unavailable, write $\bot$. Uniformity at one peer and time
does not imply independence across peers or times.

\section{A periodic counterexample under either convention}

At $c=\ell=1$, if the initiator is $p$, its neighbor is $q$, and
$q$'s cached identity is $r$, with $p,q,r$ distinct, the unique
exchange outcome is
\begin{equation}
 (f(p),f(q))=(q,r)\quad\longmapsto\quad
 (f'(p),f'(q))=(r,p).
 \label{eq:move}
\end{equation}
All other identities stay fixed. Both update conventions agree here.

\begin{theorem}[Failure with fixed periodic phases]
\label{thm:periodic}
There is a four-peer enhanced-CYCLON execution with a strongly connected
initial overlay, full caches, $c=\ell=1$, and one activation per peer
per round, in which a peer's sample never converges to uniform.
No exchange has a tie, unusable replacement, or empty cache.
\end{theorem}
\begin{proof}
Activate peers in order $1,2,3,4$ in every round, using distinct fixed
phases within a common period. Initialize all descriptor ages to zero
and the identity state to $(3,4,2,1)$. Equation~\eqref{eq:move} gives
\begin{center}
\begin{tabular}{lrrrr}
\toprule
After activation & $f(1)$ & $f(2)$ & $f(3)$ & $f(4)$\\
\midrule
Initial & 3 & 4 & 2 & 1\\
1 & 2 & 4 & 1 & 1\\
2 & 2 & 1 & 1 & 2\\
3 & 3 & 1 & 2 & 2\\
4 & 3 & 4 & 2 & 1\\
\bottomrule
\end{tabular}
\end{center}
The initial graph is the directed cycle $1\to3\to2\to4\to1$.
In every exchange $p,q,r$ are distinct. The initiator's sole descriptor
is incremented from zero to one, then discarded and replaced in its
outgoing message by its own age-zero descriptor. The responder's
age-zero descriptor moves to the initiator. Thus all stored ages are
again zero after each exchange. The entire state and scheduler phase
repeat after four activations.

At each round boundary, peer $1$ returns peer $3$ deterministically.
For all $k\ge0$,
\[
 \left\|\Law(S_1^{(1)}(4k))-\Unif\{2,3,4\}\right\|_{\TV}=\frac23.
\]
Moreover peer $1$ never stores peer $4$, even between round boundaries.
Every mixture over observation times therefore remains at least
$1/3$ from uniform. Time averaging cannot repair this execution.
\end{proof}

\section{Frozen multi-entry states with random scheduling}

This section uses retention and discards already-known descriptors
\emph{without refreshing their ages}. It allows random initiators,
random subsets, and any tie-breaking rule.

\begin{theorem}[Strongly connected frozen states]
\label{thm:twins}
For every even $n\ge4$, every $2\le c\le n-2$, and every
$1\le\ell\le c$, there is a strongly connected, full-cache initial
state for which all cache identities remain fixed forever under
enhanced CYCLON with retention. Every peer always has a unique oldest
descriptor. This holds for every activation sequence and every outcome
of the protocol's random subset selections.

For every $1\le b\le c$ and every observer $i$, its sample distance is
the positive constant
\begin{equation}
 \left\|\Law(S_i^{(b)}(t))
 -\Unif((V\setminus\{i\})_b)\right\|_{\TV}
 =1-\frac{(c)_b}{(n-1)_b},
 \qquad t\ge0,
 \label{eq:frozenTV}
\end{equation}
where $(a)_b=a(a-1)\cdots(a-b+1)$.
\end{theorem}
\begin{proof}
Partition peers into pairs $(p_k,q_k)$, $1\le k\le m=n/2$.
For each pair choose a set $S_k$ of $c-1$ identities outside the pair
that contains $p_{k+1}$, with indices modulo $m$.
This is possible since $m\ge2$ and $c-1\le n-3$.
Set
\[
 C_{p_k}=\{q_k\}\cup S_k,\qquad
 C_{q_k}=\{p_k\}\cup S_k.
\]
Give the reciprocal partner descriptor age $1$ and every other
descriptor age $0$. Each pair is internally mutually reachable;
the edges to $p_{k+1}$ connect the pairs around a directed ring.
Consequently the overlay is strongly connected.

Suppose $p_k$ initiates. Its unique oldest descriptor is $q_k$.
The outgoing message consists of its own fresh identity $p_k$ and
$\ell-1$ elements of $S_k$, all already in $q_k$'s cache. The responder
therefore discards every incoming descriptor.
Every descriptor it returns is either $p_k$, which is a self entry
at the initiator, or an element of $S_k$, already known there.
The initiator also discards every incoming descriptor. No cache
identity and no stored age is overwritten.

Incrementing all ages at the initiating cache preserves its strict
age difference between the partner and the common neighbors. Passive
responses do not increment ages. The same argument applies when
$q_k$ initiates. Induction proves that the identities remain frozen
and the designated partners remain uniquely oldest for all time.

The output is uniform on the $(c)_b$ ordered tuples contained in the
fixed cache, whereas the target is uniform on all $(n-1)_b$ tuples.
For uniform laws on a subset $A$ of a finite set $B$, their
total-variation distance is $1-|A|/|B|$, proving~\eqref{eq:frozenTV}.
\end{proof}

The smallest example, with $n=4,c=2$, is
\[
\begin{array}{ll}
C_1=\{2^{[1]},3^{[0]}\},&C_2=\{1^{[1]},3^{[0]}\},\\
C_3=\{4^{[1]},1^{[0]}\},&C_4=\{3^{[1]},1^{[0]}\},
\end{array}
\]
where brackets denote ages. It works for both $\ell=1$ and $\ell=2$.
Peer $1$ never samples peer $4$; its singleton distance is $1/3$.
For a realistic illustration, $c=20,\ell=5$ is covered whenever
even $n\ge22$.

This counterexample uses a prescribed nonzero-age initialization.
It does not apply to mandatory partner removal or duplicate-processing
rules that refresh stored ages. The next proposition addresses
reachability from zero ages for one parameter choice, while carefully
distinguishing pathwise and distributional statements.

\begin{proposition}[A frozen state reachable from zero ages]
\label{prop:zero}
Under retention, $n=4,c=2,\ell=1$, independent uniform initiators,
uniform age-tie resolution, and uniform responder subsets, there is
a strongly connected full-cache, zero-age initialization from which
a frozen identity state is reached with probability at least $2^{-24}$.
Consequently almost-sure time-averaged uniform sampling fails for that
initialization.
\end{proposition}
\begin{proof}
Start from
\[
C_1=\{2^0,4^0\},\quad C_2=\{3^0,4^0\},\quad
C_3=\{1^0,2^0\},\quad C_4=\{1^0,2^0\}.
\]
Superscripts denote ages. This graph is strongly connected.
Execute the following triples of initiator, selected partner, and
returned identity:
\[
(2,4,1),\ (1,4,2),\ (3,1,2),\
(3,1,4),\ (4,2,1),\ (4,2,3).
\]
For $\ell=1$, an accepted incoming descriptor replaces the single
sent entry; a filtered descriptor makes no replacement. Direct
substitution in these rules gives
\[
\begin{array}{ll}
C_1=\{4^1,3^0\},& C_2=\{3^1,4^0\},\\
C_3=\{2^2,4^1\},& C_4=\{1^2,3^1\}.
\end{array}
\]
The pairs $(1,4)$ and $(2,3)$ now satisfy the frozen-state condition,
with common sets $\{3\}$ and $\{4\}$ respectively.
Each of the six specified initiators has probability $1/4$,
each specified oldest-partner choice is between two tied entries,
and each specified response has probability $1/2$.
The exact prefix probability is therefore $(1/16)^6=2^{-24}$.
On this event, peer $1$ never again samples peer $2$, so its empirical
frequency for that identity tends to zero rather than $1/3$.
\end{proof}

Positive-probability absorption by itself does \emph{not} prove that
the unconditional local marginal law fails to approach uniform from
the zero-age initialization: different frozen outcomes could balance
in their mixture. Theorem~\ref{thm:twins} directly disproves that
marginal-law guarantee when the specified nonzero initial ages are
allowed.

\section{An availability obstruction under mandatory removal}

Now choose initiators independently and uniformly, remove the partner
unconditionally, and let an activation of an empty cache do nothing.

\begin{proposition}[Persistent failure of full-capacity requests]
\label{prop:availability}
Let $c=n-1$ and $1\le\ell\le c$. For every peer $i$ and $t\ge1$,
even starting from complete caches,
\[
 \Pr\{S_i^{(c)}(t)=\bot\}\ge\frac1n,\qquad
 \left\|\Law(S_i^{(c)}(t))
       -\Unif((V\setminus\{i\})_c)\right\|_{\TV}\ge\frac1n.
\]
\end{proposition}
\begin{proof}
If activation $t$ selects $i$, either its cache is already empty or
it removes some partner $q$. Since $q$'s cache contains no self entry
$q$, its response cannot restore $q$ to $i$. Thus $i$ has at most
$n-2<c$ entries at the end of that activation. The event that $i$ was
selected has probability $1/n$. The distance bound follows by testing
the event $\{\bot\}$, which has target probability zero.
\end{proof}

This concerns unconditional $b=c$ requests. It does not establish
nonuniformity for smaller samples conditional on availability.
Retrying or conditioning changes the sampling model.

\section{A positive capacity-one theorem}

\subsection{Precisely specified variant}
Set $c=\ell=1$, retain the current neighbor when a reply is unusable,
and choose the initiator independently and uniformly at each step.
With $q=f(p)$ and $r=f(q)$, the complete identity update is
\begin{equation}
 T_p f=
 \begin{cases}
 f,&r=p,\\
 f\text{ with }f(p)\gets r,\ f(q)\gets p,&r\ne p.
 \end{cases}
 \label{eq:kernel}
\end{equation}
Initial ages may be arbitrary: the single cached entry is necessarily
oldest, so ages have no effect on this identity projection.
Let $\Omega_n$ be the nonself functions $f:V\to V$ whose underlying
undirected graph is connected, and let
\[
 P(f,g)=\frac1n\#\{p:T_p f=g\}.
\]

\begin{theorem}[Convergence of the capacity-one retention variant]
\label{thm:positive}
For every $n\ge2$, $f\in\Omega_n$, and observer $i$,
\[
 \Law_f(f_t(i))\longrightarrow\Unif(V\setminus\{i\}).
\]
For $n\ge4$, with $D=2n(n-2)$ and $\delta=n^{-D}$,
\begin{equation}
 \left\|\Law_f(f_t(i))-\Unif(V\setminus\{i\})\right\|_{\TV}
 \le(1-\delta)^{\lfloor t/D\rfloor}.
 \label{eq:positive}
\end{equation}
Consequently $D\lceil n^D\log(1/\varepsilon)\rceil$ activations
suffice for distance at most $\varepsilon$.
For $n=3$, $(8/9)^{\lfloor t/2\rfloor}$ is sufficient.
For $n=2$, the sample is already uniform.
\end{theorem}
This is an extremely weak quantitative bound and gives no efficient
mixing claim. The theorem concerns identity marginals, not stationarity
of the complete age process.

\subsection{Proof for \texorpdfstring{$n\ge4$}{n >= 4}}

\begin{lemma}[Inverse moves and connectivity]
\label{lem:inverse}
$\Omega_n$ is closed under $T_p$. Every nontrivial move initiated
by $p$ with $q=f(p)$ is inverted by initiating $q$ in the resulting
state. In particular $P$ is symmetric.
\end{lemma}
\begin{proof}
A nontrivial move replaces a path on the three distinct vertices
$p,q,r$ by another path on the same vertices, preserving undirected
connectivity and avoiding self entries. In the resulting state
$g(q)=p,g(p)=r$; initiating $q$ restores $f(q)=r,f(p)=q$.
This pairs each marked transition $(f,p,g)$ with its unique reverse
marked transition $(g,q,f)$, so the transition counts agree.
\end{proof}

\begin{lemma}[Insertion into a cycle]
\label{lem:insert}
Every state in $\Omega_n$ reaches a directed Hamiltonian cycle in at
most $n-2$ moves.
\end{lemma}
\begin{proof}
A connected functional digraph has exactly one directed cycle:
outgoing paths eventually reach a cycle, and an edge cannot join
basins of distinct cycles. The cycle length is at least two.
While an off-cycle vertex remains, choose $p$ whose parent $q=f(p)$
lies on the cycle. Initiating $p$ replaces its cycle edge $q\to r$
by $q\to p\to r$, adding one vertex to the cycle.
\end{proof}

\begin{lemma}[Adjacent swap]
\label{lem:swap}
If a Hamiltonian cycle has consecutive distinct vertices $x,a,b,y$,
the activation word $x,x,a,b$ swaps $a,b$, leaving every other
outgoing edge unchanged.
\end{lemma}
\begin{proof}
The only changing outgoing edges are
\[
\begin{array}{c|ccc}
 & f(x)&f(a)&f(b)\\ \hline
\text{initial}&a&b&y\\
x&b&x&y\\
x,x&y&x&x\\
x,x,a&a&y&x\\
x,x,a,b&b&y&a
\end{array}
\]
All moves are valid for $n\ge4$. At $n=4$ the temporary residual
cycle $x\leftrightarrow y$ is allowed.
\end{proof}

\begin{lemma}[Diameter]
\label{lem:diameter}
For $n\ge4$, the transition graph on $\Omega_n$ is connected with
diameter at most $D=2n(n-2)$.
\end{lemma}
\begin{proof}
Reach a Hamiltonian cycle from each of two states in at most $n-2$
steps. Anchor the cyclic notation at one fixed vertex. The other
$n-1$ vertices form a linear ordering; any two orders are related
by at most $(n-1)(n-2)/2$ adjacent swaps. Each costs four moves by
Lemma~\ref{lem:swap}. Reverse the second insertion path using
Lemma~\ref{lem:inverse}. The total is at most
$2(n-2)+2(n-1)(n-2)=2n(n-2)$.
\end{proof}

\begin{proof}[Proof of Theorem~\ref{thm:positive} for $n\ge4$]
The uniform law $\pi_n$ on $\Omega_n$ is stationary because $P$ is
symmetric. Choose the state
$z(1)=2,z(2)=1,z(j)=1$ for $j\ge3$. Initiating $1$ at $z$ is a
self-loop. Every state reaches $z$ along a path of at most $D$ moves.
Pad shorter paths with that self-loop to exactly $D$ activations.
The prescribed sequence has probability $n^{-D}=\delta$, so
\[
 P^D(f,\cdot)\ge\delta\,\delta_z(\cdot)\quad(f\in\Omega_n).
\]
Writing $P^D(f,\cdot)=\delta\,\delta_z+(1-\delta)R(f,\cdot)$ yields
\[
 \|\mu P^D-\nu P^D\|_{\TV}
 \le(1-\delta)\|\mu-\nu\|_{\TV}
\]
for any probability laws $\mu,\nu$. Iterate with $\nu=\pi_n$,
then use total-variation contraction for the remaining steps:
\[
 \|P^t(f,\cdot)-\pi_n\|_{\TV}
 \le(1-\delta)^{\lfloor t/D\rfloor}.
\]
Relabelling bijects $\Omega_n$ onto itself and preserves $\pi_n$.
Permutations fixing $i$ act transitively on the other identities, so
$\pi_n\{f:f(i)=j\}=1/(n-1)$ for each $j\ne i$.
Projection gives~\eqref{eq:positive}. The time bound follows from
$(1-\delta)^k\le e^{-\delta k}$.
\end{proof}

\subsection{Small network sizes}
For $n=2$ the sole state is the reciprocal pair and its local target
distribution has one element. For $n=3$ the identity chain has two
four-state classes. One consists of
\[
 H=(2,3,1),\quad L_1=(3,1,1),\quad
 L_2=(2,1,2),\quad L_3=(3,3,2).
\]
The other is obtained by relabelling a reversed triangle.
The center $H$ moves to each leaf with probability $1/3$.
Each leaf returns to $H$ with probability $1/3$ and stays with
probability $2/3$. These classes contain all eight nonself maps.
Uniform law on each class is stationary, and every observer points
to either other identity in exactly two of its four states.

Choose a leaf $L_\star$ in the appropriate class. From every state,
$P^2(f,L_\star)\ge1/9$: from another leaf go through $H$;
from $H$ use the target leaf's self-loop; from the target leaf use
two self-loops. The same minorization argument proves the stated
bound. This also shows why irreducibility of the full identity
space is stronger than necessary for local uniformity.

\section{A general symmetry reduction, not a mixing proof}

\begin{lemma}[Symmetrized reference process]
\label{lem:symmetry}
Let a Markov transition kernel commute with peer relabelling, and
suppose the sample extraction rule is equivariant and always defined.
Fix initial state $x$, observer $i$, and the group $H_i$ of
permutations fixing $i$. Start a reference process $Y_t$ from
\[
 \bar\mu_x=\frac1{|H_i|}\sum_{\sigma\in H_i}\delta_{\sigma x}.
\]
Then $S_i^{(b)}(Y_t)$ is exactly uniform on the ordered distinct
$b$-tuples of other peers at every time. For any coupling,
\[
 \|\Law_x(S_i^{(b)}(X_t))-\Unif((V\setminus\{i\})_b)\|_{\TV}
 \le\Pr\{S_i^{(b)}(X_t)\ne S_i^{(b)}(Y_t)\}.
\]
\end{lemma}
\begin{proof}
The initial law of $Y$ is invariant under $H_i$. Equivariance
preserves this invariance at every time and transfers it to the
sample law. Since $H_i$ acts transitively on the output tuples,
that law must be uniform. The remaining claim is the coupling
inequality.
\end{proof}

This lemma requires no finite state space, stationary age law,
or full-state irreducibility. If samples can be unavailable, the
uniformity conclusion holds only conditional on availability.
The lemma provides no forgetting bound by itself. In
Theorem~\ref{thm:twins}, its reference law is uniform while the original
local law remains permanently nonuniform, so the required disagreement
probability cannot tend to zero.

\section{Exact model for the mandatory-removal results}
\label{sec:mandatorymodel}

The positive results that follow use the enhanced exchange in the
published CYCLON paper~\cite[Section 2.2]{cyclon}, with the following
explicit completion of its choices. Membership is fixed, exchanges
are reliable and atomic, and $1\le\ell\le c\le n-1$. Each cache
$C_p$ contains at most $c$ distinct nonself identities, each with a
finite nonnegative integer age.

At each global event, choose an initiator independently and uniformly
from $V$. An empty initiator does nothing. Otherwise, the initiator
$p$ increments all its cached ages and selects a maximum-age partner
$q$, breaking ties uniformly. It selects
$\min(\ell-1,|C_p|-1)$ other entries uniformly. Write $S$ for these
other entries with their incremented ages. Its outgoing buffer is
$\{p^0\}\cup S$, and it unconditionally removes $q$ from its cache.
The responder chooses a uniform subset of size
$\min(\ell,|C_q|)$ of its \emph{pre-exchange} cache and returns it.
Both sides discard incoming self entries and incoming identities
already present locally, without refreshing duplicate ages. They
insert every remaining incoming entry, using vacancies first and
then replacing sent entries still resident locally. Every eligible
choice of victims has positive probability, independent of identities.
For uniform quantitative bounds one may choose uniformly among the
eligible sets. Entries not selected as victims remain in place.

The two buffers are formed before either cache processes incoming
entries. Removing $q$ creates a vacancy at $p$; it does not remove
the other sent entries. At $q$, all sent entries are initially
resident. These rules always provide enough eligible slots. If the
pre-exchange cache sizes are $a,b$, the initiator has
$c-a+1+\min(\ell-1,a-1)\ge\ell$ vacancies plus eligible residents,
and the responder has $c-b+\min(\ell,b)\ge\ell$.
The incoming buffer size at either side is at most $\ell$.

This defines a countable-state Markov chain including ages.
Independent equal-rate Poisson activation clocks give this same
embedded chain. Fixed periodic phases, which the original paper
uses, give a different scheduler. Thus none of the positive results
below silently asserts convergence for every periodic schedule.
The underfull-buffer convention and uniform tie-breaking above are
also mathematical specifications, rather than claims that the paper
uniquely prescribes them.

\section{From full-state convergence to local uniformity}
\label{sec:stationary}

The following facts isolate what a general proof must establish.
A Markov kernel is \emph{equivariant} if relabelling peers in a state
and then taking a step has the same law as taking a step and
relabelling its outcome.

\begin{lemma}[A common atom gives convergence, including ages]
\label{lem:atom}
Let $P$ be a Markov kernel on a countable state space. Suppose there
are $L\ge1$, $\delta>0$, and a state $z$ such that
\[
 P^L(x,\{z\})\ge\delta\qquad\text{for every state }x.
\]
Then $P$ has a unique stationary probability law $\pi$, and
\[
 \sup_x\|P^t(x,\cdot)-\pi\|_{\TV}
 \le(1-\delta)^{\lfloor t/L\rfloor}.
\]
Irreducibility of the whole state space is not required.
\end{lemma}
\begin{proof}
The case $\delta=1$ is immediate. Otherwise split
$P^L(x,\cdot)=\delta\delta_z+(1-\delta)R(x,\cdot)$, where $R$ is a
Markov kernel. Hence $\mu\mapsto\mu P^L$ is a contraction by
$1-\delta$ in total variation. Probability measures on a countable
space are complete in this metric, so the contraction has a unique
fixed point $\pi$. Since
$(\pi P)P^L=(\pi P^L)P=\pi P$, uniqueness gives $\pi P=\pi$.
Every $P$-stationary law is also $P^L$-stationary, proving uniqueness
for $P$. Iterating the contraction and applying the remaining
$t\bmod L$ transitions proves the bound.
\end{proof}

\begin{theorem}[Stationary symmetry and availability]
\label{thm:stationary}
Let $P$ be equivariant on a relabelling-invariant state space.
Suppose it has a unique stationary law $\pi$ and
$\eta_t(x)=\|P^t(x,\cdot)-\pi\|_{\TV}\to0$.
Fix $i$ and $1\le b\le c$, and suppose
$a=\pi\{|C_i|\ge b\}>0$. For uniform sampling without replacement
from $C_i$, followed by uniform ordering,
\[
 \Law_\pi(S_i^{(b)}\mid |C_i|\ge b)
 =\Unif((V\setminus\{i\})_b).
\]
The same conditional sample law converges to this uniform law from
every initial state $x$. If $\eta_t(x)\le a/2$, its distance is at most
$4\eta_t(x)/a$.
\end{theorem}
\begin{proof}
For every permutation $\sigma$ of $V$, equivariance makes
$\sigma_*\pi$ stationary. Uniqueness gives $\sigma_*\pi=\pi$.
The event $A=\{|C_i|\ge b\}$ is invariant under permutations fixing
$i$, and the output rule is equivariant under this group. Its action
on distinct ordered $b$-tuples of other identities is transitive.
An invariant probability law on that finite set is uniform.

For the quantitative claim include the extra randomness of sample
extraction in the transition from a state to $(A,S_i^{(b)})$.
Total variation cannot increase under this transition. Put
$a_t=\Pr_x(A\text{ at }t)$. Then $|a_t-a|\le\eta_t(x)$.
For any set $B$ of available tuples,
\[
 \left|\frac{\Pr_x(A,S_i^{(b)}\in B)}{a_t}
       -\frac{\Pr_\pi(A,S_i^{(b)}\in B)}a\right|
 \le\frac{2\eta_t(x)}{a_t}
 \le\frac{4\eta_t(x)}a .
\]
Taking the supremum over $B$ proves the result.
\end{proof}

At stationarity the \emph{unconditional} output law, including the
unavailable symbol, is
$(1-a)\delta_\bot+a\Unif((V\setminus\{i\})_b)$.
Its distance from the always-available uniform target is exactly
$1-a$. Conditioning at a fixed observation time is the sampling
rule in Theorem~\ref{thm:stationary}. Waiting or retrying until a cache
can answer introduces a stopping-time rule and requires a separate
argument. Nor does this theorem assert independent samples across
peers or time, or a uniform distribution on global overlay graphs.

\section{A general age-reset result for shuffle length one}
\label{sec:agereset}

\begin{theorem}[Finite-set regeneration for $\ell=1$]
\label{thm:agereset}
Use the mandatory-removal model of Section~\ref{sec:mandatorymodel}
with arbitrary $n,c$ and $\ell=1$. Label the peers $1,\ldots,n$ and set
\[
 K=n+c-1,\qquad L=nK,\qquad \delta=(nc)^{-L}.
\]
Let $\mathcal F$ be the finite set of states with all ages zero and
$C_i\subseteq\{i+1,\ldots,n\}$ for every $i$. Then, from every state
with arbitrary finite ages,
\[
 P^L(x,\mathcal F)\ge\delta.
\]
If $\tau_{\mathcal F}$ is its first hitting time, then
\[
 \Pr_x\{\tau_{\mathcal F}>mL\}\le(1-\delta)^m,\qquad
 \mathbb E_x\tau_{\mathcal F}\le L/\delta .
\]
\end{theorem}
\begin{proof}
First activate one fixed peer $p$ for $K$ consecutive events.
Only its cache ages are incremented during this block. When its
chosen partner $q$ already knows $p$, prescribe that $q$ return the
descriptor for $p$. This occurs with probability at least $1/c$.
The reply is a self entry at $p$, so mandatory removal reduces
$|C_p|$ by one. The request contains only $p^0$, a duplicate at
$q$; hence $q$ retains its descriptor for $p$.

If $q$ does not know $p$, its receipt of $p^0$ makes it learn $p$.
It either uses a vacancy or replaces its single returned entry.
Subsequent exchanges in this block cannot make it forget $p$:
it never initiates, and every request from $p$ contains only the
already-known identity $p$. There are at most $n-1$ such first-learning
events. The size of $C_p$ never increases, because one descriptor is
removed and at most one is received. Every other nonempty activation
in the prescribed block shrinks it. At most $c$ shrinkings empty it.
Consequently $n-1+c=K$ activations suffice, independently of all ages.
Remaining activations of an empty $p$ do nothing.

Now execute this block for $p=1,2,\ldots,n$ in order. Immediately
after its block, $p$ is empty. During later blocks it never initiates,
and every incoming request contains only a fresh self descriptor
$j^0$ from a later peer $j>p$. Thus its final cache consists solely
of age-zero descriptors of later peers. The final state belongs
to $\mathcal F$.

Prescribing the $L$ initiators has probability $n^{-L}$. Each reply
that must be prescribed has conditional probability at least $1/c$;
there are at most $L$ of them. No particular oldest-partner or
victim choice was required. The resulting lower bound is
$(nc)^{-L}$. Applying it to successive $L$-event blocks, using the
Markov property, gives the tail bound, and summing that bound gives
the expectation.
\end{proof}

This proves that unbounded initial ages can be forgotten often
enough to hit a finite family. It does not prove that the different
members of $\mathcal F$ communicate or have identical limiting
local marginals. In particular, $\mathcal F$ need not have a unique
empty sink. Weak connectivity is preserved for $\ell=1$: the
contacted edge is retained in reverse, and if the responder deletes
its sole returned edge, that edge is transferred to the initiator
or is already present there. Hence the same hitting statement
holds for $\mathcal F$ intersected with the weakly connected states
when that is the initial state space.
The common-atom construction in Section~\ref{sec:generalone}
will connect this family for arbitrary $n,c$. First we record a
general criterion showing what the age reset alone implies.

\subsection{Convergence without assuming a unique recurrent class}

\begin{lemma}[A finite reset family with self-loops]
\label{lem:finiteclasses}
Let $P$ be a countable-state Markov kernel. Suppose a finite
nonempty family $F$, an integer $L\ge1$, and $\delta>0$ satisfy
\[
 P^L(x,F)\ge\delta\quad\text{for every }x,
 \qquad P(f,\{f\})>0\quad\text{for every }f\in F.
\]
Then there are finitely many closed communicating classes
$C_1,\ldots,C_m$, with $m\le |F|$, each having a unique stationary
probability law $\pi_k$. Every initial state enters their union
almost surely. There exist $B\ge1$ and $0<\alpha<1$ such that,
writing $h_k(x)$ for the probability of eventually entering $C_k$,
\begin{equation}
 \sup_x\left\|P^t(x,\cdot)-
                 \sum_{k=1}^m h_k(x)\pi_k\right\|_{\TV}
 \le 2(1-\alpha)^{\lfloor t/(2B)\rfloor}.
 \label{eq:multiclass}
\end{equation}
The weights sum to one. This conclusion permits the limiting
stationary law to depend on the initial state.
\end{lemma}
\begin{proof}
On the finite vertex set $F$, draw $f\leadsto g$ when some
finite positive-probability path of the full chain goes from $f$
to $g$. Consider the terminal strongly connected components of
this reachability graph and choose an atom $z_k$ in each.
Every $f\in F$ reaches at least one such atom.

Let $C_k$ be the set of all full-chain states reachable from $z_k$.
If $y\in C_k$, the reset bound gives a path from $y$ to some
$g\in F$. Terminalness implies that $g$ and $z_k$ lie in the
same strongly connected component of the finite graph. Hence
$y$ can return to $z_k$. Thus $C_k$ is a communicating class,
and it is closed because every state reachable from one of its
states is reachable from $z_k$. Different terminal components
give disjoint classes.

For each $f\in F$, prescribe a finite positive-probability path to
one selected $z_{k(f)}$. Since $F$ is finite, the maximum path
length $M$ is finite. Pad all paths to length $M$ using the
self-loop at their terminal atoms. Let $\varepsilon>0$ be the
minimum probability of these finitely many padded paths.
Set $B=L+M$ and choose
$0<\alpha\le\min(\delta\varepsilon,1/2)$.
Then, with $C=\bigcup_k C_k$,
\[
 P^B(x,C)\ge\alpha\quad\text{for every }x.
\]
Moreover, if $x\in C_k$, its reset states lie in $F\cap C_k$,
and their prescribed paths necessarily lead to $z_k$. Therefore
\[
 P^B(x,\{z_k\})\ge\alpha\quad (x\in C_k).
\]
Lemma~\ref{lem:atom}, applied to the closed space $C_k$, gives
its unique stationary law and convergence bound
$(1-\alpha)^{\lfloor t/B\rfloor}$.
The first inequality and closedness give
$\Pr_x(\tau_C>s)\le(1-\alpha)^{\lfloor s/B\rfloor}$.
In particular $\sum_k h_k(x)=1$.

Split $t=s+r$, where $s=\lfloor t/2\rfloor$. Conditional on
$X_s\in C_k$, evolve for $r$ more steps and compare with $\pi_k$.
On $X_s\notin C$, compare the remaining evolution with the
mixture of its eventual class laws; this contribution is bounded
by the probability of that event. The Markov property gives
$\mathbb E_x h_k(X_s)=h_k(x)$, so the comparison law is exactly
$\sum_k h_k(x)\pi_k$. Hence the total distance is at most
\[
 (1-\alpha)^{\lfloor s/B\rfloor}
 +(1-\alpha)^{\lfloor r/B\rfloor}
 \le2(1-\alpha)^{\lfloor t/(2B)\rfloor}.
\]
All closed communicating classes must intersect $F$ by the reset
bound, and so are among the classes just constructed.
\end{proof}

\begin{theorem}[Arbitrary-capacity convergence for singleton shuffles]
\label{thm:elloneconvergence}
For every $n\ge2$ and $1\le c\le n-1$, mandatory-removal CYCLON
with $\ell=1$ and independent uniform activations converges in
full-state total variation from every weakly connected initial
state with arbitrary finite integer ages. Its limit has the form
$\sum_k h_k(x)\pi_k$ in Lemma~\ref{lem:finiteclasses}, and the
uniform geometric bound \eqref{eq:multiclass} holds for some
constants depending on $n,c$ and the specified transition rule.
\end{theorem}
\begin{proof}
Restrict to the closed weakly connected state space.
Theorem~\ref{thm:agereset} supplies its finite reset family
$F=\mathcal F\cap\{\text{weakly connected states}\}$ and a uniform
reset probability. Every $f\in F$ has $C_n=\varnothing$,
so activating $n$ gives a self-loop of probability $1/n$.
Apply Lemma~\ref{lem:finiteclasses}.
\end{proof}

This is a general convergence statement for the age-and-cache
process. By itself it does not prove forgetting of the initial
configuration or uniform local samples. The constants $B,\alpha$
above are existential: their construction uses finite paths between
the reset states, without an explicit bound on those path lengths.
To invoke Theorem~\ref{thm:stationary} one must still establish a
unique stationary regime, or prove that all possible limiting
regimes and their initial-state-dependent weights give the desired
uniform local law. The next two sections resolve this issue for
every capacity and provide explicit convergence bounds.

\section{A complete mandatory-removal theorem at capacity one}
\label{sec:mandatoryone}

\begin{theorem}[Every network size, capacity one]
\label{thm:mandatoryone}
Let $n\ge2$ and $c=\ell=1$ in
Section~\ref{sec:mandatorymodel}. On the weakly connected state
space, the full cache-and-age chain has a unique stationary law
$\pi$: the uniform law on all labelled rooted spanning trees,
with all descriptor ages zero and all arcs directed toward the
empty root. Set $D=2n(n-1)$ and $\beta=n^{-D}$. Then
\[
 \sup_x\|P^t(x,\cdot)-\pi\|_{\TV}
 \le(1-\beta)^{\lfloor t/D\rfloor}.
\]
At stationarity, every observer $i$ has
\[
 \pi(C_i=\varnothing)=\frac1n,\qquad
 \pi(C_i=\{j^0\})=\frac1n\quad(j\ne i).
\]
Thus the sample converges to uniform on the other peers conditional
on a nonempty cache, with limiting availability $1-1/n$.
\end{theorem}
\begin{proof}
Represent a singleton cache by its identity $f(p)$ and an empty
cache by $f(p)=\bot$. If $p$ is nonempty, put $q=f(p)$ and
$r=f(q)$. The identity update is
\[
 f'(q)=p,\qquad
 f'(p)=
 \begin{cases}
  \bot,& r\in\{p,\bot\},\\
  r,& r\notin\{p,\bot\}.
 \end{cases}
\]
Other pointers do not change. A path on three distinct vertices
$p\to q\to r$ becomes $q\to p\to r$. Contact with an empty
responder reverses the edge. A reciprocal exchange removes only
one of its two opposite arcs. Weak connectivity is therefore
preserved.

A weakly connected digraph of outdegree at most one consists
either of a directed cycle with trees feeding into it, or of a
rooted spanning tree with one empty root. Indeed, following outgoing
pointers reaches a cycle or an empty vertex, and different terminal
cycles or empty vertices cannot have intersecting weak basins.
On a directed cycle of length $k\ge3$, activating a cycle vertex
shortens that cycle by one. On a two-cycle, activating either
endpoint produces a rooted tree. At most $n-1$ prescribed
activations thus reach a rooted tree, which remains a rooted tree.

In a rooted tree the initiator's descriptor is deleted, the
responder's descriptor transfers to the initiator if present, and
the responder receives a fresh age-zero descriptor. There is no
copying, and the only incremented descriptor is deleted.
Consequently the number of positive-age descriptors decreases by
one whenever a positive-age carrier initiates, and never increases.
At most $n-1$ further prescribed activations clear every positive
age. The zero-age rooted-tree family is closed.

We next prove that its identity transition graph is connected.
Fix an anchor $a$ and follow its directed path to the root.
While an off-path vertex remains, choose one, $p$, whose parent
$q$ is on the path. Activating $p$ inserts $p$ immediately after
$q$ on the path; if $q$ is the root, $p$ becomes the new root.
Such a vertex exists by connectivity of the tree.
At most $n-1$ moves give a Hamiltonian path beginning at $a$.

On a Hamiltonian path with consecutive vertices $x,u,v$, write
$y=f(v)$, allowing $y=\bot$ when $v$ is the root.
The word $x,x,u,v$ swaps $u,v$ and changes no other outgoing edge:
\[
\begin{array}{c|ccc}
 &f(x)&f(u)&f(v)\\ \hline
 \text{initial}&u&v&y\\
 x&v&x&y\\
 x,x&y&x&x\\
 x,x,u&u&y&x\\
 x,x,u,v&v&y&u
\end{array}
\]
All four moves are valid, including when $y=\bot$.
Adjacent swaps sort the $n-1$ vertices following $a$, using at most
$(n-1)(n-2)/2$ swaps.
Every nontrivial tree move initiated by $p$ is inverted by
initiating its old parent $q$ in the resulting tree, including
when that parent was the empty root. Therefore we may reverse
the insertion path of a target tree. The tree transition graph
has diameter at most
\[
 2(n-1)+4\,\frac{(n-1)(n-2)}2=2(n-1)^2.
\]

Choose one fixed zero-age rooted star $z$.
The cycle reduction, age clearing, and tree diameter bound give
a path to $z$ of length at most
$2(n-1)+2(n-1)^2=D$ from every initial aged state.
Pad to exactly $D$ events using the empty root's no-op.
Every prescribed initiator has probability $1/n$; there are no
nontrivial buffer or tie choices at capacity one. Hence
$P^D(x,\{z\})\ge n^{-D}$.
Lemma~\ref{lem:atom} gives uniqueness and the displayed bound.

Pairing each marked tree transition $(f,p,g)$ with its inverse
$(g,q,f)$ proves symmetry of the identity kernel on the finite
tree family. Its uniform law, with all ages zero, is therefore
stationary and must equal $\pi$. Relabelling makes the root
uniform on $V$. Permutations fixing $i$ make the remaining
$n-1$ possible cached identities equiprobable, each with mass
$(1-1/n)/(n-1)=1/n$. Conditional convergence follows from
Theorem~\ref{thm:stationary}.
\end{proof}

\section{A complete arbitrary-capacity theorem for singleton shuffles}
\label{sec:generalone}

We now close the remaining common-regime gap for $\ell=1$.
For an ordered set $V=\{1,\ldots,n\}$, define $H(V,c)$ by
giving peer $i$ the $\min(c,n-i)$ greatest labels larger than $i$,
all at age zero. Its largest vertex is empty, and every other
vertex points to that largest vertex.

\begin{theorem}[Arbitrary population and cache capacity, $\ell=1$]
\label{thm:generalone}
Let $n\ge2$ and $2\le c\le n-1$ in
Section~\ref{sec:mandatorymodel}, with $\ell=1$.
Define, using this fixed capacity $c$,
\[
 B_1=0,\qquad B_m=2m(m+c-1)+2B_{m-1}\quad(2\le m\le n),
\]
\[
 T=n(n+c-1)+B_n,\qquad \gamma=(nc^2)^{-T}.
\]
From every weakly connected state with arbitrary finite integer ages,
\[
 P^T(x,\{H(V,c)\})\ge\gamma .
\]
The full cache-and-age chain therefore has a unique stationary
law $\pi$ and
\[
 \sup_x\|P^t(x,\cdot)-\pi\|_{\TV}
 \le(1-\gamma)^{\lfloor t/T\rfloor}.
\]
For every observer and every $1\le b\le c$, availability has a
positive stationary probability, and the local ordered $b$-sample
converges to uniform conditional on availability.
Together with Theorem~\ref{thm:mandatoryone}, this covers every
allowed capacity for shuffle length one.
\end{theorem}

\subsection{Clearing blocks and separating one source}

Write $\mathcal B(p)$ for the clearing block in
Theorem~\ref{thm:agereset}: activate only $p$; whenever its partner
already knows $p$, prescribe reply $p$, and otherwise choose any
reply. With $m$ active vertices it empties $p$ within $m+c-1$
events. No unselected resident of the initiator can be replaced
when $\ell=1$, since the sole usable reply fits the vacancy made
by partner removal. If all ages are zero at the block's start,
they are zero at its end: only $p$ increments ages, and its cache
is empty at the end; passive peers receive only fresh $p^0$.

\begin{lemma}[An unknown source with a connected remainder]
\label{lem:source}
From any zero-age weakly connected state on $V$, clear its
smallest vertex $a$, then clear each vertex of
$W=V\setminus\{a\}$, keeping $a$ passive. At the end all ages
are zero, no cache in $W$ contains $a$, $C_a$ is nonempty, and
the graph induced by $W$ is weakly connected.
\end{lemma}
\begin{proof}
Immediately after clearing $a$, every component of the induced
graph on $W$ has a descriptor for $a$: $C_a$ is empty while the
whole graph is connected.
During the remaining blocks the component partition on $W$
can only coarsen. A $W$-internal exchange preserves its internal
connectivity by edge reversal and transfer; if its reply is $a$,
it deletes no $W$ edge. An exchange with passive $a$ likewise
deletes no $W$ edge and may add one through the reply.

The first contact with $a$ inserts its first cached identity and
defines an anchor component in $W$. Thereafter every cached
identity at $a$ belongs to that anchor, which may grow.
Indeed a new initiator $p$ contacting $a$ receives an anchor
identity, joining the anchor, or receives itself and was already
in it. A nonempty cache at passive $a$ cannot become empty:
inserting a new initiator either fills a vacancy or replaces one
entry, while a duplicate changes nothing.

Every initial component joins the anchor. A descriptor for $a$
can disappear either when its holder contacts $a$, joining the
anchor directly, or when a full responder returns it and replaces
it by the active peer's fresh identity. In the latter case it
transfers to, or is already present at, that active peer. The
reversed edge links the two within $W$, and its clearing block
must eventually contact $a$. Thus eliminating a component's
last such descriptor forces that component into the anchor.

At the end no $W$ cache contains $a$: each vertex was emptied
and later passive requests contain only other initiators'
identities. Consequently all components have joined, and at
least one contact has made $C_a$ nonempty. Since $a$ is now
globally unknown, it cannot reappear while $a$ stays inactive.
All block boundaries have zero ages.
\end{proof}

\subsection{Extending the canonical remainder}

\begin{lemma}[Setting the cache of an unknown source]
\label{lem:extend}
Let $c\ge2$. Suppose all ages are zero, the smallest vertex $a$
is unknown in $W=V\setminus\{a\}$, the induced state on $W$
is $H(W,c)$, and $C_a$ is any nonempty cache in $W$.
There is a sequence of $|V|$ clearing blocks ending with all
ages zero, $a$ still unknown in $W$, $W$ weakly connected,
and $C_a$ equal to the $\min(c,|W|)$ greatest members of $W$.
\end{lemma}
\begin{proof}
Let $m=\max W$, $d=\min(c,|W|)$, let $J$ be the $d$
greatest members of $W$, and put $S=J\setminus\{m\}$.

First clear $a$, using reply $a$ from a responder that knows
$a$, and otherwise its smallest cached identity.
Every $W$ vertex other than $m$ retains its edge to $m$.
A full responder has at least two entries, so its smallest
entry is not $m$; an underfull responder evicts nothing.
The block must visit $m$. Otherwise let $q$ be the greatest
visited vertex. At its first visit its untouched canonical cache
returns a label greater than $q$. That label either is already
pending at $a$ or is inserted, and cannot disappear without being
visited, contradicting maximality of $q$.
Since $m$ initially is empty and remains passive, this phase
ends with $C_m=\{a^0\}$ and $C_a=\varnothing$.
The remainder stays connected through its edges to $m$.

Next, for each $j\in S$, start its clearing block with $j\to m$
and prescribe reply $a$. Preserve the descriptor for $m$ at every
unprocessed vertex: at a full responder not knowing $j$, choose
any reply other than $m$; at one knowing $j$, reply $j$.
Underfull responders need not evict anything.
Thus each unprocessed $j$ still knows $m$ when its block begins.
All ages are zero, so its first choice of $m$ is a legal tie.
Before that contact $C_m$ consists of $a$ and the previously
processed members of $S$, with size at most $d-1<c$.
Fresh $j$ therefore fills a vacancy and preserves $a$ there.
On later visits, $m$ already knows $j$ and replies $j$.
The active $j$ obtains or retains a descriptor for $a$, so must
contact $a$ before its block ends.
Passive $a$ has fewer than $c$ entries during this phase.
After these blocks,
\[
 C_a=S,\qquad C_m=\{a\}\cup S,
\]
all at age zero, and every processed $j$ is free of $a$.

Now clear all $h\in W\setminus J$.
Keep $a$ at passive $m$: if $m$ is full and does not know $h$,
choose a reply other than $a$, possible because $c\ge2$;
otherwise use the ordinary clearing policy.
At passive $a$, protect every member of $S$.
If this phase is nonempty, $|S|=c-1$, so there is at most one
unprotected slot. Fill a vacancy for a new $h$, or return and
replace the unique unprotected entry. For a duplicate $h$,
reply $h$. These choices respect the clearing-block rule.
At the end the only remaining descriptor for $a$ is at $m$:
every other vertex was cleared after its last possible
acquisition of $a$, and later fresh requests cannot reintroduce
that identity. Cache $C_a$ consists of $S$ and at most one
unprotected entry.

Finally start $m$'s clearing block with $m\to a$, a zero-age tie.
At $a$, reply with and replace its unprotected entry, if present;
otherwise fill its vacancy. Fresh $m$ is new there, since $m$
has never initiated in these phases. Hence $C_a$ becomes
exactly $J$. This removes the globally last descriptor for $a$.
Finish clearing $m$; it cannot contact $a$ again, so $C_a$ stays
fixed. All ages are zero at the end.

After the first phase $W$ was connected. Thereafter $a$ is
passive; exchanges with it only add $W$ edges, while internal
exchanges preserve $W$ connectivity. Thus $W$ remains connected.
The number of clearing blocks is
$1+|S|+|W\setminus J|+1=|V|$.
\end{proof}

\subsection{Induction and the convergence bound}

\begin{proof}[Proof of Theorem~\ref{thm:generalone}]
For the induction allow a fixed configured capacity $c\ge2$
even on a smaller active population. Excess cache slots cause
no difficulty; the cache is still a set of nonself identities
from that population. The one-vertex case is empty.

On a zero-age connected $m$-vertex state, apply
Lemma~\ref{lem:source}, then freeze its unknown smallest vertex
$a$. Use induction to make $W$ canonical. Nothing can contact
$a$, so its cache stays fixed. Apply Lemma~\ref{lem:extend} to
set $C_a$ to the greatest $\min(c,|W|)$ labels, leaving $a$
unknown, $W$ connected, and all ages zero.
Freeze $a$ again and make $W$ canonical a second time.
This reaches exactly $H(V,c)$.
The two lemmas use $m$ clearing blocks each, of at most
$m+c-1$ events, and the two recursive calls cost at most
$2B_{m-1}$. Thus the stated recurrence bounds the entire path.

From arbitrary ages first apply the clearing sweep of
Theorem~\ref{thm:agereset}, costing $n(n+c-1)$ events and
leaving a connected zero-age state. The total path length is
at most $T$. Pad shorter paths at the empty largest vertex of
$H(V,c)$. At each prescribed event the initiator has probability
$1/n$, the chosen oldest tie has probability at least $1/c$,
and the prescribed nonempty reply has probability at least $1/c$.
Empty replies and no-ops require fewer choices.
Therefore the path probability is at least $(nc^2)^{-T}$.
Lemma~\ref{lem:atom} proves the full-state conclusion.

For any observer, relabel the same construction so that it is
the smallest vertex of the canonical state and therefore has
$c$ entries. Stationarity and the same minorization imply that
its full-cache probability is at least $\gamma$.
The transition rule and connected state space are equivariant,
so Theorem~\ref{thm:stationary} proves conditional uniformity
for every $1\le b\le c$.
\end{proof}

\begin{proposition}[A robustness consequence]
\label{prop:domination}
Let $P_1$ be the singleton-shuffle kernel in
Theorem~\ref{thm:generalone}. Suppose another equivariant
kernel $Q$ preserves the same weakly connected aged state
space and satisfies $Q(x,A)\ge\varepsilon P_1(x,A)$ for every
state $x$ and event $A$, for a fixed $\varepsilon>0$.
Then $Q$ also converges to a unique stationary law, with
conditional uniform local samples, and has the bound
\[
 \sup_x\|Q^t(x,\cdot)-\pi_Q\|_{\TV}
 \le(1-\varepsilon^T\gamma)^{\lfloor t/T\rfloor}.
\]
\end{proposition}
\begin{proof}
Kernel domination gives $Q^T\ge\varepsilon^T P_1^T$.
Use the same common atom and then
Lemmas~\ref{lem:atom} and Theorem~\ref{thm:stationary}.
Relabelled canonical atoms again give positive availability.
\end{proof}

For example, this covers a protocol that independently chooses a
singleton shuffle with probability $\varepsilon$ and otherwise
takes any equivariant step preserving the connected state space.
It is a corollary for that explicitly randomized-length protocol;
it does not assert domination for fixed $\ell>1$ CYCLON.

\section{A complete multi-entry mandatory-removal theorem}
\label{sec:three}

\begin{theorem}[Three peers with two cache entries]
\label{thm:three}
In the mandatory-removal model of Section~\ref{sec:mandatorymodel},
let $n=3$ and $c=\ell=2$. Start from any weakly connected identity
overlay, with arbitrary finite nonnegative integer ages. The complete
cache-and-age chain has a unique stationary law $\pi$ on this state
space, and
\begin{equation}
 \sup_x\|P^t(x,\cdot)-\pi\|_{\TV}
 \le(1-6^{-20})^{\lfloor t/20\rfloor}.
 \label{eq:threebound}
\end{equation}
For every observer and $b=1,2$, the stationary probability of
availability is positive and the conditional local sample converges
to $\Unif((V\setminus\{i\})_b)$.
\end{theorem}

The bound is deliberately crude. This is a proof for actual
oldest-first, age-carrying exchanges with multi-entry caches, but it
does not establish a theorem for arbitrary $n,c,\ell$. At $n=3$ the
conditional $b=2$ output is automatically uniform when both possible
neighbors are present; convergence and positive availability still
require proof.

\subsection{Three-peer transition rules}

Write $q^a$ for a descriptor of identity $q$ and age $a$.
A full cache contains both possible other identities. The responder
never needs to delete an entry: it simply adds any new nonself
identity, since only two such identities exist. A full initiator
ends with exactly its other, uncontacted entry, whose age was
incremented. It cannot receive the contacted identity from that
identity's own cache. A singleton initiator ends with the responder's
third identity, if present, and is otherwise empty.

The weakly connected state space is closed. If a full cache is
present after an exchange, all three vertices are connected.
Otherwise an exchange with an empty responder reverses an edge,
and a reciprocal singleton exchange deletes only one of the two
reciprocal arcs. A path through distinct singleton initiator and
responder creates a full responder. These cases cover every exchange
and preserve weak connectivity.

\subsection{Resetting arbitrary ages on a chain}

For three distinct role names $F,S,T$, call the state a \emph{chain}
if
\begin{equation}
 C_F=\{S^a,T^b\},\qquad C_S=\{F^d\},\qquad C_T=\{S^e\}.
 \label{eq:chain}
\end{equation}
Define the zero-age state $Z(F,S,E)$ by
\[
 C_F=\{S^0,E^0\},\qquad C_S=\{E^0\},\qquad C_E=\varnothing .
\]

\begin{lemma}[Chain age reset]
\label{lem:chainreset}
Every chain reaches a state of form $Z$ in at most eight exchanges,
regardless of its finite initial ages.
\end{lemma}
\begin{proof}
Suppose first $b\ge a$. Prescribe the following five exchanges:
\[
 F\to T,\quad S\to F,\quad F\to S,\quad T\to S,\quad T\to F.
\]
The complete trace is
\[
\begin{array}{c|lll}
 \text{after}&C_F&C_S&C_T\\ \hline
 \text{initial}&\{S^a,T^b\}&\{F^d\}&\{S^e\}\\
 F\to T&\{S^{a+1}\}&\{F^d\}&\{S^e,F^0\}\\
 S\to F&\{S^{a+1}\}&\varnothing&\{S^e,F^0\}\\
 F\to S&\varnothing&\{F^0\}&\{S^e,F^0\}\\
 T\to S&\varnothing&\{F^0,T^0\}&\{F^1\}\\
 T\to F&\{T^0\}&\{F^0,T^0\}&\varnothing
\end{array}
\]
The fourth choice is legal because the two ages at $T$ after its
increment are $e+1$ and $1$; $S$ is oldest, possibly tied when $e=0$.
All other partner choices are forced by singleton caches, apart
from the first, which is allowed by $b\ge a$.
The final state is $Z(S,F,T)$.

If instead $a>b$, activate $F$ toward $S$. This produces a chain
with roles $(S,T,F)$ and full-cache ages $(b+1,d)$ for its near
and far vertices, respectively. If $d\ge b+1$, the preceding reset
applies. Otherwise integrality gives $d\le b$; rotate again by
activating the new full vertex toward its near vertex.
The roles are now $(T,F,S)$ and the full ages are $(d+1,e)$.
If $e\ge d+1$, reset. Otherwise $e\le d\le b$; a third rotation
gives roles $(F,S,T)$ with full ages $(e+1,b+1)$.
Now the far vertex is oldest, and the five-exchange reset applies.
At most three rotations and five further exchanges are needed.
\end{proof}

\subsection{Reaching the chain topology}

\begin{lemma}[Identity reduction compatible with oldest-first choices]
\label{lem:chainreach}
Every weakly connected three-peer state reaches a chain in at most
seven exchanges.
\end{lemma}
\begin{proof}
Classify states by the numbers of full, singleton, and empty caches.
All named activations below use an oldest entry; where it matters,
the age inequalities are given explicitly.

\emph{Two full caches and one singleton.}
Let $A,B$ be full and $C$ singleton. Activate $A$.
If it contacts $B$, there is now one full cache and two singletons.
Otherwise it contacts $C$, which becomes full, and $A$ becomes
singleton. If $C$ previously pointed to $B$, its full cache has
ages $B^u,A^0$ for some $u\ge0$. Activate $C$ toward the full
vertex $B$; this is oldest or tied and reduces to one full cache.
If $C$ previously pointed to $A$, its cache is
$\{A^u,B^{v+1}\}$, where $v$ was the age of $B$ at $A$.
If $B$ is oldest, again activate $C$ toward $B$.
Otherwise activate $C$ toward the singleton $A$. Then $A$ is full
with cache $\{B^{v+1},C^0\}$. It can contact $B$, its strictly
oldest entry. Thus at most three exchanges reach one full cache
and two singletons.

\emph{Other states with at least two full caches.}
Three full caches become two full and one singleton after one
activation. With two full and one empty, activate a full cache.
It either contacts the other full cache, yielding one full, one
singleton, and one empty, or contacts the empty vertex, yielding
two full caches and a singleton. Use the preceding case thereafter.

\emph{One full cache and two singletons.}
If the singleton arrows form a chain ending at the full vertex,
the required form already holds. If both point to the full vertex,
activate either singleton; it now points to the other singleton,
giving a chain. The remaining case has reciprocal singleton arrows.
Let $Q$ be an oldest neighbor of the full vertex $F$ and $R$ the
other vertex. Activate $Q\to R$, emptying $Q$, then $F\to Q$.
The resulting state is a chain with full vertex $Q$.
At most two exchanges are needed.

\emph{One full $F$, one singleton $S$, and one empty $E$.}
If $S\to F$, first activate it, giving $S\to E$.
If the oldest neighbor of $F$ is $S$, activate $S\to E$, emptying
$S$; this does not change the ages at $F$.
Now activate $F$ toward an empty oldest neighbor.
This gives one full cache and two singletons, handled above.
There are at most three preliminary exchanges.

\emph{One full and two empty.}
One full activation gives the preceding pattern.

\emph{No full caches.}
With three singletons, weak connectivity implies a directed path
$P\to Q\to R$ on distinct vertices. Activate $P\to Q$; $Q$ becomes
full and the other caches are singleton.
With two singletons and one empty, there is either a path toward
the empty root or two arrows into that root. In the path case,
activate the leaf toward its singleton parent, producing one full,
one singleton, and one empty. In the two-arrow case, reverse one
arrow by activating it toward the empty root, then use the path
case. This takes at most two preliminary exchanges.
There are no other weakly connected size patterns.

The longest route just described is the last pattern: at most two
exchanges to one full, one singleton, and one empty; at most three
to one full and two singletons; and at most two to a chain.
It costs seven. The two-full/empty route costs at most
$1+3+2=6$, the three-full route at most $1+3+2=6$, and the
other routes are no longer.
\end{proof}

\subsection{A common zero-age state}

\begin{lemma}[Relabelling the zero-age roles by legal exchanges]
\label{lem:roles}
Every $Z(F,S,E)$ reaches every prescribed labelled state of form
$Z$ in at most four exchanges.
\end{lemma}
\begin{proof}
The exchange $S\to E$ maps
$Z(F,S,E)$ to $Z(F,E,S)$ in one step.
The two exchanges $F\to S$ followed by $F\to E$ map
$Z(F,S,E)$ to $Z(S,E,F)$. The first choice is an age-zero tie;
before the second exchange the remaining entry of $F$ has age one.
Direct substitution shows that both final states have all ages zero.

If the desired full vertex is currently $S$, use the two-step
rotation. If it is currently $E$, first swap $S,E$ and then
rotate, costing three steps. If it is already $F$, do nothing.
One final swap, if needed, sets the two remaining roles.
\end{proof}

\begin{proof}[Proof of Theorem~\ref{thm:three}]
Lemmas~\ref{lem:chainreach} and~\ref{lem:chainreset} reach some
zero-age $Z$ in at most $7+8=15$ exchanges.
Lemma~\ref{lem:roles} reaches the fixed state $z=Z(1,2,3)$
within four more. Pad the resulting path to exactly twenty
activations by repeatedly activating the empty vertex at $z$.
Every prescribed initiation has probability $1/3$ and every
prescribed oldest choice has probability at least $1/2$.
The buffers are deterministic in this parameter setting, and the
responder never chooses victims. Therefore
\[
 P^{20}(x,\{z\})\ge6^{-20}\quad\text{for every weakly connected }x.
\]
Lemma~\ref{lem:atom} proves existence, uniqueness, and
\eqref{eq:threebound}, without truncating ages or assuming a finite
state space.

The same construction works for any relabelled $z$.
Integrating its minorization against $\pi$ shows
$\pi(\{z'\})\ge6^{-20}$ for each labelled $Z$ state $z'$.
For each observer choose one with that observer full. Thus
$\pi\{|C_i|\ge b\}\ge6^{-20}>0$ for $b=1,2$.
Theorem~\ref{thm:stationary} gives the conditional local conclusion.
\end{proof}

\section{Controlled progress for general shuffle lengths}
\label{sec:fresh}

The next lemma addresses one difficulty caused by larger buffers:
copying a newly created descriptor can evict another such descriptor.
It concerns a prescribed legal exchange, rather than the average
behavior of the random process.

Attach an auxiliary color, old or fresh, to each descriptor.
Copies inherit their source color, and every new age-zero self
descriptor is fresh. Discarding a duplicate does not change the
stored descriptor's color. Colors are proof bookkeeping and do
not change the protocol or its oldest-first selection.

\begin{lemma}[A controlled fresh-count inequality]
\label{lem:fresh}
Use Section~\ref{sec:mandatorymodel} with arbitrary $n,c,\ell$.
Suppose a nonempty initiator $p$ selects an old partner descriptor
$q$. For every prescribed extra initiator buffer $S$, there exist
a responder buffer and eligible victim choices for which the global
number $N_f$ of stored fresh descriptors does not decrease.
It increases by at least one if the responder accepts any fresh
incoming descriptor. All these choices have positive probability
under uniform buffer and eligible-victim selection.
\end{lemma}
\begin{proof}
Let $P=C_p$ and $Q=C_q$ denote the pre-exchange identity sets,
$a=|P|$, $b=|Q|$, $k=|S|=\min(\ell-1,a-1)$, and
$m=\min(\ell,b)$. Let $S_f$ be the fresh descriptors in $S$,
using their identities when taking intersections. The responder
must insert
\[
 J=(\{p\}\cup S)\setminus Q,\qquad j=|J|.
\]
Of these, $f_J=1_{\{p\notin Q\}}+|S_f\setminus Q|$ are fresh.
There are $v_Q=c-b$ vacancies, so
$r_Q=(j-v_Q)_+$ entries must be replaced.

Call a descriptor in $Q$ protected if it is fresh and its identity
lies in $P\cup\{p\}$. Put $d_0=|Q\cap(\{p\}\cup S)|$.
Since $q\notin Q$, at most $d_0+(a-1-k)$ descriptors are protected.
Thus the number of nonprotected entries is at least
\[
 b-d_0-a+1+k
 \ge k+1-d_0-(c-b)=j-v_Q .
\]
Choose $r_Q$ nonprotected victims $D$. Choose the responder buffer
$T$ to include these victims and the entries with identities
\[
 B=(S_f\cap Q)\cup(\{p\}\cap Q).
\]
This is possible: the union lies in $Q$, and if $r_Q>0$ its size
is at most $d_0+r_Q=k+1-v_Q\le\ell$.
If $r_Q=0$, its size is at most $k+1\le\ell$.
Complete $T$ arbitrarily to size $m$. Fill vacancies and then
replace exactly $D$. Every fresh victim at the responder has
identity absent from $P\cup\{p\}$; its fresh copy in $T$ is
therefore accepted by the initiator.

The initiator accepts $I=T\setminus(P\cup\{p\})$.
Write $i=|I|$, $f_I$ for its fresh count, $f_S=|S_f|$,
$d=|S_f\cap Q|$, and $h=1_{\{p\in Q\}}$.
Because $T$ includes $B$, $i\le m-d-h$.
Deleting the old partner leaves $v_P=c-a+1$ vacancies.
Choose old extra-buffer victims before fresh ones. The number of
fresh descriptors the initiator must remove is at most
\[
 \bigl(i-v_P-(k-f_S)\bigr)_+
 \le(f_S-d-h)_+=(f_J-1)_+.
\]
Here $m-v_P\le k$ follows from both
$m-v_P\le\ell-1$ and $m-v_P\le a-1$.
The responder loses at most $f_I$ fresh descriptors, all of which
have fresh copies accepted at the initiator. Combining the two
endpoints gives
\[
 \Delta N_f\ge f_J-(f_J-1)_+
 =1_{\{f_J>0\}} .
\]
This proves the controlled inequality and its strict version.
\end{proof}

This is neither a pathwise monotonicity theorem for every random
exchange nor an expected-drift inequality. It does not handle an
activation that removes a fresh partner. It also does not show that
a strict increase can always be arranged within a bounded number
of activations. That strict-progress problem remains a gap.

There is a limited bridge to numerical ages: over $H$ global events,
any descriptor lineage gains at most $H$ age units, and copying
preserves its age. Let $A$ be the maximum of zero and all initial
fresh-descriptor ages. If every initially old descriptor has age
greater than $A+H$, an initiator containing
an old descriptor selects an old partner throughout this horizon;
new fresh descriptors start at age zero. A bounded colored
clearing strategy compatible with actual oldest choices would
therefore help construct an age reset. Such a strategy has not
been proved here for arbitrary parameters.

\section{Numerical age control and initial-lineage erasure}
\label{sec:lineages}

\begin{lemma}[Eviction under consecutive activations]
\label{lem:forcing}
Use Section~\ref{sec:mandatorymodel} with arbitrary $n,c,\ell$.
Suppose a descriptor $r$ of age $a$ is present at $p$.
Let $M_s$ be the maximum initial age of identity $s$ anywhere in
the network, using zero if it is absent, and let $B$ be the
initial maximum age. If only $p$ is activated, the resident
descriptor $r$ is removed, either as partner or as a replacement
victim, within
\[
 U=1+\sum_{s\notin\{p,r\}}(M_s-a+1)_+
 \le1+(n-2)(B+1)
\]
activations. If the replacement choices retain $r$ until it is
selected, it is selected within this bound. In particular, an
initially globally maximum-age descriptor cannot remain resident
through $n-1$ consecutive activations of its holder.
\end{lemma}
\begin{proof}
Let $t$ count completed activations of $p$, and put
$K_s(t)=M_s(t)-t$, with $M_s(t)$ the current global maximum
age of identity $s$, again zero if absent. For $s\ne p$,
$K_s$ never increases: only $p$ increments ages, copies retain
their ages, and no new self descriptor of identity $s$ is born.
When $s$ is selected as partner, its sole descriptor at $p$ is
deleted and is absent from the outgoing buffer. Responder $s$
cannot return its own identity. Every other descriptor of identity
$s$ retains its age, so $K_s$ decreases by at least one.

While the original $r$ remains resident, its age after $t$
activations is $a+t$. A different partner $s$ selected next must
therefore satisfy $K_s(t)\ge a$. Each $s$ can be selected at most
$(M_s-a+1)_+$ times before $r$ is removed. Summing these counts
proves the first bound. If $a=B$, every summand is at most one,
giving $U\le n-1$.
\end{proof}

\begin{lemma}[Protecting one resident descriptor]
\label{lem:protect}
If $\ell<c$, suitable positive-probability buffer and victim
choices can retain any one designated nonpartner descriptor at
the initiator and any one designated descriptor at the responder.
If $c=n-1$, no nonpartner descriptor is ever replaced, for any
$\ell\le c$.
\end{lemma}
\begin{proof}
For an initiating cache of size $a>\ell$, choose the extra buffer
to exclude its protected entry. If $a\le\ell$, all $a-1$
nonpartner entries are sent, but mandatory removal creates
$c-a+1\ge\ell-a+2$ vacancies. At most $\ell$ usable replies
require at most $a-2$ replacements, leaving room to preserve
one of those $a-1$ entries.

For a responder of size $b>\ell$, exclude its protected entry
from its reply buffer. If $b\le\ell$, it sends all entries, but
$c-b\ge\ell+1-b$ vacancies permit at most $b-1$ replacements.
These buffer and victim choices can be made simultaneously at
the two endpoints.

For $c=n-1$, any cache union with usable nonself incoming
identities has size at most $c$. At the initiator the removed
partner cannot be returned by that partner. Thus all usable
incoming identities fit vacancies, without replacing a resident.
\end{proof}

\begin{lemma}[Following a target through forced replacements]
\label{lem:tracking}
For every $1\le\ell\le c$, a descriptor of identity $r$ and
age $a$ can be followed through a positive-probability sequence
until its current holder selects $r$ as partner. If the initial
global maximum age is $B$, this takes at most
\[
 1+\sum_{s\ne r}(M_s-a+1)_+\le n(B+1)
\]
events, where $M_s$ is the initial global maximum age of identity
$s$, with value zero if absent. The followed copy gains exactly
one age unit at every event preceding its deletion.
\end{lemma}
\begin{proof}
Always activate the current holder $p$.
If $\ell<c$, retain the tracked nonpartner by
Lemma~\ref{lem:protect}. Consider the remaining case $\ell=c$,
when the outgoing buffer contains every nonpartner resident.
Retain the target whenever possible. If this is impossible,
every one of the $a_p-1$ nonpartner residents must be replaced,
where $a_p$ is the current cache size. Since mandatory deletion
creates $c-a_p+1$ vacancies, the number $i$ of usable replies
must satisfy $i-(c-a_p+1)=a_p-1$, hence $i=c$.
The responder's full cache is consequently disjoint from the
initiator's pre-exchange cache and self identity.
In particular it does not know $r$. The outgoing tracked copy
of $r$ is accepted there, with its incremented age. Follow it.

For arbitrary changing holders, put $K_s(t)=M_s(t)-t$ for every
identity, taking its maximum age as zero if absent.
This quantity never increases: ages gain at most one per event,
copies preserve ages, and fresh self descriptors start at zero.
When $s$ is selected as partner, all other descriptors of identity
$s$ keep their ages, while the incremented one is removed and
cannot be copied or returned by responder $s$. Thus $K_s$
decreases by at least one.
Until the tracked $r$ is selected, its age is $a+t$.
Each different selected identity $s$ must have $K_s(t)\ge a$,
and can be selected at most $(M_s-a+1)_+$ times.
Sum over all $s\ne r$; unlike Lemma~\ref{lem:forcing}, the
holder changes, so no fixed holder identity can be excluded.
\end{proof}

\begin{theorem}[Accessibility of initial-lineage erasure, all parameters]
\label{thm:lineages}
For every $1\le\ell\le c\le n-1$, from every state of
Section~\ref{sec:mandatorymodel} with finite
maximum age $M_0$, there is a finite positive-probability execution
that erases every initially present descriptor and all descendants
obtained by copying it. With $D_0=2n(n-1)$, the execution can be
chosen to have length at most
\[
 T_0=((n+1)^{D_0}-1)(M_0+1).
\]
Every descriptor remaining at its end was generated during this
execution and has age at most its length.
Connectivity is not needed for this accessibility statement.
\end{theorem}
\begin{proof}
Mark all initial descriptors old, let their copies inherit the
mark, and mark every newly generated self descriptor fresh.
We eliminate the old descriptors one identity at a time.
Fix an identity $r$. Within its stage, keep a growing set of
immunized peers, each storing a fresh descriptor for $r$.
Keep immunized peers passive and preserve their fresh $r$.
For $\ell<c$ this is possible by Lemma~\ref{lem:protect}.
For $\ell=c$, whenever a received buffer contains an identity
already known at its responder, there are at most $c-1$ usable
incoming entries. A responder of size $b$ then replaces at most
$(c-1-(c-b))_+\le b-1$ residents, permitting protection of one.

Start with any old copy of $r$ and apply the moving-target
construction of Lemma~\ref{lem:tracking}.
No immunized peer is chosen as a holder. In the full-shuffle
case, every nonterminal outgoing packet contains the tracked
old $r$. If its responder is immunized, that entry is a
duplicate, allowing preservation of the fresh resident $r$.
Moreover such a responder replies with its own fresh $r$,
which is a duplicate at the old holder. It gives at most $c-1$
usable replies, so cannot force migration of the tracked copy
into that immunized responder.

When $r$ is finally selected by holder $q$, that holder loses
its descriptor for $r$, and responder $r$ cannot return it.
The outgoing fresh self descriptor ensures that $r$ stores $q$
afterward. If $q$ was already known, its duplicate leaves at
most $c-1$ other usable incoming entries in the full-shuffle
case, so its existing descriptor can be protected. If new,
the inserted descriptor is not an eligible victim in that exchange.

Now activate only the source peer $r$, keeping its cached $q$
until it is selected or until retention is impossible.
The eviction bound of Lemma~\ref{lem:forcing} makes this block
finite. If $q$ is selected, it has remained passive and free
of $r$, so accepts fresh $r$ and becomes immunized.
If $q$ is instead unavoidably evicted, the proof of
Lemma~\ref{lem:tracking} shows that the responder has $c$
usable replies and is missing identity $r$. It accepts the
outgoing fresh self descriptor of $r$ and becomes immunized.
Stop the block at this event. Either outcome adds a new
immunized peer. Every previously immunized responder sees the
duplicate fresh self descriptor $r$, so can preserve its resident
fresh $r$ even under full shuffling.

Each iteration adds a distinct immunized peer, and none
reacquires an old $r$. After at most $n-1$ iterations all old
descriptors for $r$ are gone. Extinction is irreversible:
copying preserves identity and new self descriptors are fresh.
Repeat for all identities, using at most $D_0$ routing or
single-source forcing blocks.

For the length bound let $S$ be elapsed execution length and
put $R=M_0+S+1$. Current ages are at most $M_0+S$.
Each block has length at most $nR$ by the two numerical-age
lemmas, so it multiplies $R$ by at most $n+1$.
After at most $D_0$ blocks,
$M_0+S+1\le(n+1)^{D_0}(M_0+1)$, proving the bound.
Every remaining descriptor starts at age zero during this
execution and gains at most one per event. All prescribed choices
have positive probability, proving accessibility.
\end{proof}

The bound depends on $M_0$; this is not a uniform finite-set
regeneration or a recurrence theorem. Erasing initial descriptor
lineages alone does not erase the influence of the initial topology
on the locations of the fresh descriptors.

\subsection{Maximal capacity and connectivity}

\begin{proposition}[Monotone connectivity at maximal capacity]
\label{prop:maxcapacity}
For $c=n-1$ and every $1\le\ell\le c$, the underlying undirected
graph never loses an edge. From every weakly connected state
with finite ages, a complete underlying graph is accessible by
a finite positive-probability path. Any stationary probability
law on this state space, if one exists, is supported on complete
underlying graphs.
\end{proposition}
\begin{proof}
By Lemma~\ref{lem:protect}, the only removed resident descriptor
is the mandatory partner at the initiator. Its undirected edge
is preserved by the reverse descriptor at the responder.
All other existing edges persist.

If the connected underlying graph is not complete, it has an
induced two-edge path $u-v-w$ with $u,w$ nonadjacent.
Lemma~\ref{lem:forcing} permits forcing contact with any
currently cached target by activating only its holder; no
nonpartner can be replaced in this regime.
If the orientation contains a directed path $u\to v\to w$,
force $u$ to contact $v$ and include $w$ in the reply. This
adds the missing edge, unless it was already added along the way.
If $u\to v\leftarrow w$, first force $u\to v$, creating
$v\to u$, then force $w\to v$ and reply with $u$.
If $u\leftarrow v\to w$, activate only $v$ until its first
contact with either endpoint. The other descriptor remains
resident, and the contacted endpoint learns $v$. Freeze $v$
and use the resulting directed path as in the first case.
Selecting the desired reply entry has positive probability for
every $\ell\ge1$. These cases cover the possible orientations,
including bidirectional edges. At most two forcing blocks add
an edge. Monotonicity allows repetition until the graph is complete.

Finally, edge count is bounded and nondecreasing. Under a
stationary law it cannot increase with positive probability
over any finite horizon. Every noncomplete state has a finite
positive-probability path increasing it, so cannot have stationary
mass. The state space is countable, which proves the support claim.
\end{proof}

By itself this proposition proves accessibility, rather than
almost-sure completion or existence of a stationary law.
The forcing path lengths may depend on ages. For $\ell=2$,
Corollary~\ref{cor:twocomplete} below upgrades this to uniform
almost-sure completion. For full shuffling
$\ell=c=n-1$, there is also the exact knowledge-set relation
\[
 K_q'=K_p\cup K_q,\qquad
 K_p'=(K_p\cup K_q)\setminus\{q\},
 \qquad K_i=C_i\cup\{i\}.
\]
A bookkeeping token moved to each responder accumulates knowledge
monotonically. Turning this into a uniform age-reset construction
remains an open step here.

\section{Uniform age control for shuffle length two}
\label{sec:two}

\begin{lemma}[Two-entry exchanges preserve weak connectivity]
\label{lem:twoconnect}
In Section~\ref{sec:mandatorymodel}, $\ell=2$ preserves every
weak component for every $c\ge2$, including underfull caches.
\end{lemma}
\begin{proof}
Let $p$ contact $q$, and let $r$ be its one extra sent identity,
if present. First suppose the responder ends with $p$ in its
cache. The removed partner edge is then preserved in reverse.
A responder victim $v$ was in its reply. If $v\ne p$ is new
at $p$, it is inserted there. If it was already present,
its duplicate consumes one of at most two reply entries, so
at most one usable reply remains. This fits the mandatory
vacancy, preventing any initiator-extra replacement.
Thus every lost $q$--$v$ edge has the path $q$--$p$--$v$.

If the initiator replaces its extra $r$, it must accept two
new reply entries; therefore $r$ is absent from the reply.
If $q$ already knew $r$, that descriptor was not sent and
cannot be a victim. Otherwise the outgoing $r$ is inserted
there. The lost $p$--$r$ edge has the path $p$--$q$--$r$.

If instead $q$ ends without $p$, it initially knew $p$ and
evicted that descriptor after discarding incoming $p$ as a
duplicate. Its only possible usable incoming entry is $r$.
Thus it inserts $r$, and $p$ belongs to its reply. That self
entry leaves at most one usable reply at $p$, so $p$ retains
$r$. The endpoints remain connected through $p$--$r$--$q$.
There is at most one responder victim, namely $p$, and no
initiator-extra victim in this case. Singleton and empty
cases obey the same reasoning with missing entries omitted.
No new identity can cross an initially separate component.
\end{proof}

\begin{lemma}[Uniform routing at shuffle length two]
\label{lem:tworouting}
For $\ell=2$ and every $c\ge2$, a chosen descriptor of identity
$r$ can be followed until its holder selects $r$ as partner
within $(nc+1)(2n+c)$ controlled exchanges, independently of
all initial ages. During the construction, any passive peer's
designated resident descriptor for $r$ can be preserved.
\end{lemma}
\begin{proof}
Use a separate auxiliary color for this routing block: mark
every currently stored descriptor old, and only newly generated
self descriptors fresh. These colors differ from any older
lineage markings used elsewhere. Keep following the chosen old
$r$, always activating its current holder $p$ and using $r$
as the extra entry when a different partner is selected.
At time $t$ its age is $a+t$, whereas any auxiliary fresh
descriptor has age at most $t-1$. Consequently every
nonterminal selected partner is auxiliary old.

We choose replies and victims so that the responder ends
knowing both $p,r$, the global auxiliary fresh count never
decreases, and it increases by at least one whenever that
responder did not know $p$.
Write $P=C_p$, $Q=C_q$, $a_p=|P|$ and $b=|Q|$.
Call a responder descriptor nonprotected for this paragraph
if it is auxiliary old or is auxiliary fresh with identity
outside $P\cup\{p\}$.
If $Q$ already knows both $p,r$, reply with that pair:
both are filtered at $p$, and there is no responder update.
If $Q$ knows exactly one of them, reply with it and, when
needed, a nonprotected victim different from it. Such a victim
exists when $Q$ is full: among its $c-1$ other descriptors,
at most $a_p-2\le c-2$ can be protected fresh identities.
An underfull responder needs no victim in this one-new-entry
case. Retain the known member and insert the missing one.

If $Q$ knows neither, it must insert $p,r$.
The number of required victims is $(2-(c-b))_+$.
There are enough nonprotected descriptors: at most $a_p-2$
of $Q$'s identities can be protected fresh, so their complement
has size at least $b-c+2$ when victims are needed.
Include these chosen victims in its two-entry reply.
Each fresh responder victim is thus copied to a previously
unknown location at $p$. The initiator's only eligible extra
victim is its old tracked $r$, and its mandatory partner is
also old. Therefore fresh losses at the responder are
compensated at $p$; accepting new fresh $p$ gives a strict
increase. The responder retains or inserts both members
$p,r$ in every case.

If $p$ cannot retain $r$, its cache is full and both reply
entries are new, so the responder knew neither $p$ nor $r$.
Follow the outgoing copy there. This migration strictly
increases the global auxiliary fresh count.
There can be at most $nc$ such migrations.

Between migrations the holder is fixed. Every partner visited
in this segment learns and retains the pair $p,r$.
On any repeat visit prescribe reply $\{p,r\}$; it is entirely
filtered and decreases the holder's cache size by one.
There are at most $n-2$ distinct nonterminal partners.
A first visit can increase its cache size by at most one;
therefore at most $c-1+(n-2)$ repeat-visit shrinkings occur
before only $r$ remains. Including the terminal selection
or migration, a segment has at most $2n+c$ events.
At most $nc+1$ segments prove the bound.
Every responder choice retains its existing $r$, so it also
preserves designated passive copies, regardless of their
other lineage markings.
\end{proof}

\begin{theorem}[A uniform finite-age reset for every capacity, $\ell=2$]
\label{thm:tworeset}
Choose eligible victim sets uniformly in
Section~\ref{sec:mandatorymodel}. For arbitrary $n,c$ with
$2\le c\le n-1$ and $\ell=2$, put
\[
 L=n(n-1)(nc+2)(2n+c),\qquad
 \delta=(nc^5)^{-L}.
\]
Let $F_L$ be the finite family of weakly connected states
whose stored ages are at most $L$. From every weakly
connected state with arbitrary finite integer ages,
\[
 P^L(x,F_L)\ge\delta,\qquad
 \mathbb E_x\tau_{F_L}\le L/\delta .
\]
In particular, no bound on initial ages is assumed.
\end{theorem}
\begin{proof}
Use the identity-by-identity immunization proof of
Theorem~\ref{thm:lineages}, replacing each old-copy routing
block by Lemma~\ref{lem:tworouting}. Its pair-retention choices
preserve all previously immunized passive peers.
At the terminal routing exchange, retain or insert the
holder's identity at its source, as in that immunization proof;
if it was already known, there is at most one other usable
incoming entry and a victim different from it can be chosen.

The subsequent source-only callback also has a uniform bound.
Let $s$ be the source identity being immunized and $q$ its
cached target. Whenever the partner is not $q$, send
$\{s^0,q\}$. Stop if the responder does not know $s$:
it accepts fresh $s$ and becomes immunized.
Otherwise include $s$ in its reply, and include $q$ too if
known. The self reply leaves at most one usable entry,
which fits the mandatory vacancy; hence $q$ is retained
and the active cache never grows. The responder learns or
retains both $s,q$. A repeat visit replies with that pair,
shrinking the active cache. There are at most $n-2$ first
partners and $c-1$ shrinking repeats before $q$ is selected.
The block thus costs at most $n+c-1$ events. Selection of
$q$ also immunizes it, since it remained passive and free of $s$.

At most $n(n-1)$ immunization iterations erase every initial
lineage. The resulting total length is at most
\[
 n(n-1)\bigl[(nc+1)(2n+c)+(n+c-1)\bigr]\le L.
\]
All remaining descriptors were born during this execution.
Pad to exactly $L$ by arbitrary valid events: old lineages
cannot reappear, and every remaining age is still at most $L$.
Connectivity holds by Lemma~\ref{lem:twoconnect}.

Every prescribed event has probability at least
\[
 \frac{1}{2nc(c-1)\binom c2}\ge\frac{1}{nc^5}:
\]
this accounts for initiator, oldest tie, chosen extra,
responder subset, and a victim subset among at most two
sent entries. Thus the controlled path has probability
at least $\delta$. Repeating $L$-event blocks gives the
geometric hitting tail and the expectation bound.
\end{proof}

For $c=\ell=2$ the proof sharpens to $L=2n^2(n-1)$:
a forced migration lands in the two-entry cache
$\{p^0,r^{a+t+1}\}$, so $r$ is selected next.
In a nonmigrating segment the full two-entry holder cannot
grow, giving at most $n$ events per old-copy block.
The callback block likewise takes at most $n$.

\begin{corollary}[Uniform age tails and stationary moments]
\label{cor:twotails}
Let $M_t$ be the largest stored age after $t$ activation
events, with value zero when there are no descriptors.
Under Theorem~\ref{thm:tworeset}, for every integer $k\ge1$
and $t\ge kL$,
\[
 \sup_x\Pr_x\{M_t>kL\}\le(1-\delta)^k,
\]
where the supremum is over weakly connected states with
arbitrary finite initial ages. Every stationary probability
law $\pi$ therefore satisfies
\[
 \pi\{M>kL\}\le(1-\delta)^k,\qquad
 \mathbb E_\pi M\le\frac L\delta .
\]
Moreover, for $0\le\theta<-L^{-1}\log(1-\delta)$,
\[
 \mathbb E_\pi e^{\theta M}
 \le\frac{\delta e^{\theta L}}
 {1-(1-\delta)e^{\theta L}}.
\]
\end{corollary}
\begin{proof}
The largest age increases by at most one at each event.
Partition the last $kL$ events before $t$ into $k$ blocks
of length $L$. Conditionally on the past, each block ends
with maximum age at most $L$ with probability at least
$\delta$. If even the earliest block succeeds, subsequent
age growth leaves $M_t\le L+(k-1)L=kL$.
Thus the event $M_t>kL$ requires every block to fail,
which has probability at most $(1-\delta)^k$.
Apply this bound to a stationary initial law. The integer
variable $\lceil M/L\rceil$ is stochastically dominated by
a geometric variable on $\{1,2,\ldots\}$ with success
probability $\delta$. Its mean and exponential moment give
the stated bounds.
\end{proof}

\begin{proposition}[Stationary regimes and continuous-time convergence]
\label{prop:tworegimes}
Under Theorem~\ref{thm:tworeset}, the embedded chain has
finitely many closed communicating classes, each with a
stationary probability law. With independent equal-rate
Poisson activation clocks, the full state converges in
total variation to an initial-state-dependent mixture of
these laws, uniformly at some geometric rate in elapsed time.
This does not yet assert a unique regime or uniform local samples.
\end{proposition}
\begin{proof}
Rescale time so the total activation rate is one.
The one-time-unit transition kernel is
\[
 R=e^{-1}\sum_{j=0}^\infty\frac{P^j}{j!}.
\]
It has a self-loop of probability at least $e^{-1}$ at every
state and dominates $e^{-1}P$. Hence
$R^L(x,F_L)\ge e^{-L}\delta$.
Lemma~\ref{lem:finiteclasses} applies to $R$.
Its communicating and closed classes are exactly those of
$P$, because every finite number of embedded transitions
has positive Poisson probability.
If $\pi_k$ is the unique $R$-stationary law of one class,
then $\pi_k P$ is another such law: $P$ and $R$ commute.
Uniqueness gives $\pi_k P=\pi_k$.
The eventual class probabilities are unchanged by Poisson
time change. Extend the integer-time convergence bound to
arbitrary elapsed times by Markov contraction for the remaining
fractional interval. Rescaling back to any common clock rate
only changes the constants.
\end{proof}

\begin{corollary}[Almost-sure topology completion at maximal capacity]
\label{cor:twocomplete}
Let $n\ge3$, $c=n-1$, and $\ell=2$, with uniform choices as
in Theorem~\ref{thm:tworeset}. Let $G$ be the set of aged
states with a complete underlying undirected graph, and put
\[
 h=(n-1)(n-2),\qquad
 T=(n+1)^h(L+1)-1,\qquad
 \alpha=(nc^5)^{-T}.
\]
Then, uniformly over weakly connected initial states with
arbitrary finite ages,
\[
 P^T(x,G)\ge\alpha,\qquad
 \Pr_x\{\tau_G>kT\}\le(1-\alpha)^k,\qquad
 \mathbb E_x\tau_G\le T/\alpha.
\]
After this almost-sure finite hitting time the underlying
graph stays complete. This does not assert a complete
directed graph, full caches, or uniqueness of the aged regime.
\end{corollary}
\begin{proof}
First perform the $L$-event reset of
Theorem~\ref{thm:tworeset}. With probability at least
$(nc^5)^{-L}$, all ages are now at most $L$.
Proposition~\ref{prop:maxcapacity} then adds each missing
underlying edge with at most two forcing blocks.
A connected graph initially lacks at most
$\binom n2-(n-1)$ edges, so at most $h$ blocks suffice.

If the current maximum age is at most $B$, a forcing block
has at most $n(B+1)$ events by Lemma~\ref{lem:forcing};
maximal capacity prevents nonpartner eviction. If $s$ events
have elapsed since the reset, use $B\le L+s$.
Across one such block the quantity $L+s+1$ grows by at most
a factor $n+1$. Thus the additional number of events is at most
\[
 [(n+1)^h-1](L+1)=T-L.
\]
Every prescribed event has probability at least
$(nc^5)^{-1}$, as in the reset theorem. Completion has
conditional probability at least $(nc^5)^{-(T-L)}$ within
this horizon. Since $G$ is absorbing, any shorter path can
be padded freely. Combining with the reset gives
$P^T(x,G)\ge\alpha$. Apply the Markov property to successive
$T$-event blocks to obtain the tail and expectation bounds.
\end{proof}

For the embedded chain, an aperiodicity argument is still needed
to replace continuous-time convergence by convergence at every
activation count. For conditional uniform local sampling, the
substantive remaining issue is a common class, or sufficient
relabelling invariance of every possible limiting regime.

\section{A replacement-rule obstruction for every larger shuffle}
\label{sec:largerobstruction}

The choice of eligible victims in Section~\ref{sec:mandatorymodel}
is substantive. In this section choose victim sets uniformly and
independently at the endpoints. A resident whose identity was
received and filtered as a duplicate is still an eligible victim
if it belongs to the sent buffer. This is the precise completion
addressed by the next theorem.

\begin{theorem}[Failure of conditional uniformity for $\ell\ge3$]
\label{thm:largerobstruction}
For every fixed $3\le\ell\le c$, set
\[
 D_{c,\ell}
 =c\binom{c-1}{\ell-1}\binom c\ell(\ell-1)\binom\ell2,
 \quad M=2(c+2),\quad p_*=\frac1{M^2D_{c,\ell}^2}.
\]
For every
\[
 n\ge
 \max\{3c+5,\ (c+2)(c+1)M^2D_{c,\ell}^2+2\},
\]
there is a deterministic weakly connected initial state with
all caches full and all ages zero from which some observer's
singleton output, conditional on availability, fails to converge
to uniform on the other $n-1$ peers.
This theorem concerns the stated uniform-victim completion;
it does not assert failure for every CYCLON implementation.
\end{theorem}
\begin{proof}
Take two disjoint core sets
\[
 A=\{p,q,a_1,\ldots,a_c\},\qquad
 B=\{v,w,b_1,\ldots,b_c\}.
\]
Let the remaining peers be $O$, of size at least $c+1$, with
one designated member $o$. Set
\[
\begin{array}{ll}
 C_p=\{q,v,a_1,\ldots,a_{c-2}\},&
 C_q=\{v,a_2,\ldots,a_c\},\\
 C_v=\{w,a_1,b_1,\ldots,b_{c-2}\},&
 C_w=\{a_1,b_2,\ldots,b_c\}.
\end{array}
\]
For each $a_i$ use $C_{a_i}=A\setminus\{a_i,q\}$.
For each $b_i$ use $C_{b_i}=B\setminus\{b_i,w\}$, except
that $C_{b_c}$ replaces $b_1$ by $o$.
Give every outside peer its next $c$ cyclic successors in $O$.
All caches have $c$ distinct nonself identities, at age zero.
The core is strongly connected: each group internally reaches
its source, the sources reach their group members, and the
displayed cross references join the groups in both directions.
The edge $b_c\to o$ joins the outside component.
The whole initial graph is weakly connected; strong connectivity
of the whole graph is not claimed.

Let $\tau$ be the second initiation by a core peer.
The first two core initiators are $v$ then $p$ with probability
$1/M^2$. All intervening initiations are outside and cannot
alter a core cache: outside caches contain only outside
identities, and neither designated core exchange involves
the boundary peer $b_c$ as an endpoint.

Fix $R_A\subseteq\{a_2,\ldots,a_{c-2}\}$ and
$R_B\subseteq\{b_2,\ldots,b_{c-2}\}$ of size $\ell-3$.
Prescribe the following first exchange:
$v$ chooses $w$, sends
$\{v,a_1,b_1\}\cup R_B$, and receives
$\{a_1,b_{c-1},b_c\}\cup R_B$.
Its two usable incoming identities fill the mandatory vacancy
and replace $a_1$. At $w$, the new identities $v,b_1$ replace
$a_1,b_{c-1}$ among its sent entries.
Thus both cross references from $B$ to $A$ disappear.

For the second exchange prescribe $p\to q$, outgoing buffer
$\{p,v,a_1\}\cup R_A$, and reply
$\{v,a_{c-1},a_c\}\cup R_A$.
The initiator receives two new identities and chooses $v$
as its additional victim. The responder inserts $p,a_1$
and replaces $v,a_{c-1}$.
Now every cache in $A$ contains only identities in $A$, and
every cache outside $A$ contains only identities outside $A$.
This partition persists forever.

Given its designated initiator, each prescribed exchange has
probability $D_{c,\ell}^{-1}$: oldest selection costs $1/c$,
the extra subset costs $\binom{c-1}{\ell-1}^{-1}$,
the reply costs $\binom c\ell^{-1}$, and the victim choices
cost $(\ell-1)^{-1}\binom\ell2^{-1}$.
The exact isolation event $E$ therefore has probability $p_*$,
independent of $n$. Also $\tau<\infty$ almost surely.

On $E$, at least one cache in $A$ remains nonempty forever
after $\tau$. Passive updates never decrease cache cardinality;
an empty responder accepts its initiator's fresh self entry.
Consequently every nonempty-initiator exchange leaves a
nonempty responder, and empty initiations change nothing.

For $i\in A$, let $a_i(t)$ be its availability probability,
$\nu_{i,t}$ its conditional singleton law, and
$d_i(t)=\|\nu_{i,t}-\Unif(V\setminus\{i\})\|_{\TV}$.
These conditional laws are defined at every finite $t$:
with positive probability $i$ has not initiated, and its
initially full cache cannot shrink through passive updates.
The isolation event gives
\[
 \sum_{i\in A}a_i(t)\nu_{i,t}(A\setminus\{i\})
 \ge\Pr(E,\tau\le t).
\]
Each corresponding uniform mass is $(c+1)/(n-1)$ and
$a_i(t)\le1$. Therefore
\[
 \liminf_{t\to\infty}\sum_{i\in A}d_i(t)
 \ge p_*-\frac{(c+2)(c+1)}{n-1}>0.
\]
For at least one fixed observer,
\[
 \limsup_{t\to\infty}d_i(t)
 \ge\frac{p_*-(c+2)(c+1)/(n-1)}{c+2}>0.
\]
This proves failure of conditional marginal uniformity,
rather than only a pathwise isolation event.
\end{proof}

For $c=\ell=3$, this construction gives $p_*=1/32400$ and
the explicit threshold $n\ge648002$.
The large threshold is sufficient for the proof, not a
prediction of the smallest failing network or a practical
failure probability. The theorem includes proper shuffles
$3\le\ell<c$, not only full-cache exchanges.
Its population bound is outside the maximal-capacity regime
$c=n-1$ of Proposition~\ref{prop:maxcapacity}.

\begin{corollary}[The obstruction also starts from strong connectivity]
\label{cor:strongobstruction}
Under the same uniform-victim rule, fix $3\le\ell\le c$ and
the quantities $M,D_{c,\ell},p_*$ of
Theorem~\ref{thm:largerobstruction}. Put
\[
 p_{\mathrm s}=\frac{p_*}{c+2}
 =\frac{1}{(c+2)M^2D_{c,\ell}^2}.
\]
For every
\[
 n\ge\max\{3c+5,\ (c+2)^2(c+1)M^2D_{c,\ell}^2+2\},
\]
there is a deterministic \emph{strongly connected} initial
state, with all caches full and all ages zero, from which
conditional singleton uniformity fails. More precisely, some
fixed observer $i$ satisfies
\[
 \limsup_{t\to\infty}d_i(t)
 \ge\frac{p_{\mathrm s}-(c+2)(c+1)/(n-1)}{c+2}>0.
\]
\end{corollary}
\begin{proof}
Keep the core and its edge $b_c\to o$ from the preceding
construction. Choose an outside gateway $z$. In its cache,
replace its second cyclic successor by $b_c$, retaining its
first cyclic successor. The outside directed successor cycle
survives. Thus the outside is internally strongly connected,
and the two edges $b_c\to o$ and $z\to b_c$ join it to the
strongly connected core in both directions. No outside cache
other than $C_z$ initially contains a core identity.

Use independent rate-one Poisson clocks to realize the same
iid-initiator embedded chain. Compare the outside evolution
with a standalone process on $O$ whose initial caches are the
unmodified cyclic successor caches, with the same parameters
$c,\ell$ and the same outside clocks and random choices.
Until an outside exchange first has $z$ as initiator or
responder, $C_z$ has not been read. Therefore every other
outside cache agrees pathwise with this standalone process,
provided no core exchange has yet reached outside.
Merely copying the identity $z$ does not read $C_z$.

Let $N_z(u)$ count outside initiations by $z$ and outside
exchanges whose responder is $z$ in the standalone process
up to elapsed time $u$. Rotating all outside labels leaves
its initial state and kernel invariant. Consequently, writing
$d_z^-(u)$ for its current indegree,
\[
 \mathbb E d_z^-(u)
 =\frac1{|O|}\mathbb E\sum_{j\in O}|C_j(u)|\le c.
\]
The instantaneous initiation rate at $z$ is one, and the
incoming contact rate is at most $d_z^-(u)$: each peer
containing $z$ initiates at rate one and selects it with
probability at most one. Hence
\[
 \mathbb E N_z(u)\le(c+1)u.
\]

Let $\tau$ be the time of the second core initiation.
It is independent of the standalone outside process, has
mean $2/M$, and is independent of the core initiator labels.
Thus Markov's inequality gives
\[
 \Pr\{N_z(\tau)=0\}
 \ge1-\mathbb E N_z(\tau)
 \ge1-\frac{2(c+1)}M=\frac1{c+2}.
\]
Require also that the first two core initiators are $v,p$
and that their exchanges follow the preceding certificate.
They involve only endpoints $v,w,p,q$, so they read and
modify no outside cache. On the no-touch event the gateway
never exports its core identity, and no other outside peer
can contact the core. The coupling is therefore valid
through $\tau$. Core labels have probability $M^{-2}$,
independently of $\tau$ and the outside process; the two
prescribed exchanges each have probability $D_{c,\ell}^{-1}$.
Their intersection isolates $A$ permanently with probability
at least $p_{\mathrm s}$.

Apply the final conditional-total-variation argument of
Theorem~\ref{thm:largerobstruction} with $p_{\mathrm s}$ in
place of $p_*$. The Poisson construction only bounds the
probability of a finite embedded-chain prefix; the resulting
nonconvergence statement holds for the iid activation-count
chain as well as its Poisson time change.
\end{proof}

For $c=\ell=3$ the stronger initialization gives
$p_{\mathrm s}\ge1/162000$ and the sufficient population
threshold $n\ge3240002$. These constants again serve to prove
existence, rather than to locate the smallest failing network.
Strong connectivity at initialization therefore does not
repair this completion's general convergence claim.

\section{An explicit protection rule and a proved extension}
\label{sec:protected}

Consider the following alternative completion of cache updates:
record all received identities before filtering duplicates, and
forbid eviction of any resident whose identity appears in that
received buffer. Insert usable new identities as before.
In this section duplicate ages still remain unchanged, as in
Section~\ref{sec:mandatorymodel}; minimum-age merging is introduced
separately in Section~\ref{sec:thesis}.
Call this the received-identity protection rule.
It makes explicit the descriptor-retention condition needed by
the paper's edge-transfer connectivity argument
\cite[Section 3.1]{cyclon}. We state it as a specified rule;
the paper's abbreviated replacement description does not uniquely
choose between all completions studied here.

\begin{lemma}[Protection is feasible and preserves connectivity]
\label{lem:receiveprotect}
Received-identity protection is feasible for every $n,c,\ell$
and every underfull exchange. Under this rule every weak
component is preserved.
\end{lemma}
\begin{proof}
At the initiator let $v$ be its vacancy count after mandatory
partner removal, $k$ the number of extra sent residents, and
$e$ the number of those residents protected by received
identities. If $i$ incoming entries are usable, then
$i+e\le\ell$ and $v+k\ge\ell$.
Thus $(i-v)_+\le k-e$, so enough unprotected eligible
victims remain. At the responder use its vacancy count $v$,
reply size $m$, protected eligible count $e$, and usable
incoming count $j$. Again $j+e\le\ell$ and $v+m\ge\ell$,
so $(j-v)_+\le m-e$. Newly inserted entries are not victims.

The responder retains or inserts the initiator's identity,
which is always in the outgoing packet; hence the removed
partner edge is preserved in reverse.
Every removed initiator extra was sent and therefore is
retained or inserted at the responder.
Every removed responder entry was returned and is retained
or inserted at the initiator. The latter cannot be the
initiator's self identity, because it is protected at the
responder. Each lost edge has a path through the two
exchanging peers. New references stay within their previous
weak component.
\end{proof}

The designated cross-reference victims in
Theorem~\ref{thm:largerobstruction} were received duplicates.
This rule forbids those choices, so that theorem does not
apply to this completion.
For $\ell=1$ the protected kernel equals the earlier kernel:
filtering a duplicate leaves no usable entry that could evict it.
For $\ell=2$ the pair-preserving construction in
Theorem~\ref{thm:tworeset} uses only protected-legal moves:
both received identities are retained at responders, and an
extra received as a duplicate at the initiator leaves at most
one usable reply, fitting its mandatory vacancy.
Thus both the singleton convergence theorem and the two-entry
finite-age reset transfer to this rule.

\begin{theorem}[A protected sampler with occasional singleton shuffles]
\label{thm:protectedmixture}
Fix any $n,c,\ell$ with $2\le\ell\le c\le n-1$ and
$0<\varepsilon\le1$.
At each activation independently use shuffle length one with
probability $\varepsilon$, and shuffle length $\ell$ otherwise;
use received-identity protection and uniform eligible choices.
From every weakly connected state with arbitrary finite ages,
the full process converges uniformly to a unique stationary
law and its local samples converge to uniform conditional on
availability. With $T,\gamma$ from
Theorem~\ref{thm:generalone}, its bound is
\[
 \sup_x\|Q^t(x,\cdot)-\pi_Q\|_{\TV}
 \le(1-\varepsilon^T\gamma)^{\lfloor t/T\rfloor}.
\]
\end{theorem}
\begin{proof}
Lemma~\ref{lem:receiveprotect} supplies a closed connected
state space. The singleton kernel is unchanged, so
$Q\ge\varepsilon P_1$. Apply
Proposition~\ref{prop:domination}.
\end{proof}

This is a complete theorem for an explicitly modified sampler.
It is not a proof for fixed-length $\ell>1$ CYCLON.
For the protected fixed-length sampler, Theorem~\ref{thm:tworeset}
and Proposition~\ref{prop:tworegimes} settle age recurrence at
$\ell=2$, while a common-regime argument remains outstanding.
For fixed protected $\ell\ge3$, neither a uniform finite-age
reset nor full convergence has been established here.

\section{The adopted thesis convention: protection and minimum ages}
\label{sec:thesis}

The fixed-length proof target from this point uses both rules
suggested by the author's thesis~\cite[Figure 2.3, p.~19]{cyclonthesis}:
retain every valid received identity, and keep the smaller age
when two descriptors have the same identity.
This is a source-supported choice of CYCLON semantics, rather
than a repair claimed to be absent from the original design.
Retain the independent uniform scheduler, atomic exchanges,
mandatory partner removal, uniform maximum-size subset choices,
and underfull-cache conventions of Section~\ref{sec:mandatorymodel}.
Write $\widehat P_\ell$ for the resulting kernel.
The thesis's periodic scheduler is not silently substituted for
the independent scheduler used in the positive theorems.

Precisely, first form both buffers and remove the initiator's
partner as before. Ignore incoming self descriptors. Protect
every resident identity present in the received buffer and
replace its age by the minimum of its resident and incoming ages.
Insert new valid identities, filling vacancies and then replacing
only unprotected sent residents; newly inserted identities are
not victims. Choose any required victim set uniformly among
these eligible residents. Lemma~\ref{lem:receiveprotect} proves
feasibility, since minimum-age merging does not change the
identity counts in that argument.

\begin{proposition}[Lossless exchange of identity knowledge]
\label{prop:thesisunion}
For every shuffle length and capacity, each exchange $p\to q$
under $\widehat P_\ell$ preserves
\[
 K_p'\cup K_q'=K_p\cup K_q,\qquad K_v=C_v\cup\{v\},
\]
where caches in this formula denote identity sets.
In particular every weak component is preserved, and the
disconnection constructions of Section~\ref{sec:largerobstruction}
are excluded.
\end{proposition}
\begin{proof}
Every removed nonendpoint identity was sent to the opposite
endpoint and is retained there by received-identity protection.
The endpoint identities remain in their own knowledge sets.
Conversely, every new cached identity came from one of the
two old caches or is an endpoint's self identity.
This proves union equality. The edge-transfer proof of
Lemma~\ref{lem:receiveprotect} gives weak connectivity.
Minimum-age merging changes neither assertion about identities.
\end{proof}

\begin{theorem}[Singleton convergence survives minimum-age merging]
\label{thm:thesisone}
For $\ell=1$, the full-state convergence and conditional local
uniformity conclusions of Theorems~\ref{thm:mandatoryone}
and~\ref{thm:generalone} hold for $\widehat P_1$.
For $c\ge2$ the same canonical atom, time bound $T$, and
minorization $\gamma$ in Theorem~\ref{thm:generalone} apply.
At capacity one the stationary law remains uniform on the
zero-age rooted spanning trees.
\end{theorem}
\begin{proof}
The kernels are not identical: refreshing a duplicate can
change a later oldest-partner choice. Nevertheless the clearing
proof of Theorem~\ref{thm:agereset} is partner-agnostic.
A first contact makes the passive responder know the active
peer; a later contact can return that self identity and shrink
the active cache. Refreshing its passive copy to age zero
does not change either fact. Thus $n+c-1$ consecutive
activations still empty a chosen cache.

A clearing block starting with all ages zero again ends with
all ages zero: its active cache is empty, and its passive
inputs are fresh self descriptors of age zero.
The initial clearing sweep therefore still reaches the same
zero-age forward family. The source-separation and canonical
extension arguments use identities, cache cardinalities,
clearing completion, and zero-age ties at block boundaries.
None requires preserving a positive age difference.
Their invariants and event-probability bounds thus persist,
giving the same common-atom construction for $c\ge2$.

At capacity one, minimum-age merging can remove positive-aged
copies earlier. The reduction to the zero-age rooted-tree
class remains valid, and the transition kernel on that class
is unchanged. The stationary formula follows as before.
Finally apply Lemma~\ref{lem:atom} and
Theorem~\ref{thm:stationary}. Higher-capacity global stationary
age laws need not equal those of the no-refresh kernel.
\end{proof}

\begin{theorem}[Two-entry finite-age reset with minimum-age merging]
\label{thm:thesistwo}
For $2\le c\le n-1$ and $\ell=2$, use the adopted kernel
$\widehat P_2$. Put
\[
 L=n(n-1)(nc+2)(2n+c),\qquad
 B=(nc+1)L,\qquad \delta=(nc^5)^{-L}.
\]
Let $\widehat F_B$ be the finite family of weakly connected
states with maximum age at most $B$. Then, from every weakly
connected state with arbitrary finite initial ages,
\[
 \widehat P_2^L(x,\widehat F_B)\ge\delta,\qquad
 \mathbb E_x\tau_{\widehat F_B}\le L/\delta.
\]
Thus there are finitely many closed stationary regimes, and
the Poisson-time process converges uniformly at a geometric
rate to an initial-state-dependent mixture of their laws.
Uniqueness and conditional uniformity are not asserted here.
\end{theorem}
\begin{proof}
Directly tracking one initial copy's increasing age would fail:
an incoming younger duplicate can reset its resident age.
Instead choose a gap in the finite collection of initial ages.
There are at most $nc$ initial descriptors and $nc+1$ disjoint
intervals
\[
 (jL,(j+1)L],\qquad j=0,\ldots,nc.
\]
One interval contains no initial age. Write $a=jL$ for its
left endpoint. Color initial descriptors of age at most $a$
young and those of age greater than $a+L$ old.
Color new self descriptors young; copies inherit the color
of the descriptor retained by minimum-age merging.

During the first $L$ events, every young descriptor has age
at most $a+t\le a+L$, and every old descriptor has age
strictly greater than $a+L$. Old--old minimum merging
preserves this lower bound. A mixed merge necessarily keeps
the young copy. Thus a holder of an old tracked descriptor
always selects an old partner throughout this horizon.

Erase old colors identity by identity. For a fixed identity
$r$, initially immunize every current holder of young $r$.
Keep these peers passive during this identity's stage.
Whenever an old $r$ remains outside the immune set, use the
pair-retaining routing policy of Lemma~\ref{lem:tworouting}.
Use the present young/old colors in place of its auxiliary
fresh/old colors. If a received young $r$ replaces the
tracked old copy, stop the routing block immediately and
add its holder to the immune set.

Otherwise the same young-count argument applies. Both the
tracked extra and the selected partner at the initiator are
old. At the responder, choose victims that are old or are
young with identity outside the initiating cache and its
owner. When exactly one of the pair is known, exclude that
known identity from victim candidates first. Among the other
candidates at most $|C_p|-2$ can be protected young identities,
so the counting argument in Lemma~\ref{lem:tworouting}
still supplies enough victims. A young victim is copied to
a new young location at the initiator; minimum-age merging
can only add further young locations. A forced migration
requires that the responder knew neither member of the pair,
so its acceptance of the fresh initiator gives a strict
increase in the global young count.

All these choices obey received-identity protection.
There are at most $nc$ migrations, and a fixed-holder
segment still has at most $2n+c$ events: each first partner
learns the pair, and a repeat replies with that pair and
shrinks the initiating cache. Hence a routing block either
ends by young conversion or selects $r$ within
$(nc+1)(2n+c)$ events. The age-gap separation, rather than
a particular copy's age trajectory, justifies every use
of an old selected partner in this argument.

At a terminal contact $h\to r$, peer $h$ becomes $r$-free:
the owner cannot reply with its own identity. The owner
retains or inserts fresh $h$ by received protection.
Activate only owner $r$, retaining target $h$ as its extra.
Any contact with a nonimmune responder now inserts or
refreshes fresh $r^0$, immediately adding that responder
to the immune set. Otherwise the responder already knows
young $r$; prescribe a reply containing $r$, and $h$ too
if known. The self reply leaves at most one usable entry,
so the mandatory vacancy suffices and $h$ survives.
Every first partner learns the pair, while repeats shrink
the active cache. Within $n+c-1$ events the callback thus
adds a new immune peer, at latest by contacting $h$.
Previously immune peers retain their young $r$ because
every packet to them contains $r$, and mixed merging
cannot replace young by old.

Every iteration adds an immune holder; at most $n-1$
iterations per identity erase all old copies of that
identity. Across all identities, the event bound is
\[
 n(n-1)\bigl[(nc+1)(2n+c)+(n+c-1)\bigr]\le L.
\]
No erased old color can subsequently reappear.
The age-gap separation holds throughout these first $L$
events, so this bound closes the finite-horizon construction.
Pad shorter paths by arbitrary valid events. All remaining
descriptors are young, hence their final ages are at most
$a+L\le(nc+1)L=B$.

The same conservative probability bound $(nc^5)^{-1}$
per prescribed event as in Theorem~\ref{thm:tworeset}
applies. Connectivity holds by Proposition~\ref{prop:thesisunion}.
This proves the finite-family minorization and its geometric
hitting bound. Apply the Poissonization argument of
Proposition~\ref{prop:tworegimes} to obtain the stationary-regime
statement. Young initial lineages may survive this construction;
the conclusion is a uniform numerical age bound, not complete
initial-lineage extinction.
\end{proof}

\subsection{A general finite-age reset for every shuffle length}

The next construction removes the remaining age-recurrence gap
for the adopted algorithm. The controlled choices below describe
positive-probability outcomes of its existing random choices.
They do not require peers to observe global state or change the
algorithm of Section~\ref{sec:target}.

\begin{lemma}[A decreasing color potential for every $\ell\ge2$]
\label{lem:generalcolor}
Color descriptors young or old, color fresh self descriptors
young, and let mixed duplicate merges keep young. Assume every
peer holding an old descriptor has an old maximum-age descriptor,
so an old oldest partner can be chosen. Then, whenever any old
descriptor remains, legal controlled choices strictly decrease
lexicographically the nonnegative integer triple
\[
 \Phi=\left(nc-F,\ V,\ W\right),\qquad
 F=\sum_i f_i,\quad V=\sum_i f_i^2,\quad
 W=\sum_i o_i\ell^{c-f_i},
\]
where $f_i$ and $o_i$ are the young and old cache counts at
peer $i$. The coordinates lie in the ranges
\[
 0\le nc-F\le nc,\quad 0\le V\le nc^2,\quad
 0\le W\le nc\ell^c.
\]
\end{lemma}
\begin{proof}
Choose an initiator $p$ maximizing $f_p$ among peers with
$o_p>0$, and choose a legal old partner $q$.
Write $P,Q$ for their pre-update identity sets, $Y_p,Y_q$
for their young identity sets, $a=|P|$, $b=|Q|$, and
\[
 k=\min(\ell-1,a-1),\qquad m=\min(\ell,b).
\]
Choose the extra set $S\subseteq P\setminus\{q\}$ of size
$k$, including one member of $Y_p\setminus Y_q$ if that set
is nonempty. Such an extra exists when needed, since $q$
is old and $\ell\ge2$. Put $S_Y=S\cap Y_p$,
$S_O=S\setminus S_Y$, and $U=\{p\}\cup S$.
The fresh self identity $p$ is young in the outgoing packet.

The responder needs
\[
 r_q=\bigl(b+|U\setminus Q|-c\bigr)_+
\]
victims. Forbid victims in
\[
 E=(Q\cap U)\cup\bigl(Y_q\cap(P\cup\{p\})\bigr).
\]
The first part is required by received-identity protection;
the second prevents loss of a young copy already young at $p$.
If $d_0=|Q\cap U|$, then
\[
 |E|\le d_0+(a-1-k),\qquad
 b-|E|\ge b-d_0-a+1+k.
\]
The extra term counts only unsent identities of $P$, since
shared sent identities are already in $Q\cap U$, and
$q\notin Q$. When $r_q>0$,
$r_q=b+k+1-d_0-c$, so $b-|E|\ge r_q$ by $a\le c$.
Choose $r_q$ victims $D_q\subseteq Q\setminus E$.

Require the reply $T\subseteq Q$ to contain
\[
 D_q\cup A,\qquad
 A=(S_Y\cap Q)\cup(\{p\}\cap Q).
\]
These two sets are disjoint and $A\subseteq Q\cap U$.
If $r_q>0$, their total size is at most
$r_q+d_0=b+k+1-c\le\min(\ell,b)=m$.
If $r_q=0$, $|A|\le d_0\le m$ instead.
Fill the reply arbitrarily to size $m$.
Every chosen responder victim was sent and is unprotected,
so this is a legal joint outcome of reply and victim choices.
At the initiator, choose old eligible victims before young
ones, using only the already permitted victim sets.

We first prove that $F$ cannot decrease, with strict increase
if an outgoing young identity was not already young at $q$.
Let
\[
 d_Y=|S_Y\cap Q|,\quad \sigma=\mathbf1_{\{p\in Q\}},\quad
 e_O=|S_O\cap T|.
\]
At most $m-d_Y-\sigma-e_O$ reply identities are usable
new entries at $p$: the counted entries are duplicates or
the self identity. There are $c-a+1$ mandatory vacancies
and $k-|S_Y|-e_O$ eligible old extras. Old-first victim
selection therefore loses at most
\[
 \bigl(|S_Y|-d_Y-\sigma\bigr)_+=(j_Y-1)_+,
 \qquad
 j_Y=\mathbf1_{\{p\notin Q\}}+|S_Y\setminus Q|,
\]
young copies at $p$. Here we used
$c-a+1+k\ge\ell\ge m$.
Every young victim at $q$ is absent from $P\cup\{p\}$,
so its reply inserts a new protected young copy at $p$.
These young losses are compensated one for one.
The outgoing packet additionally creates $j_Y$ young copies
at previously absent identities, and converts $g$ old copies
at $q$ to young, where $g$ counts its outgoing young
identities previously old there. Consequently
\[
 F'-F\ge j_Y+g-(j_Y-1)_+\ge0,
\]
with strict inequality if $j_Y+g>0$.

Suppose now that $F'=F$. Then $p\in Y_q$ and every
young member of $S$ belongs to $Y_q$. By the prescribed
choice of $S$, this implies $Y_p\subseteq Y_q$, hence
\[
 f_q\ge f_p+1.
\]
In particular $q$ has no old descriptor: otherwise its
young count would contradict the maximal choice of $p$.
No young copy is lost at $p$. All young changes consist
of $t$ responder victims transferring to new young locations
at $p$, so
\[
 f_p'=f_p+t,\quad f_q'=f_q-t,\quad
 0\le t\le f_q-f_p-1.
\]
For $t>0$ this gives
\[
 V'-V=2t\bigl(t-(f_q-f_p)\bigr)<0.
\]

If also $V'=V$, then $t=0$. The responder was entirely
young and no new young copy is created at $p$, so $p$
receives no new descriptor at all. It loses exactly its
selected old partner. The responder can acquire at most
$\ell-1$ old extras, while the young counts stay unchanged.
Therefore
\[
 W'-W\le-\ell^{c-f_p}
       +(\ell-1)\ell^{c-f_q}<0,
\]
using $f_q\ge f_p+1$.
This proves strict lexicographic decrease in every case.
The coordinate bounds follow directly from capacity $c$.
\end{proof}

\begin{theorem}[Uniform finite-age reset for every fixed shuffle length]
\label{thm:thesisgeneralreset}
Use the fixed adopted algorithm in Section~\ref{sec:target},
with $2\le\ell\le c\le n-1$. Set
\[
 H=(nc+1)(nc^2+1)(nc\ell^c+1),\qquad
 B=(nc+1)H,\qquad
 \delta=(nc\,16^c)^{-H}.
\]
For the finite family $F_B$ of weakly connected aged states
whose maximum age is at most $B$,
\[
 \widehat P_\ell^H(x,F_B)\ge\delta
\]
uniformly over all weakly connected initial states with
arbitrary finite integer ages. In particular,
\[
 \Pr_x\{\tau_{F_B}>kH\}\le(1-\delta)^k,\qquad
 \mathbb E_x\tau_{F_B}\le H/\delta.
\]
The construction does not require occasional shorter shuffles.
\end{theorem}
\begin{proof}
Among the $nc+1$ intervals $(jH,(j+1)H]$, $j=0,\ldots,nc$,
one contains no initial age, since there are at most $nc$
initial descriptors. Put $a=jH$ for that interval.
Color initial ages at most $a$ young and those greater
than $a+H$ old. Fresh self descriptors are young.
For the first $H$ events, young ages are at most $a+t\le a+H$
and old ages remain strictly greater than $a+H$, even after
old--old minimum merging. A mixed merge keeps the young copy.
Thus every oldest partner of a holder of an old descriptor
is old during this horizon.

Apply Lemma~\ref{lem:generalcolor} at every event until
old descriptors are gone. Its integer triple takes at most
$(nc+1)(nc^2+1)(nc\ell^c+1)=H$ values, and strictly decreases
while any old remains. Hence old-color extinction occurs
within $H$ events. Padding a shorter path by arbitrary
valid events cannot recreate old colors. The maximum age
at event $H$ is therefore at most $a+H\le B$.
Connectivity holds by Proposition~\ref{prop:thesisunion}.

Each prescribed event has probability at least
$(nc\,16^c)^{-1}$. A conservative product counts initiator
probability $1/n$, oldest-tie probability at least $1/c$,
the extra and reply subsets, and both eligible victim subsets.
The latter four subset denominators have product at most
$2^{c-1}2^c2^{\ell-1}2^\ell\le16^c$.
Prescribing a compatible reply/victim pair in the proof
does not change their actual order of sampling.
Multiplication along the controlled path gives $\delta$;
padding after extinction is unrestricted.
The Markov property on successive $H$-event blocks proves
the hitting bounds.
\end{proof}

\begin{corollary}[General stationary regimes and exponential age tails]
\label{cor:thesisgeneralregimes}
For every fixed $1\le\ell\le c$, the adopted chain has
finitely many closed communicating classes, each with a
unique stationary probability law. With independent Poisson
clocks, its full state converges in total variation to an
initial-state-dependent mixture of these laws, uniformly
at some geometric rate in elapsed time.
For $\ell\ge2$, every stationary law $\pi$ satisfies
\[
 \pi\{M>B+(k-1)H\}\le(1-\delta)^k\quad(k\ge1),
 \qquad
 \mathbb E_\pi M\le B-H+H/\delta,
\]
and, for $0\le\theta<-H^{-1}\log(1-\delta)$,
\[
 \mathbb E_\pi e^{\theta M}
 \le\frac{\delta e^{\theta B}}
 {1-(1-\delta)e^{\theta H}}.
\]
These are age-control and stationary-regime statements;
they do not establish a single regime or conditional sampling
uniformity for fixed $\ell\ge2$.
\end{corollary}
\begin{proof}
For $\ell=1$ use Theorem~\ref{thm:thesisone}.
Otherwise apply the finite-family and Poissonization argument
of Proposition~\ref{prop:tworegimes} to
Theorem~\ref{thm:thesisgeneralreset}.
In the last $k$ blocks of $H$ events, any successful reset
leaves maximum age at most $B+(k-1)H$ at observation.
Failure of all blocks has probability at most $(1-\delta)^k$.
At stationarity this bounds $M$ by the stochastic upper bound
$B-H+HG$, where $G$ is geometric on $\{1,2,\ldots\}$ with
success probability $\delta$. Its moments give the formulas.
\end{proof}

\subsection{A structural reduction for two-entry shuffles}

Call $a$ a \emph{universal zero-age anchor} if every
$v\ne a$ stores $a^0$. Its own cache need not be empty,
and the graph induced by $V\setminus\{a\}$ need not be
connected. This is a useful intermediate configuration,
not an additional initialization assumption.

\begin{lemma}[Removing a universal anchor from controlled paths]
\label{lem:anchorreduction}
Let $2\le\ell\le c$ and suppose $a$ is a universal
zero-age anchor. Freeze $a$ and remove its descriptor from
every other cache. Every finite path of the adopted algorithm
on $W=V\setminus\{a\}$ with capacity $c-1$ and shuffle
length $\ell-1$ can be simulated by a positive-probability
path of the original algorithm, preserving every anchor
copy at age zero and leaving $C_a$ unchanged. Reduced
activations of empty caches are omitted.
\end{lemma}
\begin{proof}
For a nonempty reduced cache at $p$, choose its prescribed
oldest residual partner $q$. This is a legal actual oldest
choice: after the increment every residual age is at least
one, while the anchor's age is exactly one.
Include $a$ among the outgoing extras and require $a^0$
in the reply. The remaining buffer sizes are exactly
\[
 \min(\ell-2,|C_p\setminus\{a\}|-1),\qquad
 \min(\ell-1,|C_q\setminus\{a\}|),
\]
as required for the reduced exchange. At $p$ the returned
$a^0$ refreshes its incremented anchor; at $q$, its resident
$a^0$ survives the incoming $a^1$ by minimum merging.
The anchor is protected at both endpoints. The other
duplicate merges, vacancy counts, and eligible victim sets
are exactly those of capacity $c-1$. All prescribed subset
and victim choices have positive probability.
If a reduced cache is empty, its reduced activation is a
no-op and is skipped, rather than activating its actual
singleton cache toward $a$. Thus this is a path simulation,
not an equality of the two randomly scheduled kernels.
\end{proof}

\begin{proposition}[A universal anchor suffices when $\ell=2$]
\label{prop:anchorcanonical}
For the adopted algorithm with $\ell=2$ and any $c\ge2$,
every state with a universal zero-age anchor reaches a
zero-age canonical state $H(V,c)$ by a finite
positive-probability path. The order defining $H$ may be
chosen arbitrarily. In particular all such labelled
canonical states communicate.
\end{proposition}
\begin{proof}
First keep the anchor $a$ fixed. By
Lemma~\ref{lem:anchorreduction}, the proved singleton
constructions can canonicalize each residual weak component
on $W$, with capacity $c-1$ and any chosen component root.
For capacity one use Theorem~\ref{thm:mandatoryone}; for
higher capacities use Theorem~\ref{thm:thesisone} and its
constructive proof, which permits excess slots on smaller
active populations. Singleton components need no exchanges.

We first empty $C_a$, without any assumption on its ages.
If it is nonempty, choose an oldest $q\in C_a$ and
canonicalize $q$'s residual component with $q$ as root,
while freezing $a$. Then $C_q=\{a^0\}$ and $C_a$ is
unchanged. Activate $a$ toward $q$. Its only reply is the
self identity $a$, so $C_a$ loses $q$ and receives nothing.
At $q$, fresh $a^0$ is protected and at most one extra
descriptor is inserted, which fits since $c\ge2$.
Every other anchor stays at age zero. Repeat at most $c$
times to empty $a$.

Now canonicalize each residual component at age zero.
Roots $u,v$ of two different components have
$C_u=C_v=\{a^0\}$. The following four exchanges merge
their residual components and restore the empty owner and
all anchor copies:
\[
\begin{array}{c|ccc}
 &C_a&C_u&C_v\\\hline
 \text{initial}&\varnothing&\{a^0\}&\{a^0\}\\
 u\to a&\{u^0\}&\varnothing&\{a^0\}\\
 v\to a&\{u^0,v^0\}&\varnothing&\{u^0\}\\
 a\to v&\{u^0\}&\varnothing&\{u^0,a^0\}\\
 a\to u&\varnothing&\{a^0\}&\{u^0,a^0\}
\end{array}
\]
At the third exchange choose $v$ in the oldest tie.
Its reply $u^0$ refreshes the initiator's $u^1$, and
the responder keeps $u^0$ when receiving $u^1$.
All displayed caches have at most two entries, so the
table is legal for every $c\ge2$; other peers are untouched.
The residual edge $v\to u$ merges the components.
Repeat with the same root $u$ until $W$ is connected,
then use the reduced singleton construction to obtain
$H(W,c-1)$. Together with the empty anchor $a$, this is
exactly $H(V,c)$ with $a$ ordered last.

To change the anchor to any desired $b\ne a$, first
canonicalize $W$ with $b$ last. Then $C_b=\{a^0\}$,
and every other peer in $W$ contains $b^0$.
One exchange $b\to a$, with empty reply, gives
$C_b=\varnothing$ and $C_a=\{b^0\}$ and leaves all
other $b^0$ copies unchanged. The result is the canonical
state with the final two labels $b,a$ interchanged.
With $b$ fixed as anchor, the reduced singleton
construction sets all remaining labels in any desired
order. This proves the assertion for every order.
\end{proof}

\begin{proposition}[A universal identity need not start fresh]
\label{prop:anyanchor}
For fixed $\ell=2$, suppose every $v\ne a$ stores identity
$a$, with arbitrary finite resident ages and arbitrary
$C_a$. A positive-probability path of at most
$2(n-1)(n+c)$ exchanges makes all those copies age zero.
Consequently Proposition~\ref{prop:anchorcanonical} applies
without an initial age-zero requirement on the anchor.
\end{proposition}
\begin{proof}
Choose a holder $p\ne a$ whose copy is not zero and
activate only $p$. Until it contacts $a$ or receives $a^0$,
always include $a$ as the extra descriptor. Every nonowner
responder $q$ knows $a$; require it to return $a$, together
with $p$ if it already knows $p$, or any other entry
otherwise. A singleton reply contains only $a$.
At most one incoming identity is new at $p$, so it fits
the mandatory vacancy and cannot evict its anchor.
At the responder the incoming identities $p,a$ are
protected. Thus every visited responder permanently knows
$p$ during this block, and a repeated visit returns only
duplicates or self to $p$, shrinking its cache by one.
There are at most $n-1$ first visits. Within $n+c$ events
$p$ either becomes a zero holder or contacts owner $a$.
No other zero holder is lost, since every such responder
receives an anchor descriptor and keeps its resident zero.

In the latter case $p$ loses $a$, and the owner receives
and protects fresh $p^0$. All other peers still store $a$.
Now activate only $a$, always using $p$ as extra while its
partner differs from $p$. Require every such responder to
return $a$ and, if known, $p$. The incoming self identity
$a$ is ignored at the owner, so at most one usable new
entry fits its vacancy and target $p$ cannot be evicted.
Each responder learns $p$ on its first visit; repeated
visits return $a,p$ and shrink the owner's cache.
Within $n+c$ events the owner therefore contacts $p$,
which receives protected $a^0$. During this block all other
anchor holders receive fresh $a^0$ whenever touched.
Universality is restored and at least one additional
holder is zero. At most $n-1$ iterations prove the bound.
\end{proof}

\subsection{Uniform contact across a frozen boundary}

\begin{theorem}[A uniform first-contact bound for $\ell=2$]
\label{thm:firstcontact}
Use the adopted algorithm with $\ell=2$.
Fix a nonempty proper set $G\subset V$, put $S=V\setminus G$
and $s=|S|$, and suppose some cache in $S$ contains an
identity $r\in G$. Define
\[
 \begin{split}
 D_s&=s(s+c),\qquad J_s=(sc+1)(2s+c),\\
 T_s&=D_s(J_s+1)+J_s,\qquad
 \eta_s=(nc^5)^{-T_s}.
 \end{split}
\]
There is a positive-probability path, using only initiators
in $S$, that contacts $G$ within $T_s$ events.
Under the original uniform scheduler, if $\tau$ is the
first exchange between $S$ and $G$ in either direction,
\[
 \Pr_x\{\tau\le T_s\}\ge\eta_s,\qquad
 \Pr_x\{\tau>kT_s\}\le(1-\eta_s)^k.
\]
No bound on initial ages or on the frozen caches is needed.
The constructed path always activates an $r$ holder,
sends $r$ as extra when its partner is not $r$, and
requires $r$ in the reply whenever the responder has it.
The total number of $r$ holders cannot decrease along
this path except at its final contact if that partner
is owner $r$.
\end{theorem}
\begin{proof}
We construct a path while keeping $G$ inactive; any
contact with $G$ terminates the construction successfully.
Before such contact, identity $r$ remains in at least one
$S$ cache by received-identity protection and the pair-union
identity in Proposition~\ref{prop:thesisunion}.
No fresh $r$ can be generated inside $S$.
Its minimum outside age $m$ is therefore nondecreasing.

First we show that $m$ can be raised by at least one,
or $G$ contacted, within $D_s$ events. Pick a holder $p$
of $r^m$ and activate only $p$. Until success or loss of
this minimum copy, send $r$ as extra. If its actual oldest
partner is in $G$, we have succeeded. If a responder $q$
does not store $r^m$, require $r$ in the reply if it is
known there. The increment at $p$ destroys that
minimum copy and none is created: the outgoing copy has
age $m+1$, and the reply has no $r^m$.
If $q$ does store $r^m$, require it to return this
descriptor, plus $p$ when known or one other residual
entry otherwise. The initiator's copy is refreshed to $m$;
the responder's is protected and stays at $m$.
There is at most one usable new reply, so $p$ has no
extra eviction. Every such responder learns and thereafter
retains $p$. A repeat returns $r^m,p$, shrinking $C_p$.
Thus within $s+c$ events either a minimum copy is lost,
or $p$ is left with the singleton $r^m$ and contacts $G$.
There are at most $s$ minimum copies. Repeating this
operation proves the bound $D_s=s(s+c)$.

Repeat the minimum-raising construction $J_s+1$ times.
Unless $G$ has already been contacted, every outside
$r$ copy now has age greater than $J_s$.
Mark all descriptors currently present auxiliary old,
and all subsequent fresh self descriptors auxiliary young.
For this proof-only marking, a duplicate is young whenever
\emph{either} input is young, regardless of which supplies
the minimum numerical age. This marking is sound: the
minimum age is at most the young input's age, so after
$t$ further events every young descriptor has age at most
$t$. No $r$ descriptor becomes young, since its owner is
frozen. For the next $J_s$ events every selected oldest
partner of an $r$ holder is therefore auxiliary old.

Apply the pair-retaining routing choices of
Lemma~\ref{lem:tworouting}, stopping on contact with $G$.
Its victim argument still applies under minimum merging:
exclude incoming identities $p,r$ from responder victims,
and compensate every young victim by its copy at a new
initiator location. There are enough victims that are old
or have young identities outside $C_p\cup\{p\}$, since
$C_p$ contains the old target $r$ and old selected partner.
The initiator can evict only its old extra $r$, and the
proof-only young marking cannot be destroyed by a duplicate
merge. Thus the global young count never decreases.
If $r$ must migrate, the responder knew neither $p$ nor
$r$, so accepting fresh $p^0$ strictly increases that count.
There are at most $sc$ migrations. With a fixed holder,
first visits establish the pair $p,r$ and repeats reply
with that pair and shrink its cache, so each segment
has at most $2s+c$ events. Therefore at most
$(sc+1)(2s+c)=J_s$ routing events contact $G$.
Together with minimum raising this proves the bound $T_s$.

Each prescribed event has probability at least
$(nc^5)^{-1}$, by the two-entry choice bound in
Theorem~\ref{thm:thesistwo}. The denominator
is $n$, even though the witness path uses only $S$.
The holder-count claim follows directly from the packets:
an existing responder copy is protected, and its reply
protects the initiator's copy; a new responder copy
compensates any initiator loss. Only owner $r$ cannot
retain or return its own identity at the terminal contact.
This proves the one-block bound. Before any cross exchange,
$G$-internal events cannot change the outside caches, and
$S$-internal events preserve some outside $r$ copy.
The same hypothesis therefore holds at every block start
on $\{\tau>kT_s\}$. The Markov property gives the tail bound.
\end{proof}

\begin{corollary}[Every cut is crossed repeatedly]
\label{cor:repeatedcontact}
For $\ell=2$, every fixed nontrivial partition of a weakly
connected overlay is crossed infinitely often almost surely.
Its successive waiting times have a uniform geometric
block-tail bound depending only on $n,c$.
\end{corollary}
\begin{proof}
Connectivity supplies a directed edge across each cut.
Apply Theorem~\ref{thm:firstcontact} on the side holding
that descriptor, interchanging the two sides if necessary.
Use the largest $T_s$ over $1\le s<n$ and the corresponding
smaller common probability bound. Connectivity is preserved
after every exchange, so the argument restarts after each
crossing and at every subsequent observation time.
\end{proof}

\begin{proposition}[A prescribed peer can reach the whole overlay]
\label{prop:rootreachability}
For fixed $\ell=2$, every weakly connected initial state
can reach a state in which every vertex is reachable by
a directed cache path from any prescribed peer $a$.
The construction takes at most $\sum_{s=1}^{n-1}T_s$
controlled events, with $T_s$ from
Theorem~\ref{thm:firstcontact}.
When $c=2$, at most $2(n-1)$ controlled exchanges suffice,
independently of all initial ages.
\end{proposition}
\begin{proof}
Let $R$ be the vertices currently reachable from $a$.
If $R\ne V$, no cache edge points from $R$ to its
complement, while weak connectivity supplies a cache edge
from the complement into $R$.
Freeze $R$ and use Theorem~\ref{thm:firstcontact} until
the first exchange $p\to q$ with $p\notin R$, $q\in R$.
Before that exchange every $R$ cache is unchanged.
The responder retains fresh $p$, and every old responder
edge $q\to t$ that is evicted is copied into $C_p$.
For $t\ne p$ the copy is protected at $p$; for $t=p$
the incoming fresh $p$ already preserves the edge.
Consequently every old directed path from $a$ can be
rerouted through $q\to p\to t$ if needed, and $p$
is now reachable. Thus $R$ strictly grows. Repeat at
most $n-1$ times. Its complement size strictly decreases,
so the sum of the applicable $T_s$ bounds the total.

For $c=2$ the first-contact step costs at most two events.
Choose a boundary descriptor $q\in C_p\cap R$.
If the actual oldest partner of $p$ is in $R$, one
event suffices. Otherwise its cache is exactly
$\{u,q\}$ with oldest partner $u\notin R$.
Sending $p^0,q$ forces $C_u$ to become exactly
$\{p^0,q^b\}$ for some $b\ge0$ by received protection
and minimum merging. Activate $u$ and choose $q$,
which is oldest or tied after its increment.
The first event was entirely outside $R$; the second
strictly enlarges $R$ by the preceding argument.
\end{proof}

\begin{corollary}[A precise remaining accessibility criterion]
\label{cor:anchorcriterion}
For fixed $\ell=2$, it suffices to prove that every member
of the finite reset family $F_B$ in
Theorem~\ref{thm:thesisgeneralreset} can reach \emph{some}
universal identity, with no restriction on its ages.
That additional accessibility
statement would imply a uniform common-atom bound, unique
stationarity, full-state geometric convergence of the
embedded event chain, and the conditional sampling target
of Section~\ref{sec:target}.
\end{corollary}
\begin{proof}
By Propositions~\ref{prop:anyanchor} and
\ref{prop:anchorcanonical}, every such anchor reaches the
same prescribed labelled $H(V,c)$.
There are finitely many reset states, so the chosen paths
have a finite maximum length $J$ and a positive minimum
probability $\varepsilon$. Pad them to length $J$ by
activating the empty root of $H$. The general age reset
then gives
\[
 \widehat P_2^{H+J}(x,\{H(V,c)\})
 \ge\delta\varepsilon\qquad\text{for every }x.
\]
Apply Lemma~\ref{lem:atom}. Relabelled canonical states
give positive stationary availability of every $b\le c$,
and Theorem~\ref{thm:stationary} gives conditional uniformity.
This criterion does not assume access to an anchor for
the adopted algorithm. Theorem~\ref{thm:thesisc2} below
closes the criterion when the cache capacity is also two.
\end{proof}

\subsection{Full uniform sampling with two-entry caches}
\label{sec:thesisc2}

Throughout this subsection the \emph{cache capacity and shuffle
length are both two}: $c=\ell=2$, with arbitrary $n\ge3$.
All exchanges use the fixed algorithm of Section~\ref{sec:target}.
The preceding structural lemmas now suffice to close this case.
The proof uses extremal states within a closed communicating
class; it does not assume that peers can detect such states.

\begin{lemma}[Source-only gain of a fresh copy]
\label{lem:c2sourcegain}
Fix $a$, and let $I=\{v\ne a:a^0\in C_v\}$.
If a vertex outside $I\cup\{a\}$ is reachable from $a$
by a directed cache path, a positive-probability path of
at most $(n-1)(n+1)$ activations of $a$ increases $|I|$.
No previously existing $a^0$ copy is lost.
\end{lemma}
\begin{proof}
Every source-only exchange delivers protected fresh $a^0$
to the responder, so $I$ cannot lose a member. Stop if a
responder is outside $I$, since that immediately gives a gain.
Until then every responder stores $a^0$ and returns it;
at most one other reply identity is usable at $a$.
It fits the mandatory vacancy, so an extra retained at
the initiator cannot be evicted.

Choose a shortest directed path from $a$ to outside
$I\cup\{a\}$, with first head $z$. Its internal vertices
belong to $I$. While contacting any other partner, send
$z$ as extra. An off-path responder learns and retains
$z$; if revisited it returns exactly $a,z$, causing the
source cache to shrink. The source cache never grows,
and $z$ is retained. Thus, within $n+1$ events, either
a gain occurs or a path vertex is contacted.

When a path vertex is contacted, its reply transfers its
next path head to $a$, strictly shortening the path.
This holds both for $z$ and for a later internal vertex:
its cache contains $a^0$ and that next head. The remaining
suffix is unchanged, since neither endpoint of the exchange
is on the suffix. Off-path contacts had left the original
path intact. Repeat with the shortened path. Its initial
length is at most $n-1$, proving the bound.
\end{proof}

\begin{lemma}[Extremal fresh-copy states]
\label{lem:c2extremal}
In any closed communicating class, fix identity $a$ and
maximize $k=|\{v\ne a:a^0\in C_v\}|$ over its states.
Then $k\ge1$, and every maximizing state has the following
properties:
\begin{enumerate}
\item Every vertex directed-reachable from $a$, other
than $a$ itself, belongs to $I=\{v\ne a:a^0\in C_v\}$.
\item There is no positive-age descriptor for $a$ anywhere.
Thus $I$ is exactly the set of peers storing identity $a$.
\end{enumerate}
Moreover a maximizing state with $C_a=\varnothing$ exists
in the same closed class.
\end{lemma}
\begin{proof}
The integer maximum is attained because its possible
values are $0,\ldots,n-1$. The first property follows from
Lemma~\ref{lem:c2sourcegain} and closedness.

For the second, suppose $p$ stores $a^A$ with $A>0$.
If its oldest partner is $a$, contact $a$. No old zero
holder is touched, and the owner now caches the nonzero
holder $p$. Lemma~\ref{lem:c2sourcegain} increases $k$.
Otherwise $C_p=\{a^A,u^B\}$ with oldest partner $u$.
The request $\{p^0,a^{A+1}\}$ forces the responder cache
to be $\{p^0,a^b\}$. If $u$ previously stored $a^0$,
it keeps that zero and its full reply also gives $p$ a
zero, a contradiction. Otherwise $b>0$ and neither $p$
nor $u$ is a zero holder. Activate $u$: its uniquely
oldest partner is $a$. The owner learns fresh $u^0$,
with no old zero holder affected, and again the source-only
gain lemma contradicts maximality. This proves both properties.
They also imply $k\ge1$: if $k=0$, the absence of positive
copies means no peer caches $a$. Weak connectivity then
forces $C_a$ to be nonempty, and an owner activation
creates a protected $a^0$ copy, a contradiction.

Let $R$ be the vertices reachable from $a$ in a maximizing
state. Its caches contain only identities in $R$, and
every nonowner in $R$ stores $a^0$. If $|R|\ge3$, freeze
the complement and apply Proposition~\ref{prop:anchorcanonical}
within $R$, with $a$ as the final root. All old $R$
nonowners end with $a^0$ and the owner empty; all outside
caches are untouched. If $|R|=1$ the owner is already
empty. If $R=\{a,b\}$, then $C_b=\{a^0\}$ and the
single exchange $a\to b$ empties $a$ without changing
$b$'s zero. These paths preserve $k$ at their endpoints
and stay in the closed class.
\end{proof}

\begin{lemma}[Four-exchange fresh-copy gain]
\label{lem:c2infection}
Suppose $C_a=\varnothing$, $C_p=\{a^0,x^A\}$,
and $x$ does not store $a^0$. Then four legal exchanges
increase the number of $a^0$ holders, preserving all
previous such holders at the endpoint:
\[
 p\to x,\qquad x\to a,\qquad a\to p,\qquad x\to p.
\]
\end{lemma}
\begin{proof}
The first partner $x$ is oldest or tied with $a^0$.
Its received pair forces $C_x=\{p^0,a^1\}$: if an
$a$ descriptor was already there, its age was a positive
integer and minimum merging with one gives one.
The initiator $p$ cannot retain $x$ after this exchange,
since it removes $x$ and a responder never returns itself.
Next $x$ necessarily contacts $a$, giving
\[
 C_a=\{x^0,p^1\},\qquad C_x=\{p^1\}.
\]
The next owner activation necessarily contacts $p$,
sending $a^0,x^1$, so $C_p=\{a^0,x^1\}$ exactly.
Finally singleton $x$ contacts $p$ and receives $a^0$
and its own identity. It retains only $a^0$; at $p$,
incoming fresh $x^0$ refreshes $x^1$ and requires no
eviction. The final caches are
\[
 C_p=\{a^0,x^0\},\qquad C_x=\{a^0\}.
\]
All other old zero holders were untouched.
\end{proof}

\begin{theorem}[Full uniform sampling for $c=\ell=2$]
\label{thm:thesisc2}
For every $n\ge3$, the adopted CYCLON algorithm with
cache capacity two and shuffle length two converges from
every weakly connected initial state with arbitrary finite
integer ages to a unique full-state stationary law $\pi$.
There exist $T=T(n)<\infty$ and $\gamma=\gamma(n)>0$ such that
\[
 \inf_x\widehat P_2^T(x,\{H(V,2)\})\ge\gamma,
 \qquad
 \sup_x\|\widehat P_2^t(x,\cdot)-\pi\|_{\TV}
 \le(1-\gamma)^{\lfloor t/T\rfloor}.
\]
For every observer and $b=1,2$, stationary availability
is positive and its ordered $b$-sample converges to uniform
over $(V\setminus\{i\})_b$ conditional on availability.
No efficient bound on $T$ or $\gamma$ is asserted.
\end{theorem}
\begin{proof}
We first show that every closed communicating class
contains a universal zero-age anchor for a fixed identity
$a$. Fix such a class and the maximum $k$ of
Lemma~\ref{lem:c2extremal}. Among its maximizing states
with $C_a=\varnothing$, let
\[
 I=\{v\ne a:a^0\in C_v\},\qquad
 J=V\setminus(I\cup\{a\}).
\]
Lemma~\ref{lem:c2infection} rules out every cache edge
from $I$ to $J$: it would increase $k$. The extremal
lemma also rules out any descriptor for $a$ in $J$.
Thus each $I$ cache consists of $a^0$ and at most one
residual identity from $I$.

Choose such a state minimizing the number $d$ of weak
components of the residual graph on $I$. This minimum
exists since $1\le d\le k$. Suppose $J$ is nonempty.
Weak connectivity, the empty owner, and absence of
$I\to J$ and $J\to a$ edges imply that a peer
$P\in J$ caches some member of $I$. Inspect an actual
oldest partner of $P$ and keep $P$ frozen for the next
normalization.

By Lemma~\ref{lem:anchorreduction}, residual exchanges
on $I$ simulate capacity-one singleton shuffles, keeping
all anchors at zero and the owner empty. Within each
residual component, the capacity-one theorem makes it
a zero-age rooted star with any prescribed root.
The component partition and $d$ do not change.
In particular the selected root $Q$ has $C_Q=\{a^0\}$;
every other member of its component has
$\{a^0,Q^0\}$. None of $P$'s ages or oldest choices changes.

If $P$'s inspected oldest partner $u$ belongs to $J$,
its other entry is a target $Q\in I$. Normalize $Q$'s
component with root $Q$, and exchange $P\to u$ using
$Q$ as extra. Protection forces
$C_u=\{P^0,Q^b\}$ for some $b\ge0$; choose $u\to Q$
next, which is oldest or tied. Root $Q$ replies only
$a^0$, leaving $C_u=\{a^0,P^1\}$ while $Q$ loses its
zero. Peer $P$ still has no anchor: both its original
endpoints were in $J$ and their pre-exchange caches had
none. The zero count stays $k$, the owner stays empty,
and new holder $u$ has an edge to nonholder $P$.
Lemma~\ref{lem:c2infection} contradicts maximality.

Otherwise $P$'s inspected oldest partner is $Q\in I$.
Normalize its component with root $Q$ and exchange
$P\to Q$. If $P$ is a singleton, it receives $a^0$
while the one incoming fresh $P$ fits alongside the
root's $a^0$, increasing $k$ directly.
Otherwise write $C_P=\{Q,X\}$, with $X\ne a,Q$.
Root $Q$ replies only $a^0$ and is forced to retain
fresh $P$ and the extra $X$, losing its anchor. The new
zero set is exactly $I'=(I\setminus\{Q\})\cup\{P\}$,
and $C_P'=\{a^0,X^{A+1}\}$ for the old age $A$ of $X$.
If $X\in J$, the fresh-copy gain lemma applies immediately.
If $X\in I$ and $Q$'s former residual component had
another member $v$, that member still has
$C_v=\{a^0,Q^0\}$ with $Q$ now a nonholder, giving
the same contradiction.

The only remaining case is that $Q$ was a singleton
residual component and $X$ belongs to another component.
No other $I$ cache points to $Q$. Removing this isolated
vertex and adding new holder $P$ with residual edge
$P\to X$ reduces the component count to $d-1$.
All residual identities still lie in $I'$, the owner
is empty, and the zero count is still maximal. This
contradicts minimality of $d$. Hence $J$ is empty and
$k=n-1$.

Proposition~\ref{prop:anchorcanonical} now sends every
closed class to the same prescribed labelled atom
$h=H(V,2)$. Therefore there can be only one closed
class. The uniform finite-family reset of
Theorem~\ref{thm:thesisgeneralreset} supplies $F_B$, its
horizon $H$, and probability $\delta$.
By the finite-class argument in
Corollary~\ref{cor:thesisgeneralregimes}, every member
of $F_B$ has a finite positive path to a closed class,
and hence to $h$. Choose one such path for each member.
Their lengths have a finite maximum $L$ and their
probabilities a positive minimum $\varepsilon$.
Pad them to length $L$ at $h$'s empty-root self-loop.
Then
\[
 \widehat P_2^{H+L}(x,\{h\})\ge\delta\varepsilon
 \qquad\text{for every }x.
\]
Lemma~\ref{lem:atom} gives the theorem with
$T=H+L$ and $\gamma=\delta\varepsilon$.
Relabelling equivariance gives the same minorization
for every relabelled canonical atom. For any observer,
one such atom gives it two cached identities, proving
positive stationary availability for $b=1,2$.
Theorem~\ref{thm:stationary} proves conditional uniformity.
\end{proof}

\begin{corollary}[Uniform identity caches conditional on size]
\label{cor:c2caches}
Under Theorem~\ref{thm:thesisc2}, for any observer $i$
and $k\in\{0,1,2\}$ having positive stationary probability,
\[
 \Pr_\pi\!\left\{\operatorname{ids}(C_i)=A
                 \mid |C_i|=k\right\}
 =\binom{n-1}{k}^{-1}
 \quad\text{for every }A\subseteq V\setminus\{i\},\ |A|=k.
\]
This describes the stationary identity set, without
asserting independence across peers or successive times.
\end{corollary}
\begin{proof}
The unique stationary law is invariant under every
permutation of peer identities. Permutations fixing $i$
act transitively on all its possible $k$-element cache
sets and preserve the event $|C_i|=k$.
\end{proof}

\subsection{An explicit mixing bound for two-entry caches}

The preceding class argument has a constructive quantitative version.
Its bound is uniform over initial ages and applies to a fixed small
cache capacity with arbitrarily large networks.
It still charges for a prescribed sequence of random outcomes, so
the resulting mixing estimate is very pessimistic.

\begin{theorem}[Explicit mixing from a quadratic controlled witness]
\label{thm:c2explicitmix}
For the adopted algorithm with $c=\ell=2$ and $n\ge3$, put
\[
 T_n=16n^2,\qquad \Gamma_n=(4n)^{-T_n}.
\]
For any prescribed labelled canonical state $h=H(V,2)$,
\[
 \inf_x\widehat P_2^{T_n}(x,\{h\})\ge\Gamma_n,
 \qquad
 \sup_x\|\widehat P_2^t(x,\cdot)-\pi\|_{\TV}
 \le(1-\Gamma_n)^{\lfloor t/T_n\rfloor},
\]
where the infimum permits every weakly connected initial
state with arbitrary finite integer ages.
Consequently
\[
 t_{\rm mix}(\varepsilon)
 \le T_n\left\lceil\Gamma_n^{-1}
                  \log(1/\varepsilon)\right\rceil.
\]
For $0<\varepsilon\le1/2$, this upper bound is
$\exp(O(n^2\log n))\log(1/\varepsilon)$.
The polynomial controlled-path length does not assert a
polynomial mixing time: the probability of the prescribed
path is exponentially small.
\end{theorem}

\begin{proof}
Choose the largest label of the prescribed canonical order
as the owner identity $a$. Write
\[
 I=\{v\ne a:a^0\in C_v\},\qquad k=|I|.
\]
We construct a universal zero-age anchor through at most
$n-1$ strict increases of $k$. Temporary losses during a
macro are permitted; every completed macro strictly
increases its endpoint value.

We first record two elementary services. If $a$ caches
a nonholder target $z$, activate only $a$, retaining $z$
as extra while its partner differs from $z$. At any holder
responder require the full reply, which contains $a$ and,
once learned, $z$. The self identity $a$ is filtered at
its owner, so at most one usable new reply fits the
mandatory vacancy. Thus $z$ survives, the owner's cache
does not grow, first visits establish $z$ at responders,
and repeats shrink the cache. Within $n+2$ events some
nonholder, at latest $z$, receives fresh $a^0$.
All previously zero holders are passive and protected.
If all other $a$ copies were zero at the start, they
remain zero throughout this owner-only service.

Second, suppose $Q$ holds $a^0$ and the owner is frozen.
Activate only $Q$, choosing its residual partner rather
than $a$ in a tie, and send $a$ as extra. Stop before
an exchange to a nonholder if one is encountered.
At every internal holder contact the full reply refreshes
the initiator's anchor to zero, and the passive anchor
stays zero by minimum merging. After deleting the anchor,
these are exactly capacity-one singleton exchanges.
Every visited responder learns $Q$; a repeat returns
$Q$ as self and clears the residual cache.
Within $n$ events $C_Q=\{a^0\}$ or a nonholder is
encountered. The owner's cache stays unchanged.

If a positive-age copy $a^A$ exists at $p$, choose an
actual oldest partner. If it is $a$, one exchange makes
the owner cache the nonholder $p$; the first service gives
a strict $k$ gain within another $n+2$ events.
Otherwise $C_p=\{a^A,u^B\}$ and the request forces
$C_u=\{p^0,a^b\}$. If $u$ was a zero holder, its full
reply gives $p$ a zero and the gain is immediate.
Otherwise $b>0$, so $u$'s next oldest partner is $a$.
After $u\to a$, the owner caches the nonholder $u$,
and the same service gives a gain. This costs at most
$n+4$ events, and no earlier zero holder is lost.
Repeat this rule whenever positive copies remain.
If no $a$ descriptor exists anywhere, weak connectivity
forces a nonempty owner cache; one owner activation
creates a fresh zero. This also counts as a strict gain.

It remains to give an $O(n)$ gain macro when all $a$
descriptors are zero and $k<n-1$. If the owner caches any nonholder, the first service
gives a gain within $n+2$ events. Otherwise, if its cache
is nonempty, inspect an actual oldest partner $Q\in I$.
Freeze the owner and perform the stopped
residual clearing at $Q$. If it clears, exchange
$a\to Q$. The sole reply $a^0$ is self-filtered,
so the owner loses $Q$ and gains no usable entry.
At $Q$, fresh $a^0$ and at most one extra fit capacity
two. Thus all zeros survive, and the owner loses
one cache entry. Repeat if necessary, at most twice.

If instead the stopped clearing encounters a nonholder
$x$, perform $Q\to x$ and then immediately $a\to Q$.
The first event forces
\[
 C_x=\{Q^0,a^1\}.
\]
Although $Q$ may temporarily lose its zero, the second
event restores it with protected fresh $a^0$.
The owner's choice of $Q$ is still legal because
its cache was frozen during the clearing.
Every other old zero holder was passive and protected.
Now $x\to a$ is forced, since $a^1$ is older than $Q^0$.
This deletes the only new positive copy and makes the
owner cache fresh $x^0$. All old $k$ zero copies are
restored, and all $a$ copies are again zero.
The cached-target owner service adds a zero within
$n+2$ further events. There can have been at most one
earlier owner-head clearing before this branch.
A safe bound for either emptying the owner or gaining
a zero in this phase is therefore $3n+6$ events.

Suppose it ends with the owner empty and no gain.
If a holder caches a nonholder, the four-exchange
gadget of Lemma~\ref{lem:c2infection} gives a gain
with all $a$ copies zero at its endpoint.
Otherwise each holder cache is $a^0$ plus at most
one residual identity from $I$. Weak connectivity,
the empty owner, and the absence of any $a$ descriptor
in the nonholder set
$J=V\setminus(I\cup\{a\})$
supply an incoming boundary descriptor from $J$ to $I$.

Use the incoming-boundary distinction in
Theorem~\ref{thm:thesisc2}, with one modification:
normalize only its selected root $Q$ by the residual
clearing service, rather than making its whole component
a star. This costs at most $n$ events, keeps the outside
boundary initiator frozen, and ends with
$C_Q=\{a^0\}$. The residual weak-component partition
is preserved. If $Q$'s component is nontrivial, its
connectivity and $Q$'s empty residual cache imply
that some other member $v$ still has
$C_v=\{a^0,Q^B\}$ for a finite $B\ge0$.
This is sufficient for Lemma~\ref{lem:c2infection};
that lemma does not require $B=0$.

For completeness, a boundary initiator $P\in J$ whose
oldest partner lies in $J$ first sends the target
$Q\in I$ as extra. That responder is forced to cache
$P^0,Q$, so it can next contact $Q$.
After the root reply $a^0$, it replaces $Q$ as a zero
holder and has a residual edge to nonholder $P$.
The four-exchange gadget therefore gains a zero.
If $P$ instead contacts $Q$ directly, a singleton
$P$ gains immediately. For $C_P=\{Q,X\}$ the exchange
replaces zero holder $Q$ by $P$, which caches
$a^0,X$. If $X\in J$, apply the infection gadget.
If $Q$'s residual component was nontrivial, use the
incoming neighbor $v$ of the now-nonholder $Q$
and apply the same gadget.

The only remaining branch has $Q$ as an isolated
residual component and $X$ in another holder component.
The exchange replaces isolated $Q$ by $P\to X$,
strictly reducing the number of residual components.
It preserves $k$, the empty owner, the all-zero
property of every $a$ copy, and residual closure.
Crucially, this branch costs only one event:
the isolated $Q$ needed no residual clearing.
At most $n-1$ such reductions precede a gain.
The final gain attempt costs at most $n+6$ events,
including root clearing, any preparatory outside
exchange, and the four-event infection.
Thus a safe bound for this incoming phase is $2n+6$.
Together with the owner phase, each all-zero strict-gain
macro costs at most $5n+12$.

The positive-copy and all-zero macros all increase
the same integer $k$. At most $n-1$ such gains produce
a universal zero anchor, with conservative total
budget
\[
 1+(n-1)(5n+12).
\]
The extra initial event is harmless overcounting.

Finally put $D=2n(n-1)$. The capacity-one theorem
normalizes any residual component to a prescribed
zero-age rooted star in at most $D$ events, and the
sum of these budgets across a partition is also at
most $D$. Proposition~\ref{prop:anchorcanonical}
can therefore be implemented with at most
\[
 2(D+1)+D+4(n-2)+D
 \le 4D+4n+2=8n^2-4n+2
\]
events: empty the owner, normalize the residual
components once, merge their roots by the four-event
table, and normalize the connected remainder once.
The root kept during merging stays singleton
$\{a^0\}$ and unprocessed component roots are
untouched, so no intermediate renormalization is
needed. Choose the residual star root as the
second-largest prescribed label. Together with
empty $a$, the result is exactly the desired
$H(V,2)$; no root-change step is needed.

The entire controlled path has length at most
\[
 1+(n-1)(5n+12)+(8n^2-4n+2)
 =13n^2+3n-9<16n^2.
\]
Pad to $T_n$ by activating the canonical empty root.
Every prescribed event has probability at least
$1/(4n)$: the initiator costs $1/n$, an oldest tie
at most a factor two, and a possible passive victim
choice at most a factor two. For $c=\ell=2$ the
request extra and complete reply are forced, and
an active victim, when needed, is unique.
Multiplying these probabilities gives the asserted
common-atom bound. Lemma~\ref{lem:atom} proves the
total-variation estimate, and
$(1-\gamma)^m\le e^{-\gamma m}$ gives the mixing bound.
\end{proof}

\begin{proposition}[Polynomial availability bounds]
\label{prop:c2availabilityquant}
For the stationary law at $c=\ell=2$ and every observer $i$,
\[
 \alpha_1:=\pi\{|C_i|\ge1\}
 \ge\frac{n-1}{2n}\ge\frac13,\qquad
 \alpha_2:=\pi\{|C_i|=2\}\ge n^{-3}.
\]
More generally, for any proved unique equivariant stationary
law with capacity $c$, connectivity gives
$\pi\{|C_i|\ge1\}\ge(n-1)/(nc)$.
\end{proposition}
\begin{proof}
Every connected state has at least $n-1$ directed descriptors.
Label symmetry gives $\mathbb E_\pi|C_i|\ge(n-1)/n$,
and $|C_i|\le c\,\mathbf1_{\{|C_i|\ge1\}}$ proves
the first and the general bounds.

For $c=2$, from every state the event that some cache is
full has probability at least $n^{-2}$ after two events.
If a full cache $q$ already exists, prescribe the same
other initiator twice. A passive cache never decreases
its cardinality, so $q$ remains full.
Otherwise all caches have size at most one.
If a directed path $u\to v\to w$ on three distinct
vertices exists, activating $u$ makes $v$ full; choose
an initiator other than $v$ for the second event.
If no such path exists, the connected functional graph
is an in-star toward an empty root: a longer cycle,
a two-cycle with an attached vertex, or a rooted path
of length two would all give the excluded path.
Activate two distinct leaves toward the root.
The first inserts its address there; the second makes
that root full. These prescribed initiators have total
probability $n^{-2}$, with no additional random choice.
Stationarity implies $\pi\{\text{some cache is full}\}\ge n^{-2}$.
Symmetry then gives
$\alpha_2=\mathbb E_\pi[\#\{\text{full caches}\}]/n\ge n^{-3}$.
\end{proof}

\begin{corollary}[Explicit conditional single-sample mixing]
\label{cor:c2samplemixquant}
For $0<\varepsilon<1$, one returned identifier is within
$\varepsilon$ of uniform, conditional on a nonempty cache
at the deterministic observation time, at every event time
at least
\[
 T_n\left\lceil\Gamma_n^{-1}\log(12/\varepsilon)\right\rceil.
\]
For a request of two identifiers, the corresponding
sufficient horizon is
\[
 T_n\left\lceil\Gamma_n^{-1}\log(4n^3/\varepsilon)\right\rceil.
\]
\end{corollary}
\begin{proof}
Use Theorem~\ref{thm:stationary} with full-state error
at most $\varepsilon/12$ and $\varepsilon/(4n^3)$,
respectively, and the preceding availability bounds.
Both choices also ensure that the conditioning event
has positive probability and the theorem's error
threshold is satisfied.
\end{proof}

\begin{proposition}[A baseline lower bound for one returned identifier]
\label{prop:c2samplelower}
For every $n\ge3$ there is a zero-age initial state and
observer whose conditional sampling error at event time $t$ is at least
\[
 \left((1-1/n)^t-\frac1{n-1}\right)_+.
\]
In particular the worst-case single-sample mixing time
is $\Omega(n)$ global events for a fixed error below one,
as $n$ grows.
\end{proposition}
\begin{proof}
Use the in-star with empty center $a$ and every leaf
storing only $a^0$. Fix a leaf $i$.
Until $i$ initiates, its identity is unknown elsewhere,
so it cannot be contacted and its cache stays $\{a^0\}$.
This event has probability $(1-1/n)^t$.
Conditioning on nonemptiness can only increase this
lower bound on the output probability of $a$.
The uniform target assigns $a$ probability $1/(n-1)$.
\end{proof}

The exponential upper bound is a consequence of the
common-atom proof's prescribed-schedule probability.
The polynomial length of its controlled path is not
a polynomial expected mixing time. A fast bound needs
a more quantitative analysis of ordinary random histories.

\subsection{Ordinary random exchanges and a small exact benchmark}

\begin{lemma}[Young components equal interaction components]
\label{lem:younginteractioncomponents}
Initially mark every descriptor auxiliary old. Mark every fresh
self descriptor young, propagate marks with copied descriptors,
and mark a duplicate young if either input is young, independently
of which input supplies its numerical minimum age.
Let $G_t^Y$ be the undirected graph of young cache descriptors
after $t$ activations, and let $J_t$ contain every undirected
initiator--responder pair that has exchanged during those activations.
Under the adopted protection and sent-victim rules, for every
capacity and shuffle length, $G_t^Y$ and $J_t$ have exactly the
same connected-component partition. In particular this partition
can only coarsen.
\end{lemma}
\begin{proof}
Consider an exchange $p\to q$. Its fresh self descriptor creates
the young edge $pq$. A removed young selected-partner edge is
therefore immediately replaced. A removed young extra edge $pr$
is copied and protected at $q$, so it is replaced by the young
path $p,q,r$. A removed young responder edge $qr$ was sent and
is protected at $p$, so it is replaced by $q,p,r$.
The apparent self-filtering exception $r=p$ cannot be a responder
victim, since incoming fresh $p$ is protected. Young-OR duplicate
merging cannot erase a young connection. Thus old young components
remain connected. Conversely each new young edge is either the
fresh edge $pq$, or comes from an old young edge at one endpoint
and can only join the old components of $p$ and $q$.
Consequently the young partition is updated precisely by joining
those two components, which is also the update of the interaction
partition. Both partitions are discrete initially. Induction proves
the assertion. Numerical ages play no role in this argument.
\end{proof}

\begin{theorem}[Linear expected contact time across every cut]
\label{thm:c2readycontact}
Consider the adopted algorithm with $c=\ell=2$ under its
independent uniform global initiator scheduler. Fix a
nontrivial partition $V=S\sqcup G$, and suppose some
cache in $S$ contains an identity $r\in G$.
Let $\tau$ be the first exchange crossing this cut,
in either direction. Then, uniformly over arbitrary
finite initial ages and all cache contents,
\[
 \mathbb E_x\tau\le7n,\qquad
 \Pr_x\{\tau>14nk\}\le2^{-k}\quad(k\in\mathbb Z_{\ge0}).
\]
No controller, frozen scheduling rule, or prescribed
sequence of initiators is used.
\end{theorem}

\begin{proof}
Until the first crossing, all events are internal to
$S$ or internal to $G$. Let $N$ count the caches in
$S$ containing identity $r$. It cannot decrease before
$\tau$. Indeed, in any $S$-internal exchange the owner
$r$ is not an endpoint. If both endpoints hold $r$,
the request's extra and the complete reply protect it
at both endpoints. If just one holds $r$, it is
forwarded to and protected at the opposite endpoint,
whether or not the original location survives.
Events internal to $G$ do not change the $S$ caches.
Thus $N\ge1$ throughout the pre-crossing trajectory.

Call a copy \emph{ready} when its cache is full and
has the form
\[
 C_q=\{r^A,s^B\},\qquad A\ge B.
\]
Let $R$ count such caches in $S$. Their initiators
select $r$ with conditional probability at least
$1/2$. Consequently, if $R>0$, the conditional
probability $a$ of a crossing at the next global
event satisfies
\[
 a\ge\frac{R}{2n}.
\]

We need two elementary facts about losses of readiness:
a noncrossing event decreases $R$ by at most one,
and its conditional probability $d$ of doing so obeys
\[
 d\le\frac{2R}{n}.
\]
Here is the complete charging argument. Only the
initiator $p$ and responder $q$ can change.
If $p$ is ready and does not cross, its other head
$q$ must be tied with $r$. Its request contains
fresh $p^0$ and $r$, so the responder ends with
a full ready cache. At worst the initiator loses
its ready status; even if the responder was already
ready, the total loss is at most one.
Charge all these events to the clocks of the
$R$ ready initiators, of total probability $R/n$.

It remains to inspect a ready responder
$C_q=\{r^A,s^B\}$ when the initiator was not ready.
If the request includes $r$, the responder ends
with $\{p^0,r^C\}$ and stays ready.
If it does not include $r$ and $p\ne s$, the
complete reply contains two usable identifiers
$r,s$. A singleton initiator receives both and
is ready. A full initiator either discards its
novel extra to receive both, or its extra is
$s$, whose minimum-merged age is at most
$B\le A$. Again readiness transfers to the
initiator whenever it is lost at the responder.
The sole remaining possibility is $p=s$:
the reply entry $s$ is self-filtered, and the
initiator's retained extra may be older than $r$.
Thus a loss not already charged to a ready
initiator requires activation of the ready
responder's unique other cached identity $s$.
Union-bounding these at most $R$ clocks costs
at most another $R/n$. The same cases prove
that the total loss is never larger than one.

Use the nonnegative potential
\[
 \Psi(R)=n\bigl(2+5(4/5)^R\bigr)
\]
before a crossing, and set its value to zero after
one. Put $q_*=4/5$. If $R>0$, readiness gains
can only decrease this potential, a loss of one
costs exactly $nq_*^{R-1}$, and a crossing removes
$n(2+5q_*^R)$. Its conditional drift is therefore
at most
\[
 \begin{split}
 \mathbb E[\Delta\Psi\mid\text{present state}]
 &\le-a n(2+5q_*^R)+d n q_*^{R-1}\\
 &\le-\frac R2(2+5q_*^R)+2R q_*^{R-1}
 =-R\le-1.
 \end{split}
\]

If $R=0$, some $r$ holder still exists because
$N\ge1$. Select any such holder mathematically
for a hazard bound; this does not alter the
scheduler. With probability at least $1/n$ it
is the next initiator. A singleton holder
contacts $r$ and crosses. A full nonready
holder has its other entry strictly older
than $r$ and selects that entry. If this
crosses, we are done. Otherwise its request
forces the responder to cache
$\{p^0,r^C\}$ with $C\ge0$, creating a ready
copy. Thus with probability at least $1/n$
the potential falls from $7n$ to at most
$6n$ (or zero); other events cannot increase
it above its value at $R=0$. Consequently
its conditional drift is again at most $-1$.

For any integer $t\ge0$, summing this drift
until $\tau\wedge t$ and using nonnegativity
gives
\[
 \mathbb E_x(\tau\wedge t)
 \le\Psi(R_0)\le7n.
\]
Monotone convergence proves
$\mathbb E_x\tau\le7n$, with no assumption
of finite expected contact time in advance.

Markov's inequality gives
$\Pr_x\{\tau>14n\}\le1/2$ uniformly.
Before a crossing, $N\ge1$ still holds,
so the same estimate applies at every
$14n$-event block boundary.
The Markov property proves the stated tail.
\end{proof}

\begin{corollary}[Ordinary temporal connectivity in quadratic time]
\label{cor:c2readytemporal}
Start from any weakly connected overlay. Let
$\mathcal I_t$ be the undirected graph of pairs
that have exchanged during the first $t$ events,
and let $\tau_{\rm conn}$ be its connectivity time.
Then
\[
 \mathbb E_x\tau_{\rm conn}\le7n(n-1),
 \qquad
 \Pr_x\{\mathcal I_t\text{ is disconnected}\}
 \le(2^{n-1}-1)\,
       2^{-\lfloor t/(14n)\rfloor}.
\]
Thus
\[
 \Pr_x\!\left\{
 \tau_{\rm conn}>
 14n\left\lceil n-1+\log_2(1/\varepsilon)\right\rceil
 \right\}\le\varepsilon.
\]
These are communication bounds; they do not
assert uniform peer sampling at that time.
\end{corollary}

\begin{proof}
For each fixed unordered nontrivial cut, initial
weak connectivity supplies a descriptor across it.
Orient the cut so its tail is in $S$ and its
head $r$ is in $G$. The theorem bounds its
first observed crossing. There are
$2^{n-1}-1$ unordered nontrivial cuts; a union
bound proves the tail.

For the expectation, whenever the observed
interaction graph is disconnected, select one
of its components and keep its vertex set
fixed until the next crossing of that cut.
Conditional on the present overlay, orient the cut as
needed and apply the theorem to bound the expected wait by $7n$.
This crossing merges observed components,
and at most $n-1$ such phases are needed.
Summing their conditional mean bounds gives
the claimed expectation.
\end{proof}



\begin{corollary}[Every peer emits a fresh self descriptor quickly]
\label{cor:c2successfulsourcecover}
Let $\zeta_r$ be the first event at which peer $r$
initiates a nonempty-cache exchange. Uniformly over
all weakly connected states,
\[
 \mathbb E_x\zeta_r\le8n,\qquad
 \Pr_x\{\zeta_r>16nk\}\le2^{-k}\quad(k\in\mathbb Z_{\ge0}).
\]
Consequently, for $0<\varepsilon<1$,
\[
 \Pr_x\!\left\{\max_r\zeta_r>
 16n\left\lceil\log_2(n/\varepsilon)\right\rceil
 \right\}\le\varepsilon.
\]
This means every peer has sent its fresh self
descriptor; it does not mean all old descriptor
copies have disappeared or samples are uniform.
\end{corollary}
\begin{proof}
If $C_r$ is nonempty, passive exchanges cannot
empty it, because at least one valid fresh self
identifier is received and protected. Until $r$
initiates, it therefore stays nonempty.
The next independent initiator clock of $r$
has mean waiting time $n$ events and produces
the desired exchange.

If $C_r$ is empty, weak connectivity supplies
some other cache containing $r$.
Apply Theorem~\ref{thm:c2readycontact} to the cut
$\{r\},V\setminus\{r\}$. Its first exchange has
mean waiting time at most $7n$ and must arrive
at $r$, since an empty-cache initiator is a
no-op. It makes $C_r$ nonempty.
The preceding argument then costs at most
another $n$ events, giving $\mathbb E\zeta_r\le8n$.
Markov's inequality and restart before $\zeta_r$
give the $16n$ block tail. A union bound over
the $n$ peers proves the source-cover statement.
\end{proof}


These are bounds for ordinary random executions, rather than
probabilities of one long prescribed execution. They establish
temporal connectivity, \emph{not} a sampling mixing bound:
connectivity does not determine the distribution of cache labels.
Even complete descriptor refreshing is insufficient by itself.
For example, begin on three peers at
$C_0=\varnothing$, $C_1=C_2=\{0^0\}$ and mark all initial
descriptors old. After one uniformly random permutation of the
three initiators, every initial descriptor has been replaced by
a young one, but observer $1$ has availability $7/12$ and its
conditional sample probabilities for identities $0,2$ are
$(3/7,4/7)$, rather than $(1/2,1/2)$.
Conditioning the same experiment on all three initiations being
successful nonempty-cache exchanges still gives a biased law:
this event has probability $2/3$, conditional availability is
$5/8$, and the conditional sample probabilities are $(3/5,2/5)$.


\paragraph{An exact three-peer recurrent-class benchmark.}
The checker \texttt{verify\_c2\_recurrent\_mixing.py} enumerates
the entire family reachable from $H(V,2)$ when $n=3$ and
$c=\ell=2$. This family has exactly $75$ states and maximum
age one, with no age truncation. Its transition probabilities
are computed exactly, and integer matrix--vector powers give
the worst conditional single-identifier sampling error over
all initial states \emph{in this recurrent class}. At eight
global activations this error is exactly
\[
 \frac{9483}{1420660}\approx0.00667507,
\]
and at $25$ activations it is exactly
\[
 \frac{18359438816999}{23769834206568612112}
 \approx7.72384\times10^{-7}.
\]
This is an exact finite calculation for recurrent starts on
three peers. It supplies neither an arbitrary-$n$ rate nor a
bound from arbitrary initial descriptor ages.

\subsection{A direct target for faster single-sample mixing}

One can aim to remove the label bias of the returned identifier
without coupling the entire age-labelled overlay. The following
statements make this alternative quantitative. Throughout this
subsection $c=\ell=2$ and $n\ge3$.

\begin{lemma}[Finite-time availability without a common atom]
\label{lem:c2transientavailability}
For every allowed initial state $x$ and observer $i$, put
$\alpha_{x,i}(t)=\Pr_x\{|C_i(t)|\ge1\}$. Then
\[
 \alpha_{x,i}(2k)\ge
 \frac{1-(1-2/n-1/(2n^2))^k}{4n+1},
 \qquad
 \alpha_{x,i}(t)\ge\frac1{20n}\quad(t\ge2n).
\]
\end{lemma}
\begin{proof}
If $C_i$ is empty, weak connectivity supplies an incoming
descriptor $i\in C_j$. Prescribe $j$ as the next initiator.
If it contacts $i$, then $i$ becomes nonempty and remains so
provided the second initiator differs from $i$.
Otherwise its two-entry request forces the responder $q$
to have $C_q=\{j^0,i^b\}$ for some $b\ge0$.
Choosing $q$ as the second initiator makes $i$ oldest or tied,
and contacts it with probability at least $1/2$.
Regardless of the first oldest choice, these events make $i$
nonempty after two activations with probability at least
$b_n=1/(2n^2)$. When $i$ is initially nonempty, only its own
activation can empty its cache. The probability of such an
activation in a two-event block is at most $2/n$.
Consequently $p_k=\alpha_{x,i}(2k)$ satisfies
\[
 p_{k+1}\ge (1-2/n)p_k+b_n(1-p_k)
 =\left(1-\frac2n-\frac1{2n^2}\right)p_k+b_n.
\]
The coefficient is nonnegative. Iteration with $p_0\ge0$
gives the first bound. For $k\ge n$ its right-hand side
is at least $(1-e^{-2})/(4n+1)\ge1/(10n)$.
At an odd observation time, the preceding nonempty cache
survives with probability at least $1-1/n$.
The stated uniform bound follows.
\end{proof}

\begin{proposition}[Reduction to pairwise label bias]
\label{prop:c2localcouplingtarget}
Write $Q_{x,i,t}$ for the law of a single uniformly chosen
cache identity at observer $i$ at time $t$, using output
$\bot$ when its cache is empty. For distinct $a,b\ne i$,
let $\sigma_{ab}$ exchange labels $a,b$ everywhere in a
state while fixing $i$, and define
\[
 \beta_{x,i}(t)=\max_{a\ne b,\ a,b\ne i}
 \|Q_{x,i,t}-Q_{\sigma_{ab}x,i,t}\|_{\TV}.
\]
Whenever $\alpha_{x,i}(t)>0$,
\[
 \left\|\Law_x(S_i^{(1)}(t)\mid |C_i(t)|\ge1)
       -\Unif(V\setminus\{i\})\right\|_{\TV}
 \le\frac{n-2}{2\alpha_{x,i}(t)}\,\beta_{x,i}(t).
\]
For $t\ge2n$ the right-hand side is at most
$10n(n-2)\beta_{x,i}(t)$.
Thus a bound $\beta_{x,i}(t)\le K_n e^{-t/L_n}$,
uniform over initial states $x$ and observers $i$, with
$K_n,L_n>0$, would imply single-sample error at most $\varepsilon$
at every integer time
\[
 t\ge\max\!\left\{2n,\,
 L_n\log\!\frac{10n(n-2)K_n}{\varepsilon}\right\}.
\]
This last implication is a sufficient criterion;
polynomial values of $K_n,L_n$ are not proved here.
\end{proposition}
\begin{proof}
Put $q_a=Q_{x,i,t}\{a\}$ and
$\alpha=\sum_{a\ne i}q_a$. Equivariance gives
$Q_{\sigma_{ab}x,i,t}=\sigma_{ab}Q_{x,i,t}$,
so these two laws differ only by exchanging their masses
at $a,b$. Their total-variation distance is exactly
$|q_a-q_b|$. For each $a$,
\[
 \left|q_a-\frac\alpha{n-1}\right|
 \le\frac1{n-1}\sum_{b\ne i}|q_a-q_b|
 \le\frac{n-2}{n-1}\,\beta_{x,i}(t).
\]
Sum over $a$ and divide by $2\alpha$.
Lemma~\ref{lem:c2transientavailability} gives the
finite-time bound and the displayed sufficient horizon.
In particular, a coupling of executions from $x$ and
$\sigma_{ab}x$ bounds the corresponding query distance
by its probability of unequal outputs. It suffices to
couple their cache identity sets at $i$; equality of all
other caches or numerical ages is unnecessary.
\end{proof}


\begin{lemma}[A coupling through a symmetry set]
\label{lem:c2symmetrycoupling}
Fix a transposition $\sigma=\sigma_{ab}$ with $a,b\ne i$,
and let
\[
 F_\sigma=\{z:\sigma z=z\},\qquad
 \tau_\sigma=\inf\{t\ge0:X_t\in F_\sigma\}.
\]
For every initial state $x$,
\[
 \|\widehat P_2^t(x,\cdot)
       -\widehat P_2^t(\sigma x,\cdot)\|_{\TV}
 \le\Pr_x\{\tau_\sigma>t\}.
\]
Consequently the same bound holds for the two query laws
in Proposition~\ref{prop:c2localcouplingtarget}.
\end{lemma}
\begin{proof}
Run $X$ from $x$. Until $\tau_\sigma$, define the second
execution by $Y_t=\sigma X_t$. Equivariance makes its
transition law the correct kernel from $Y_t$. At the
hitting time, $Y_{\tau_\sigma}=X_{\tau_\sigma}$.
Thereafter use identical random choices in both executions.
The second execution still has the correct marginal law,
and they agree whenever $\tau_\sigma\le t$.
The coupling inequality gives the claim.
\end{proof}

For $n\ge4$, the set $F_\sigma$ includes states in which
$a,b$ have identical external cache identities and ages,
and no other cache contains either identity. All other
cache contents can otherwise be arbitrary. Thus this
coupling can meet without reaching the canonical overlay.
A polynomial uniform bound on the expected hitting time
of $F_\sigma$ would imply a polynomial sampling bound:
Markov's inequality and restart would give a geometric
tail, followed by Proposition~\ref{prop:c2localcouplingtarget}.
Such a hitting-time bound is not established here.

Existing low-dimensional analyses concern a different model.
Bonnet, Tronel, and Voulgaris~\cite{cyclonperformance}
study basic random shuffling rather than the timestamp-based
variant, and report an $nc\log(1/\varepsilon)+O(1)$ mixing
estimate for their single-node in-degree Markov model.
That estimate concerns the modeled degree chain. Applying it
to the adopted oldest-first algorithm would require justifying
the projection and controlling local identity bias; we do not
import that estimate into the present proof.

\subsection{A proved antisymmetric contraction on three peers}

The next result completes the direct contraction route on three
peers, including transient states with arbitrary initial ages.
It uses an exact finite projection and a computer-assisted
integer certificate. The earlier $75$-state recurrent-start
benchmark and this uniform all-initial-state bound have
different scopes.

\begin{proposition}[A finite positivity projection on three peers]
\label{prop:c2threepositivity}
For $n=3$ and $c=\ell=2$, let $\mathcal P$ be the family in
which every cache has at most one positive-age descriptor.
Every peer's cache lies in this family after its first
participation in an exchange, and this property is permanent.
Consequently, from every weakly connected initial state with
arbitrary finite ages,
\[
 \Pr_x\{X_m\notin\mathcal P\}\le3(2/3)^m.
\]
On $\mathcal P$, replacing every positive numerical age by one
is an exact Markov projection, with $487$ weakly connected
projected states.
\end{proposition}
\begin{proof}
Write the three distinct peers as $p,q,r$, and consider
$p\to q$. After mandatory removal of $q$, the initiator can
retain only identity $r$: the reply contains only $p$ and $r$,
and $p$ is self-filtered. Thus its resulting cache has at most
one entry, regardless of the initial ages. The responder
retains protected fresh $p^0$ and can have at most one other
entry. Hence both updated caches have at most one positive
entry. Untouched caches are unchanged, proving permanence.
If each peer has initiated at least once, all three caches
satisfy the property. The independent uniform initiators and
a union bound give the displayed probability estimate.

Within this family a positive descriptor is the uniquely
oldest entry whenever a second entry is present; a singleton
descriptor is also removed on initiation. Thus no previously
positive active descriptor is ever forwarded as an extra.
An extra had age zero and is sent at age one. Fresh entries
have age zero; copies preserve positivity; and minimum merges
have age zero exactly when at least one input has age zero.
Partner ties occur exactly between two zero entries.
Therefore transition probabilities of identities and
zero/positive flags depend only on those flags, proving the
projection assertion. Each peer has eight possible projected
rows: the empty row, four singleton rows, and the three
two-entry rows with at most one positive entry. Filtering
the $8^3$ triples by weak connectivity yields $487$ states.
\end{proof}

\begin{proposition}[A certified ordinary sampling bound on three peers]
\label{prop:c2threeordinarysample}
For $n=3,c=\ell=2$, the conditional single-identifier sampling
error is at most $\varepsilon$, for every weakly connected
initial state with arbitrary finite ages and every observation
time at least
\[
 \left\lceil\frac{\log(9/\varepsilon)}{\log(3/2)}\right\rceil
 +8\left\lceil\log_2(3/\varepsilon)\right\rceil,
 \qquad 0<\varepsilon<1.
\]
For $\varepsilon=10^{-6}$ the displayed horizon is $216$
global activations. This bound concerns ordinary independent
uniform activations, not a prescribed controlled execution. Its finite operator
certificate applies specifically to three peers.
\end{proposition}
\begin{proof}
Relabel the observer as $0$, and let $\sigma$ swap the other
two labels. Let $P$ be the $487$-state positivity-projection
kernel of Proposition~\ref{prop:c2threepositivity}.
The standard-library checker
\texttt{verify\_c2\_three\_peer\_antisymmetric.py}
enumerates all its transitions exactly. Writing $A=6P$, an
integer matrix, it verifies equivariance and the exact
finite inequalities
\[
 \max_z\frac12\sum_w
   \left|P^8(z,w)-P^8(\sigma z,w)\right|
   =\frac{34447}{69984}<\frac12,
 \qquad
 \min_z(P^4\mathbf1_{\{|C_0|>0\}})(z)=\frac59.
\]
These are exhaustive integer calculations, without age
truncation or floating-point inequalities: the maximum
row-difference numerator is $1653456$ over denominator
$2\cdot6^8=3359232$.

For every $\sigma$-antisymmetric function $f$,
equivariance gives
\[
 |P^8f(z)|
 =\left|\frac12\sum_w
       (P^8(z,w)-P^8(\sigma z,w))f(w)\right|
 \le\frac12\|f\|_\infty.
\]
The kernel preserves antisymmetry, so iteration and
supremum-norm contraction imply
$\|P^sf\|_\infty\le2^{-\lfloor s/8\rfloor}\|f\|_\infty$.
Choose $f(z)$ as the conditional-on-state output probability
of identity $1$ minus that of identity $2$, including the
value zero at an empty cache. This function is antisymmetric
and bounded by one. Its expected value is the difference of
the two unconditional query probabilities.
The availability certificate also persists at all $s\ge4$,
since a Markov kernel cannot decrease a function's minimum.

After $m$ initial activations, write $b=3(2/3)^m<1$.
On the event of entering $\mathcal P$, the remaining projected
process obeys these bounds; its complement has probability
at most $b$. Thus at time $m+s$, for $s\ge4$, the difference
between unconditional query probabilities is at most
$b+2^{-\lfloor s/8\rfloor}$, and availability is at least
$(1-b)5/9$. For two target identities the conditional
total-variation error is their unconditional probability
difference divided by twice availability. It is therefore
at most
\[
 \frac9{10}\,
 \frac{b+2^{-\lfloor s/8\rfloor}}{1-b}.
\]
Taking $m=\lceil\log(9/\varepsilon)/\log(3/2)\rceil$ and
$s=8\lceil\log_2(3/\varepsilon)\rceil$ makes both numerator
terms at most $\varepsilon/3$. Since $\varepsilon<1$, the
displayed error is at most $0.9\varepsilon<\varepsilon$.
The same estimates apply at every later observation time.
\end{proof}

The single-positive-entry argument is specific to three peers.
For four distinct peers $p,q,r,s$, take
$C_p=\{q^1,r^0\}$ and $C_q=\{p^0,s^1\}$.
The legal exchange $p\to q$ leaves
$C_p'=\{r^1,s^1\}$: the self reply $p^0$ is filtered,
and the old extra and other reply are both positive.
At larger network sizes, positivity flags alone do not in
general determine which positive entry is oldest.
Extending the contraction therefore requires a richer
projection or a coupling that controls these comparisons.
No arbitrary-$n$ polynomial sampling estimate follows
from the three-peer certificate.

\begin{proposition}[Capped absolute ages and ranks are not sufficient]
\label{prop:c2caprankobstruction}
For $n=4,c=\ell=2$, no fixed absolute-age cap, even augmented
with the complete current ordering and equalities of all
descriptor ages, is an exact state projection sufficient to
predict the future single-identifier query law from every
allowed initial state.
\end{proposition}
\begin{proof}
Fix a cap $B$ and an integer $K>B$. On four distinct peers
$p,q,r,s$, compare the connected initial states
\[
 C_p=\{q^{K+4},r^K\},\qquad
 C_q=\{p^0,s^{K+d}\},\qquad
 C_r=C_s=\varnothing,
 \qquad d\in\{1,2\}.
\]
Their identity sets, capped ages $a\mapsto\min(a,B)$, and
global age ranks and equalities are identical. An activation
of $p$ necessarily contacts $q$ and leaves
\[
 C_p'=\{r^{K+1},s^{K+d}\},\qquad
 C_q'=\{p^0,r^{K+1}\}.
\]
If $p$ initiates again, its oldest choice is a fair tie when
$d=1$, and is uniquely $s$ when $d=2$. Both possible
responders are empty. Thus the resulting query at $p$ is
uniform on $r,s$ for $d=1$, and returns $r$ for $d=2$.
This initiator history has probability $1/16$.
Every other two-event initiator history gives the same query
law in the two cases: a first activation of $q$ deletes the
only descriptor whose age depends on $d$; initial $r,s$
activations do nothing; and after $p\to q$, a next
activation of $q,r$, or $s$ does not change $p$'s query
probabilities. Hence the full two-event query laws differ
in total variation by $(1/16)(1/2)=1/32$.
All queries are available here. Explicitly their probabilities
on $(q,r,s)$ are $(9,16,7)/32$ and $(9,17,6)/32$.
The exact checker
\texttt{verify\_c2\_age\_projection\_obstruction.py}
reproduces these laws by rational enumeration.
\end{proof}

This excludes the particular cap-plus-rank projection on
arbitrary-age initial states. It does not exclude a richer
finite abstraction that retains critical age differences,
and does not prove that four-peer recurrent ages are unbounded.


\subsection{Zero-age normalization for arbitrary capacity at $\ell=2$}

The strengthened first-contact construction applies to every
cache capacity, and supplies a further general tool.
For an identity $a$, write
\[
 N_a(x)=|\{v\ne a:a\in\operatorname{ids}(C_v)\}|,
 \qquad Z_a(x)=|\{v\ne a:a^0\in C_v\}|.
\]

\begin{proposition}[Refresh without reducing the endpoint holder count]
\label{prop:zeroify}
For $\ell=2$ and arbitrary $c\ge2$, any state can reach
a state in which all existing descriptors for a chosen
identity $a$ have age zero, with endpoint holder count
at least its initial $N_a$. Every initially zero copy
is preserved throughout the construction.
If a positive-age copy exists, one block of at most
$T_{n-1}+n+c$ events strictly increases $Z_a$ and has
endpoint $N_a$ at least its starting value. At most
$n-1$ such blocks suffice.
\end{proposition}
\begin{proof}
Freeze owner $a$ and all its zero-age holders, and apply
Theorem~\ref{thm:firstcontact} to a positive copy outside
that frozen set. The controlled path activates only
positive holders and always sends $a$ as extra before
contacting the owner. If it first contacts a zero holder,
that responder's zero is protected and returned to the
initiator, increasing $Z_a$ without reducing $N_a$.

If it contacts owner $a$, the initiating peer $p$ loses
its positive copy, and the owner retains fresh $p^0$.
Only this terminal step can decrease $N_a$, by one.
Now activate only $a$, retaining target $p$ as extra
until it is selected. Whenever the responder holds $a$,
require it to reply with $a$ and with $p$ if known.
The returned self identity leaves at most one usable
reply, so $p$ cannot be evicted and the owner cache
never grows. Each first responder learns and retains
$p$; repeats return $a,p$ and shrink the cache.
Within $n+c$ events some nonholder, at latest $p$, is
therefore contacted. It receives protected fresh $a^0$,
restoring the lost holder and increasing $Z_a$.
Earlier positive responders may also have become zero;
all old zeros were passive and protected.

The first-contact horizon is at most $T_{n-1}$, since
the polynomial $T_s$ increases with $s$. Repeat while
a positive copy remains. Every block increases the
bounded integer $Z_a$, so at most $n-1$ blocks are
possible, and endpoint holder counts never decrease.
\end{proof}

\begin{corollary}[Maximum zero count also bounds all holders]
\label{cor:maxholders}
Fix any closed communicating class at $\ell=2$ and
identity $a$, and let $k=\max_x Z_a(x)$ in that class.
Then $N_a(x)\le k$ for every state of the class.
Every maximizing state has exactly $k\ge1$ holders,
all at age zero.
\end{corollary}
\begin{proof}
Apply Proposition~\ref{prop:zeroify} to any class state.
Its endpoint remains in the closed class and has
$Z_a=N_a\ge N_a(x)$, proving the bound.
At a maximizing state, $k=Z_a\le N_a\le k$ excludes
positive copies. If $k=0$, no peer knows $a$; weak
connectivity forces its owner cache to be nonempty,
and an owner activation creates $a^0$, a contradiction.
\end{proof}

\subsection{Owner service and a larger range of cache capacities}

\begin{lemma}[Owner service and empty-owner extremal states]
\label{lem:generalownerservice}
At $\ell=2$ and arbitrary $c\ge2$, if owner $a$ caches
a nonholder of identity $a$, a finite owner-only path
increases $N_a$ without losing any existing holder.
The same conclusion holds if a nonholder is directed-reachable
from $a$. In every closed class there is a state maximizing
$Z_a=k$ with $C_a=\varnothing$.
\end{lemma}
\begin{proof}
For a cached nonholder target $z$, activate only $a$ and
send $z$ as extra until the first nonholder contact.
At a holder responder require $a$ in the reply, and $z$
when known. The self descriptor $a$ is discarded at its
owner, leaving at most one usable reply. Thus the target
cannot be evicted and the owner cache cannot grow.
First visits establish $z$ at responders; repeats return
$a,z$ and shrink the owner cache. Within $n+c$ events
some nonholder, at latest $z$, is contacted. Every old
holder was passive and retained received fresh $a^0$;
the nonholder now accepts a new protected copy.

For a directed path to a nonholder, choose a shortest
one, whose internal vertices are holders. Retain its
first head $z$ as extra during off-path contacts.
Before a gain, the same repeat-and-shrink argument forces
a contact with a path vertex. Its reply can contain
$a$ and the next path head, shortening the path. The
remaining suffix is untouched. Repetition reaches a
nonholder in finitely many events. This extends the
argument of Lemma~\ref{lem:c2sourcegain} to arbitrary
capacity and arbitrary ages of the held $a$ descriptors.

Now take a state maximizing $Z_a=k$ in a closed class.
Corollary~\ref{cor:maxholders} says that all its holders
are zero holders, and $N_a\le k$ throughout the class.
The preceding path argument implies that the set $R$
reachable from $a$ is contained in $I\cup\{a\}$, where
$I$ is the zero-holder set. All caches in $R$ contain
only identities in $R$. Freeze its complement. For
$|R|\ge3$, Proposition~\ref{prop:anchorcanonical}, with
effective capacity $\min(c,|R|-1)$ and root $a$, leaves
the owner empty and all other $R$ peers with $a^0$.
The effective capacity changes no transition: caches
have no identities outside $R$, so a larger configured
capacity cannot cause an additional overflow.
If $R=\{a,b\}$, one exchange $a\to b$ suffices,
since $C_b=\{a^0\}$; if $R=\{a\}$ nothing is needed.
Outside zero holders are untouched. The endpoint is
therefore an empty-owner maximizing state in the same class.
\end{proof}

\begin{lemma}[Outward holder gain when $n\le2c$]
\label{lem:halfdensegain}
Suppose $\ell=2$, $c\ge3$, $n\le2c$, $C_a=\varnothing$,
and exactly $h$ peers hold identity $a$, all at age zero.
If any holder caches a nonholder, a finite legal path
increases the total holder count beyond $h$.
No connectivity assumption is needed for this lemma.
\end{lemma}
\begin{proof}
We describe a controlled path and stop as soon as the
holder count exceeds $h$. Initially track the specified
holder $p$, and let $x$ be a cached nonholder target.
Activate only $p$ while it contacts other holders.
Choose a nonanchor oldest entry in a tie, send $a$ as
extra, and require $a^0$ and $p$ when known in the reply.
The anchor is refreshed to zero; at most one other usable
identity arrives, so the active cache cannot grow.
Every other resident, including target $x$, is unsent
and cannot be a victim. First visits establish $p$ at
responders; repeats return $a,p$ and shrink the cache.
Consequently a nonholder is contacted within $n+c$ events.
All old holders have been preserved before this contact.

Consider any such contact $p\to q$, with $q$ a nonholder
and outgoing packet $\{p^0,a^A\}$, $A>0$. If $p$ is
underfull, its mandatory deletion creates at least two
vacancies; every reply fits and the old anchor survives,
while $q$ accepts a new one. If $C_q$ contains an identity
in $\operatorname{ids}(C_p)\cup\{p\}$, require one such
identity in the reply. It is duplicate or self-filtered
at $p$ (and cannot be $q$), leaving at most one usable
new identity. Again the old anchor survives and a gain
occurs. These reply choices have legal victims: when
fresh $p$ is known at $q$ only $a$ is new there; otherwise
both incoming identities are new and all returned
residents are eligible. The prescribed shared reply
does not restrict feasibility.

Thus the only remaining case has $p$ full, two usable
novel reply identities, and
\[
 \operatorname{ids}(C_q)\cap
       (\operatorname{ids}(C_p)\cup\{p\})=\varnothing.
\]
In particular $|C_q|\le n-c-1\le c-1$.
The active cache must evict its only extra, $a$, and
remains full; the passive cache accepts protected $a$
and fresh $p^0$, ending with size $\min(c,|C_q|+2)$.
Call this event a migration and track the new holder $q$.
The old holder $p$ is now a nonholder, cached as a fresh
predecessor at $q$. The total holder count is still $h$;
the only possible positive-age $a$ descriptor is at the
tracked holder. All other holders remain at zero.

After a migration, take one activation of the tracked
holder $u$, keeping all other initiators frozen. If its
partner is a zero holder, send $a$ and require $a^0$ in
the reply, with $u$ when known. Its cached nonholder
predecessor is unsent and survives. The anchor becomes
zero, the cache cannot grow, and the preceding
single-holder service applies. If the partner is a
nonholder, use the same gain-or-migration choice above.

The remaining possibility is partner $a$ itself. Let $v$
be the cached nonholder predecessor and send it as extra.
Because the owner has stayed empty, this gives
\[
 C_a=\{u^0,v^{B+1}\},\qquad B\ge0,
\]
and $u$ loses its anchor. The next owner activation must
contact $v$; send $u$ as extra. Peer $v$ acquires fresh
$a^0$, restoring the holder count to $h$. At the owner,
its one remaining extra $u$ and at most two reply entries
fit in capacity $c\ge3$, so $u$ is retained. Peer $u$
is still a nonholder. Lemma~\ref{lem:generalownerservice}
now supplies one more holder. This double callback
therefore gives a strict gain even though the first
exchange temporarily lost one holder.

If the tracked cache is ever underfull, it must lead
to a gain: its next contact either uses this owner
callback, gains at a nonholder without extra eviction,
or refreshes the anchor at a zero holder. In the last
case its cache remains underfull, and single-holder
service reaches a nonholder while retaining the anchor.

It remains to rule out infinitely many full-cache
migrations. Until a gain or a transition to an underfull
tracked cache, count full caches at all nonowners.
Passive cache cardinalities never decrease. An active
full tracked cache either stays full or invokes the
underfull finish just described. At a migration, its
nonholder responder was underfull; if the new tracked
cache is full, this strictly increases the count.
This integer is at most $n-1$, so only finitely many
such migrations can occur. Between migrations,
single-holder service takes finitely many events.
The path must therefore finish with a strict holder gain.
Every prescribed choice is an existing positive-probability
outcome of the adopted protocol.
\end{proof}

\begin{theorem}[Full uniform sampling at $\ell=2$, $n\le2c$]
\label{thm:thesishalfdense}
For $3\le c\le n-1$ and $n\le2c$, the adopted
shuffle-length-two algorithm has a unique stationary
full-state law and converges to it uniformly at a
geometric rate from every weakly connected state with
arbitrary finite integer ages. There exist
$T(n,c)<\infty$ and $\gamma(n,c)>0$ giving the
common-atom and total-variation bounds of
Theorem~\ref{thm:thesisc2}, with atom $H(V,c)$.
For every observer and $1\le b\le c$, stationary
availability is positive and its ordered local sample
converges to uniform conditional on availability.
\end{theorem}
\begin{proof}
Fix a closed class and identity $a$. By
Corollary~\ref{cor:maxholders} and
Lemma~\ref{lem:generalownerservice}, it contains an
empty-owner state maximizing $Z_a=k$, with exactly $k$
zero holders and no positive copies. Put
$I=\{v\ne a:a^0\in C_v\}$ and
$J=V\setminus(I\cup\{a\})$. No class state can have
more than $k$ total holders. Lemma~\ref{lem:halfdensegain}
therefore rules out every edge from $I$ to $J$.
All residual identities in $I$ are consequently in $I$.

If $J$ is nonempty, weak connectivity, the empty owner,
and absence of $J\to a$ and $I\to J$ edges give an
edge from $J$ to $I$. Freeze $I\cup\{a\}$ and use
Theorem~\ref{thm:firstcontact} to reach, but pause just
before, an outside initiator $P$ contacts $Q\in I$.
No outside cache contains $a$ before this first contact.
Pause all outside peers. Reserve the zero anchor at
every member of $I$ and use
Lemma~\ref{lem:anchorreduction} to simulate a
singleton clearing block at $Q$, with residual capacity
$c-1$. All its residual partners lie in $I$.
Stop when the residual cache empties, before it contacts
the owner. This leaves $C_Q=\{a^0\}$, preserves every
other zero holder, and leaves $P$'s cache and actual
oldest choice of $Q$ unchanged.

Resume $P\to Q$. Its at-most-two request identities
and $Q$'s anchor fit in capacity $c\ge3$, so $Q$
keeps its zero. Its sole reply gives $P$ another zero,
contradicting maximality of $k$. Hence $J$ is empty:
every closed class contains a universal zero anchor.

Proposition~\ref{prop:anchorcanonical} then sends every
closed class to the same labelled $H(V,c)$.
The finite-reset, finite-path, self-loop padding and
label-symmetry argument in Theorem~\ref{thm:thesisc2}
gives the stated uniform convergence and conditional
uniformity. Relabelled canonical atoms give each
observer a full cache with positive stationary
probability, proving availability for every $b\le c$.
\end{proof}

\subsection{A structural simplification for nearly full caches}

\begin{lemma}[No initiator-extra eviction at nearly full capacity]
\label{lem:densenoevict}
For any shuffle length and $c\ge n-2$, no initiator
evicts a sent extra descriptor. Its only resident
deletion is the mandatory contacted partner.
\end{lemma}
\begin{proof}
Let the nonempty initiator cache have identity set $P$
of size $d$. A usable new reply identity is outside
$P\cup\{p\}$: the removed partner cannot appear in
its own reply. At most $n-d-1$ such identities exist,
while partner removal creates $c-d+1$ vacancies.
The inequality $c\ge n-2$ gives
$n-d-1\le c-d+1$, so all usable replies fit without
evicting an extra, independently of their number.
\end{proof}

\begin{theorem}[Full uniform sampling at $\ell=2$, $c\ge n-2$]
\label{thm:thesisdense}
For every $n$ and $2\le c\le n-1$ with $c\ge n-2$,
the adopted algorithm at shuffle length two has a
unique stationary full-state law and converges to it
uniformly at a geometric rate from every weakly connected
state with arbitrary finite integer ages.
There exist $T(n,c)<\infty$ and $\gamma(n,c)>0$ giving
the same common-atom and total-variation bounds as in
Theorem~\ref{thm:thesisc2}, with atom $H(V,c)$.
For every observer and $1\le b\le c$, stationary
availability is positive and its ordered local sample
converges to uniform conditional on availability.
\end{theorem}
\begin{proof}
Capacity two is already Theorem~\ref{thm:thesisc2};
assume $c\ge3$. Fix a closed class and a state maximizing
$Z_a=k$ for a prescribed identity $a$.
By Corollary~\ref{cor:maxholders}, its holder set $I$
consists entirely of zero copies, and no class state
can have more than $k$ holders. Put
$J=V\setminus(I\cup\{a\})$.

There can be no edge from $I$ to $J$.
Indeed, if $p\in I$ caches a target in $J$, activate
only $p$, always send $a$ as extra, and choose an
oldest nonanchor partner. This is legal while any
residual entry remains, since the anchor starts at zero.
At any $I$ responder require $a^0$ in the reply and
$p$ when known. The active anchor is refreshed to zero,
the reply has at most one usable new identity, and
every other resident is unsent and retained.
First visits establish $p$ at passive responders;
repeats return $a,p$ and shrink the residual cache.
Within $n+c$ events a $J$ responder must therefore be
selected: the original $J$ target could not disappear
otherwise. At that event Lemma~\ref{lem:densenoevict}
keeps the active copy of $a$, while the nonholder
responder acquires a new protected copy. This gives
$N_a\ge k+1$, contradicting Corollary~\ref{cor:maxholders}.

The owner cannot cache a member $z$ of $J$ either.
To see this, activate only $a$, always retain $z$ as
extra, and at $I$ responders prescribe reply $a,z$
when $z$ is known, otherwise $a$ plus one other entry.
Until a $J$ contact the self reply makes $z$ survive;
first visits establish $z$ and repeats shrink the
owner cache. A $J$ contact must occur within $n+c$
events, creating another holder without losing any
old one. This is the same contradiction.
Thus $G=I\cup\{a\}$ is closed under cached identities.

If $J$ is nonempty, weak connectivity and the absence
of any $a$ descriptor in $J$ give an edge from $J$ to
$I$. Freeze $G$ and apply Theorem~\ref{thm:firstcontact}
until just before an outside initiator $P$ contacts
$Q\in I$. No outside cache has acquired $a$ before
this first contact. Pause $P$ and all other outside
peers. Within $I$, reserve the zero anchor and use
Lemma~\ref{lem:anchorreduction} to simulate a singleton
clearing block at $Q$, with residual capacity $c-1$.
All residual partners are in $I$, since it is closed.
Stop when the residual cache empties, before contacting
owner $a$. This leaves $C_Q=\{a^0\}$, preserves every
other zero holder, and does not change $P$'s cache or
its actual oldest choice of $Q$.

Resume $P\to Q$. The request has at most two identities;
together with $Q$'s old anchor they fit in capacity
$c\ge3$, so $Q$ retains its zero. Its sole reply $a^0$
also gives $P$ a new zero. This contradicts maximality
of $k$. Hence $J$ is empty: every closed class contains
a universal zero-age anchor.

Proposition~\ref{prop:anchorcanonical} sends every such
class to the same labelled $H(V,c)$. The finite-reset
family, finite paths to this shared atom, padding at
its empty root, and label-symmetry argument in
Theorem~\ref{thm:thesisc2} now apply verbatim.
Relabelled canonical atoms give every observer a
full cache with positive stationary probability,
establishing availability for all $b\le c$.
\end{proof}

\subsection{Removing the capacity restriction at shuffle length two}

\begin{lemma}[Outward holder gain for every capacity $c\ge3$]
\label{lem:generaloutwardgain}
Suppose $\ell=2$, $c\ge3$, $C_a=\varnothing$, and
exactly $h$ peers hold identity $a$, all at age zero.
If a holder caches a nonholder, a finite legal path
increases the total holder count beyond $h$.
No restriction on $n/c$ or initial ages is needed.
The path length in this lemma may depend on the ages.
\end{lemma}
\begin{proof}
Use the controlled trajectory in the proof of
Lemma~\ref{lem:halfdensegain}, replacing its full-cache
potential by the argument below. The terminal gadgets
in that proof did not use $n\le2c$: an underfull
current holder with a cached nonholder leads to a gain;
contact with the empty owner gives a gain by double
callback; and a nonholder responder sharing any identity
with the current cache or current self gives a gain
by a suitable reply. We may therefore restrict attention
to a full current holder, with its nonholder predecessor
or initial target still cached, until one of these
terminal cases occurs.

Only the current holder initiates. At every nonowner
contact it sends its fresh self and $a$. At a nonholder
contact not giving a gain, the responder cache is
disjoint from the current cache and current self.
Both reply identities are novel; the active holder
loses its anchor and stays full, and the responder
becomes the new current holder. At a zero-holder contact,
the reply contains $a^0$, so the active holder remains
current and its anchor is refreshed to zero.

Let $Z$ be the initial zero holders other than the
initial current. Throughout an unfinished trajectory,
the holder set is exactly $Z\cup\{\text{current}\}$.
Every member of $Z$ is permanently passive and keeps
its zero. Every past current is a nonholder until it
becomes current again in a migration. The cached
nonholder predecessor is unsent at zero-holder contacts
and therefore survives them. These facts justify the
terminal gadgets at every later state as well.

Initially mark every descriptor old. Mark fresh self
descriptors young and use young-OR on duplicate merges,
independently of which input supplies the minimum
numerical age. Let $F$ count young cache entries,
so $0\le F\le nc$. Identity $a$ stays old because
its owner is frozen. Every identity in $Z$ also stays
old: none of these peers ever initiates, so no fresh
self descriptor with that identity is generated.

We can make $F$ nondecreasing, with a strict increase
at each zero-holder contact. Suppose current $p$
contacts $q\in Z$. If $q$ knows $p$, reply $a,p$;
the active cache loses its mandatory partner and gains
no new identity, becoming underfull and invoking the
terminal gain gadget. Otherwise fresh $p$ is new at $q$.
Require reply $a,D$ if another resident exists.
If $q$ is underfull no victim is needed, and any such
$D$ is allowed. If $q$ is full, choose a nonanchor
$D$ absent from $C_p$ and use it as the responder's
single victim. Such a $D$ exists: its $c-1$ nonanchor
residents can share at most $c-2$ identities with $C_p$,
which also contains $a$ and the selected identity $q$.
Neither of those two identities is a nonanchor resident
of $q$. The reply is legal, protects the active anchor,
and copies a young victim to a previously absent
location at $p$ if compensation is needed. The selected
partner $q$ is old, and no young initiator extra can
be evicted since its only extra is old $a$.
Fresh $p$ supplies an additional young entry at $q$.
Thus $F$ strictly increases.

At a nonholder migration, choose every required
responder victim inside the two-entry reply. The
disjointness means that every young victim is copied
to a new location at the initiator. Its only extra
victim is old $a$, and fresh $p$ adds a young entry
at the responder. Only the selected partner can
remove a young initiator resident. Hence
\[
 F'-F\ge1-\mathbf1_{\{\text{selected partner is young}\}}.
\]
So $F$ cannot decrease, and it increases at every
migration selecting an old partner. No assumption
that the responder was underfull has been used.

Suppose the controlled path never reached a gain.
Every finite unfinished state admits a next event
of the prescribed kind. Since $F$ is a bounded
nondecreasing integer, it would eventually be constant.
There are then no zero-holder contacts, and all events
are nonholder migrations selecting young partners.
Start after one such migration, with current anchor
age $A\ge1$, and count subsequent events by $t$.
At time $t$ its age is exactly $A+t$: a migration
increments and transfers it to a previous nonholder,
with no duplicate refresh.

For every identity $s$, let $M_s(t)$ be its largest
stored age anywhere, or zero if absent, and put
$K_s(t)=M_s(t)-t$. Each event increments only active
residents by one; fresh selves have age zero, and
copies and minimum merges cannot exceed their input
ages. Thus every $K_s$ is nonincreasing. More strongly,
when $s\ne a$ is selected as partner, its sole active
copy is removed, it is not sent as extra, and its own
responder cannot return itself. All remaining $s$
copies have unchanged or reduced ages. Therefore
\[
 K_s(t+1)\le K_s(t)-1
 \quad\text{whenever partner }s\text{ is selected}.
\]
The oldest rule simultaneously requires the selected
descriptor's age to be at least $A+t$, giving
$K_s(t)\ge A$. Each identity can consequently be
selected at most $(M_s(0)-A+1)_+$ times during this
constant-$F$ period. Summing over the finitely many
$s\ne a$ bounds its length, contradicting an infinite
unfinished path. Hence some terminal gain occurs.
Every prescribed exchange has positive probability,
so its finite controlled path is accessible.
\end{proof}

\begin{theorem}[Full uniform sampling for every two-entry shuffle]
\label{thm:thesisgeneraltwo}
For every $n\ge3$ and $2\le c\le n-1$, the adopted
algorithm with fixed shuffle length two converges
from every weakly connected initial state with
arbitrary finite integer ages to a unique stationary
full-state law $\pi$. There exist finite $T(n,c)$
and positive $\gamma(n,c)$ such that
\[
 \inf_x\widehat P_2^T(x,\{H(V,c)\})\ge\gamma,
 \qquad
 \sup_x\|\widehat P_2^t(x,\cdot)-\pi\|_{\TV}
       \le(1-\gamma)^{\lfloor t/T\rfloor}.
\]
For every observer and $1\le b\le c$, stationary
availability is positive and its ordered local
$b$-sample converges to uniform conditional on availability.
No practical mixing-time bound is asserted.
\end{theorem}
\begin{proof}
Capacity two is Theorem~\ref{thm:thesisc2}.
For $c\ge3$, repeat the closed-class proof of
Theorem~\ref{thm:thesishalfdense}, using
Lemma~\ref{lem:generaloutwardgain} in place of
Lemma~\ref{lem:halfdensegain}. Zeroification still
bounds all holder counts by the class maximum $k$
of $Z_a$; an empty-owner maximizer exists by
Lemma~\ref{lem:generalownerservice}. The new lemma
rules out all outward holder-to-nonholder edges.
If any nonholders remain, a paused incoming first
contact and singleton residual normalization give
one additional zero holder, since $c\ge3$.
Thus every closed class contains a universal zero
anchor and reaches the same labelled $H(V,c)$.

Although the new outward path can have an
age-dependent length, the finite-family argument
still gives a uniform minorization: the finite reset
family $F_B$ has only finitely many chosen paths to
this common atom. Their lengths have a finite maximum
and probabilities a positive minimum. Pad at the
empty-root self-loop and apply the uniform reset,
as in Theorem~\ref{thm:thesisc2}. The same argument
gives uniqueness and geometric convergence.
Relabelling symmetry yields conditional uniformity;
relabelled canonical atoms give positive availability
for every $b\le c$.
\end{proof}

\begin{corollary}[Implicit mixing bounds and their quantitative limitation]
\label{cor:twoimplicitmixing}
Let $T,\gamma$ be the constants of
Theorem~\ref{thm:thesisgeneraltwo}, reduced if necessary so that
$0<\gamma\le1/2$. Time here counts global initiator events,
including empty-cache activations. For $0<\varepsilon<1$,
the full-state mixing time satisfies
\[
 t_{\rm mix}(\varepsilon)
 \le T\left\lceil\gamma^{-1}\log(1/\varepsilon)\right\rceil.
\]
For every observer $i$ and request size $1\le b\le c$,
the conditional local sample is within $\varepsilon$ of uniform
at every deterministic event time at least
\[
 T\left\lceil\gamma^{-1}
            \log\!\left(\frac4{\gamma\varepsilon}\right)\right\rceil.
\]
These are bounds in terms of implicit constants, rather than
an explicit rate of growth in $n,c$.
\end{corollary}
\begin{proof}
Use $(1-\gamma)^m\le e^{-\gamma m}$ in the full-state
bound. Equivariance gives the same minorization $\gamma$
to a relabelled canonical atom at which observer $i$ has
a full cache. Stationarity therefore gives availability
$\alpha=\pi\{|C_i|\ge b\}\ge\gamma$.
At the displayed local horizon, full-state distance is
at most $\gamma\varepsilon/4$. It is consequently
at most $\alpha/2$, and
Theorem~\ref{thm:stationary} bounds the conditional
distance by $4(\gamma\varepsilon/4)/\alpha\le\varepsilon$.
\end{proof}

The constants are computable in principle. Let $H,B,F_B$
be the finite-age reset data of
Theorem~\ref{thm:thesisgeneralreset}, and enumerate the finite
family $F_B$. From each member run breadth-first search over
positive-probability transitions until reaching the fixed atom.
The class and reset arguments guarantee accessibility, so each
search terminates. The kernel is finitely branching, with ages
at search depth $d$ at most $B+d$, making every search level
effectively finite. Let $L$ be the maximum of the finitely
many discovered path lengths, and pad shorter paths at the
empty-root self-loop. The uniform one-outcome probability
$q=(nc\,16^c)^{-1}$ gives the valid choices
\[
 T=H+L,\qquad \gamma=q^{H+L}.
\]
This supplies a terminating procedure for a numerical
upper bound at each fixed $n,c$, but no useful closed-form
bound on $L(n,c)$. Moreover, the factor $\gamma^{-1}$
is exponential in the prescribed path length. A useful
mixing estimate needs quantitative control of the ordinary
random evolution, or a much more efficient convergence
argument, beyond this rare-path bound.
All local sampling statements here condition on availability
at a fixed observation time.

\subsection{Three-entry shuffles: first contact and a conditional common atom}

The general lower-length convergence theorem also gives
progress at shuffle length three. The next two results
hold for every $n\ge4$ and $3\le c\le n-1$ at $\ell=3$.
They do not assert uniform sampling at length three.

\begin{theorem}[Uniform first contact for three-entry shuffles]
\label{thm:threefirstcontact}
At $\ell=3$, fix a nonempty proper
$G\subset V$, put $S=V\setminus G$, $s=|S|$, and
suppose some outside cache contains an identity $r\in G$.
There are constants $T_3(n,c,s)<\infty$ and
$\eta_3(n,c,s)>0$, independent of all initial ages,
such that an $S$-initiator path contacts $G$ within $T_3$
events. Under the original scheduler the first exchange
across the cut has the geometric bound
\[
 \Pr_x\{\tau>jT_3\}\le(1-\eta_3)^j.
\]
The constructed path always activates an $r$ holder,
sends $r$ as extra when the partner is not its owner,
and requires $r$ in the reply whenever known. Its total
$r$-holder count cannot decrease except at a terminal
contact with owner $r$.
\end{theorem}
\begin{proof}
Keep $G$ inactive until contact. Let $m$ be the minimum
outside age of $r$ and $A\subset S$ its minimum-age holders.
For an exchange internal to $A$, reserve $r$ as one
request extra and one reply entry. Removing this anchor
gives the exact length-two kernel with capacity $c-1$:
the minimum rule restores the anchor to $m$, its incoming
copy protects it, and all other updates are unchanged.
There is only an oldest-partner qualification. If the
residual oldest age is at least $m$, a reduced oldest
partner is an actual oldest partner. If all residual
ages are smaller than $m$, or the residual cache is
empty, the actual next partner is owner $r\in G$,
which is immediate success. No ages are shifted or truncated.

Suppose a residual $A$ cache has an identity outside $A$
other than $r$. On $V\setminus\{r\}$ use
Theorem~\ref{thm:firstcontact} with active set $A$;
put arbitrary empty auxiliary caches at its complement.
Until the terminal contact, only $A$ caches are read,
so those auxiliary choices are immaterial. Simulate
this lower-length path until a contact with $G$, or
until an initiator in $A$ contacts $q\in S\setminus A$.
In the latter event choose the actual reply to contain
$r$ whenever known. Its age is greater than $m$, while
the sent anchor has age $m+1$. At least one minimum
copy disappears, and no new minimum copy is created.
The auxiliary terminal reply is never needed.

Otherwise all residual $A$ identities lie in $A$.
Choose one weak component $K$ of the residual graph.
For $|K|=1$ its cache is just the anchor and contacts
$G$ on activation. For $k=|K|\ge2$ use effective
capacity $d_k=\min(c-1,k-1)$ and length
$e_k=\min(2,k-1)$. Its lower-length kernel has a
uniform common-atom horizon $U_k$. Indeed, $k=2$
uses length one, and all other capacities are covered
by Theorem~\ref{thm:thesisgeneraltwo}.
Simulate a path to a reduced canonical
state with an empty root cache. A reduced empty-cache
activation instead contacts $G$, so any such interruption
is already success. If no interruption occurs, activate
the empty reduced root once to contact $G$.

Let $A_b$ be the length-two first-contact horizon on
$n-1$ peers at capacity $c-1$ with active-set size $b$,
and choose
\[
 E=1+\max\bigl(\{A_b:1\le b\le n-2\}
        \cup\{U_k:2\le k\le n-1\}\cup\{0\}\bigr).
\]
Within $E$ events one minimum copy disappears or $G$
is contacted. All preterminal initiators and responders
in this phase store $r^m$, so their other minimum copies
are retained unchanged. At most $s$ such phases raise
the minimum outside age by at least one, or contact $G$.

Put $R_s=(sc+1)c$. Raise the outside minimum $r$ age
$R_s+1$ times using the preceding phases, unless $G$
is contacted earlier. Every outside $r$ age then
exceeds $R_s$. Mark all current descriptors old and
future fresh selves young, with young-OR duplicate
colors. For the next $R_s$ events every young age is
at most the elapsed number of events. Thus an oldest
partner of an $r$ holder is old throughout this horizon.

At each step choose an $r$ holder $p$ maximizing its
young count $f_p$ among outside $r$ holders. Use the
reply/victim choices of Lemma~\ref{lem:generalcolor},
but reserve old $r$ among its extras and require all
$Q\cap U$ in the reply, in addition to the responder
victims $D_q$. Reserving $r$ leaves room for the
required young difference extra at length three:
if the cache has only two entries, both its anchor
and selected partner are old and there is no young
difference to send. Requiring all $Q\cap U$ is feasible
by the same bound $|D_q|+|Q\cap U|\le\min(3,|Q|)$.
In particular an existing responder $r$ is returned.
The young-count estimate in that lemma still gives
$F'\ge F$, where $F$ counts young entries in $S$.

Write $M=\max\{f_v:v\in S,\ r\in C_v\}$.
If $F'=F$, the same calculation gives
$Y_p\subseteq Y_q$, $p\in Y_q$, and $f_q\ge f_p+1$.
Thus $q$ could not already have held $r$, by maximality
of $f_p=M$. Its new $r$ copy is protected. On this
plateau no young initiator copy is lost; only $t$
young responder victims transfer to new locations
at $p$, with
\[
 f_p'=f_p+t,\qquad f_q'=f_q-t,\qquad
 0\le t\le f_q-f_p-1.
\]
Consequently the new $r$ holder $q$ has
$f_q'\ge M+1$, even if $p$ loses its anchor.
The pair $(sc-F,c-1-M)$ strictly decreases
lexicographically at every nonterminal step.
There are only $(sc+1)c=R_s$ possible pairs,
since every $r$ holder contains at least one old
descriptor. Therefore $G$ is contacted within $R_s$
such steps.

The minimum-raising and final routing phases activate
only $r$ holders and send $r$ whenever the partner is
not its owner. Existing responder copies are returned
and protected; a newly created responder copy compensates
any initiator loss. This proves the holder-count claim.
One may take
\[
 T_3=sE(R_s+1)+R_s,\qquad
 \eta_3=(nc\,16^c)^{-T_3}.
\]
The probability estimate is the all-length controlled
choice bound. Before the first crossing, $G$-internal
events cannot change the outside caches, and outside
internal events preserve some $r$ there. Restarting
the one-block bound proves the geometric tail.
\end{proof}

\begin{proposition}[A universal zero anchor suffices at length three]
\label{prop:threezeroanchor}
In the same parameter range, a state in which every
nonowner stores $a^0$ has a finite path to any prescribed
canonical state $H(V,c)$. Hence, if every closed class
can reach a universal zero anchor, full convergence and
conditional uniform sampling follow at length three.
\end{proposition}
\begin{proof}
Freeze $a$ and reserve its zero anchor in every request
and reply. Lemma~\ref{lem:anchorreduction} simulates
the length-two, capacity-$c-1$ kernel on residual
components of $V\setminus\{a\}$; reduced empty-cache
activations are omitted. The lower-length common-atom
theorems apply to every component by exactly the
parameter check in Theorem~\ref{thm:threefirstcontact}.

To empty the owner, inspect an actual oldest $q\in C_a$
and keep $a$ frozen while its residual component is
normalized with root $q$. This gives $C_q=\{a^0\}$,
without changing $a$'s oldest choices. Activate $a\to q$.
Its sole reply is the self descriptor $a$, so $C_a$
loses $q$ and no other entry. At $q$, incoming fresh
$a^0$ protects its anchor, and at most two extras fit
beside it in capacity $c\ge3$. All anchors remain zero.
At most $c$ repetitions empty the owner.

Normalize each residual component. If there is more
than one, choose two roots $u,v$, each with the singleton
cache $\{a^0\}$. The same four exchanges used in
Proposition~\ref{prop:anchorcanonical},
\[
 u\to a,\quad v\to a,\quad a\to v,\quad a\to u,
\]
end with $C_a=\varnothing$, $C_u=\{a^0\}$ and
$C_v=\{a^0,u^0\}$. All involved pre-exchange caches
have size at most two, so their packets and replies
are identical at lengths two and three. Other anchors
are unchanged. This merges two residual components.
Repeat normalization and merging, then normalize the
connected residual graph once more. This gives the
canonical state with root $a$. The root-change gadget
of Proposition~\ref{prop:anchorcanonical} also uses
only caches of size at most two, so any labelled
canonical order is reachable.

If universal zero anchors are accessible from every
closed class, all classes therefore share this atom.
The finite-reset and symmetry argument of
Theorem~\ref{thm:thesisc2} gives the final implication.
\end{proof}

\begin{lemma}[A designated owner target can be retained]
\label{lem:ownertargetretain}
For any shuffle length $\ell\ge2$, suppose owner $a$ caches a
designated target $p$ and contacts some $q\ne p$.
Choose $p$ as extra and require $a$ in the reply,
assuming the responder has it. There is a legal
initiator-victim choice retaining $p$. This is a
positive-probability controlled choice, rather than
a claim about every uniform victim outcome.
\end{lemma}
\begin{proof}
If $p$ is returned, received-identity protection suffices.
Otherwise let $d=|C_a|$ before partner removal and
$e=\min(\ell-1,d-1)$ be its number of extras.
The returned self descriptor leaves at most $\ell-1$
usable identities. If $h$ of those are resident
duplicates, at most $\ell-1-h$ are new. The number
of required victims is at most
\[
 (\ell+d-c-2-h)_+.
\]
Among extras other than $p$, at least $e-1-h$ are
not protected by incoming identities. Since
\[
 \ell+d-c-2\le e-1
\]
both when $d\ge\ell$ (use $d\le c$) and when
$d<\ell$ (use $\ell\le c$), these other extras
suffice. If the displayed lower bound is negative,
the victim requirement is zero. Thus $p$ need not
be selected as victim.
\end{proof}

\begin{lemma}[Owner service when zero-holder caches have spare reply choices]
\label{lem:largezerocallback}
Suppose $2\le\ell<c$ and every zero holder
$q\in I=\{q\ne a:a^0\in C_q\}$ has
$|C_q|\ge\ell+1$. If owner $a$ caches a peer outside
$I$, an owner-only controlled path strictly increases
$Z_a$, preserves every old zero, and does not decrease
$N_a$. The same conclusion holds if some peer outside
$I\cup\{a\}$ is directed-reachable from $a$.
No common ages or shared additional descriptors are
required. The path length may depend on initial ages.
\end{lemma}
\begin{proof}
First take a cached target $p\notin I$. Before a
nonzero contact, every responder is a zero holder
with at least $\ell+1$ entries. Require its reply to
contain $a$ and exclude $p$, and retain $p$ as extra
using Lemma~\ref{lem:ownertargetretain}. Such a reply
exists even when $p$ is resident. Passive cardinalities
never decrease, and fresh $a^0$ protects every old
zero, so the size condition continues to hold.
The target is never returned and its owner-cache age
is exactly $A+t$ after $t$ controlled events.

For every identity $s$, define $M_s(t)$ as its global
maximum stored age, or zero if absent. The argument
of Lemma~\ref{lem:generaloutwardgain} applies without
colors: $K_s(t)=M_s(t)-t$ is nonincreasing and drops
by at least one whenever $s$ is selected as partner.
Indeed a selected identity is removed from the active
cache, excluded from its outgoing extras, and absent
from its own responder's reply. Any selected $s\ne p$
has age at least $A+t$, so $K_s(t)\ge A$. Each such
identity can be selected only finitely many times.
Thus an outside-$I$ contact must occur, at latest at
$p$, producing a new zero without losing any old copy.

For the path statement, choose a shortest directed
path to outside $I\cup\{a\}$, with first head $z$.
Its internal peers are zero holders. Until a path
vertex or nonzero peer is contacted, retain $z$
and exclude it from replies. The same age argument
forces such a contact. At any contacted internal
path vertex require its next path head in the reply,
along with self-$a$. This is compatible with excluding
$z$ when needed, and shortcuts the path. The untouched
suffix is preserved because no suffix vertex has yet
responded. Repeating the finite shortening argument
gives a nonzero contact. All responders before that
contact stay zero and large throughout.
\end{proof}

\begin{proposition}[An owner-only callback obstruction at every length]
\label{prop:owneronlyobstruction}
For every $m\ge3$ and $c\in\{m,m+1\}$, there is a strongly
connected state with $n=m+2$ and $\ell=m$,
in which owner $a$ caches a nonholder target $p$ but
can never contact it while $a$ is the sole initiator.
Its caches are initially full at $c=m$ and all underfull
at the proper-shuffle capacity $c=m+1$.
This refutes an auxiliary owner-only service claim;
it is not a counterexample to the original iid
all-peer scheduler or to full CYCLON convergence.
\end{proposition}
\begin{proof}
Take distinct $a,p,w$ and a set $R$ of $m-1$ other
peers. Use the caches
\[
 \begin{split}
 C_a&=\{p^0\}\cup\{r^1:r\in R\},\\
 C_q&=\{a^0,p^0\}\cup\{r^1:r\in R\setminus\{q\}\}
                    \quad(q\in R),\\
 C_p&=\{w^0\}\cup\{r^0:r\in R\},\qquad
 C_w=\{p^0\}\cup\{r^0:r\in R\}.
 \end{split}
\]
Every cache has size $m$, valid at either capacity,
and the overlay is strongly
connected: $a$ reaches $p$ and $R$, $p$ and $w$ reach
each other and $R$, and $R$ reaches $a$ and $p$.
Target $p$ does not store identity $a$.

The first owner partner is some $q\in R$. Its full
reply and minimum-age duplicate merges leave
\[
 C_a=\{p^0\}\cup\{r^1:r\in R\setminus\{q\}\}.
\]
All passive caches remain unchanged: every incoming
identity was already present, with an equal or younger
resident age. From any state of this form, the next
owner partner is some $s\in R\setminus\{q\}$,
and its full reply produces the same form with omitted
identity $s$. The target stays at age zero, while
every available $R$ partner stays at age one. Thus
the owner always selects $R$ and never $p$.
Requests use all available extras, and replies use
all $m$ residents, exactly as the specified algorithm
requires. No overflow is needed in any of these merges,
so raising capacity from $m$ to $m+1$ changes no event.
At $m=3$ the restricted dynamics is a
forced two-state cycle after the first event.
\end{proof}

\begin{lemma}[A cooperative callback with shared zero anchors]
\label{lem:threecooperativecallback}
For any $3\le\ell\le c$, let
$I=\{q\ne a:a^0\in C_q\}$.
Suppose a set $P\subset V\setminus\{a\}$ of
$\ell-2$ identities is cached by $a$, and every
$q\in I$ satisfies
\[
 \{a^0\}\cup\{p^0:p\in P\}\subseteq C_q.
\]
Then a controlled path of at most $c(n+c+1)+1$
events strictly increases $Z_a$, preserves every
original zero copy throughout, and does not decrease
$N_a$. Ages of other descriptors, the owner's other
entries, caches outside $I$, and all other resident
identities are unrestricted. The shared additional zero
anchors are a condition of this lemma, not a new model
assumption.
\end{lemma}
\begin{proof}
The no-self rule implies $P\cap I=\varnothing$.
If an actual oldest owner partner is outside $I$,
activate the owner and finish: fresh $a^0$ creates a
new zero there. Otherwise inspect an actual oldest
$Q\in I$ and freeze the owner.

Reserve $\{a\}\cup P$ as the $\ell-1$ request extras and reply
entries in $I$-internal exchanges. Removing them
simulates exactly the singleton kernel with residual
capacity $c-\ell+1$. All reserved anchors stay at zero
by minimum merging and received-identity protection.
Any residual oldest partner is also an actual oldest
partner: its nonnegative age is at least those of
the zero anchors, so choose it in a tie.
Use the singleton clearing construction with only
$Q$ initiating, but stop before any selected residual
partner $r$ outside $I$. At internal responders return
$Q$ in the residual reply when known. First visits
establish $Q$, repeats return it and shrink its
residual cache, and its size never grows. Within
$|I|+c$ events either the residual cache empties or
such an outside residual partner is selected.
All $I$ anchors are preserved and the owner's
inspected oldest choice is unchanged. In the empty
case $C_Q=\{a^0\}\cup\{p^0:p\in P\}$;
in the other case $Q$
still contains the selected $r\notin I\cup\{a\}\cup P$.

Now activate $a\to Q$, choosing every member of $P$
as extra and any required remaining extra. In the empty
case the reply is precisely the $\ell-1$ reserved
anchors; otherwise require those anchors and $r$.
The owner filters its self descriptor and refreshes
its existing targets in $P$.
It loses $Q$ and receives at most one novel identity,
namely $r\notin I$, which fits the mandatory vacancy.
Thus its cache size cannot grow and its number of
cached $I$ identities strictly decreases. At $Q$,
fresh $a^0$ and the received $P$ descriptors protect
all resident zero anchors;
any extra inserted identity uses a vacancy or a legal
residual victim. Received-identity protection and
minimum merging preserve all reserved anchors regardless
of its other residual identities.

Repeat whenever the owner's actual oldest partner
belongs to $I$. If no $I$ head remains, its cache contains
only identities outside $I$ and still contains the
nonempty set $P$. Its next activation gains a new zero.
Thus at most $c$ clearing-and-owner blocks and one
final activation suffice, giving the stated bound.
Only $I$ and the owner initiate before this final
contact; every original zero stays zero, and all
outside holders are untouched until the final
responder refresh. No copy of identity $a$ is lost.
\end{proof}

\begin{lemma}[A cooperative callback with common ages at full shuffle]
\label{lem:fullpositivecallback}
Suppose $\ell=c\ge3$, let
$I=\{q\ne a:a^0\in C_q\}$, and take a set
$P\subset V\setminus\{a\}$ of size $c-2$.
For each $p\in P$ choose an integer $m_p\ge0$.
Assume every $q\in I$ contains
$\{a^0\}\cup\{p^{m_p}:p\in P\}$, and the owner
caches every $p\in P$ with age at least $m_p$.
No restriction is imposed on residual identities.
A controlled path of at most $2(n+c+1)+1$ events
strictly increases $Z_a$, preserves all old zero
copies, and does not decrease $N_a$.
\end{lemma}
\begin{proof}
Again $P\cap I=\varnothing$. Put $M=\max_{p\in P}m_p$.
Select an outside-$I$ owner partner whenever tied for
oldest; any such contact already finishes with a new
zero. Otherwise let $Q\in I$ be the actual oldest
owner partner and freeze the owner.

Reserve $a$ and $P$ during $I$-internal exchanges.
Their fixed ages $0,(m_p)_{p\in P}$ are preserved
by minimum merging; removing them leaves a singleton
kernel of capacity one. Run its $Q$-only clearing
controller until its residual cache is empty, its
residual partner is outside $I$, or its sole residual
descriptor has age below $M$. Until one of these
conditions holds, a reduced oldest internal partner
has age at least $M$ and is an actual oldest partner.
This partial block has at most $n+c$ exchanges.

Activate $a\to Q$. Its pre-exchange cache consists
of the reserved anchors and at most one residual
descriptor. Its full reply refreshes each owner
$P$ age to exactly $m_p$ and supplies at most that
one residual identity. At the responder the sent
owner $P$ ages are at least $m_p+1$, so its common
age vector is retained. All zero copies of $a$ are
protected. Any newly inserted owner head in $I$
has age strictly below $M$: an outside residual
was outside $I$, and the only other nonempty stopping
case had residual age below $M$.

There are at most two owner cache entries outside
$P$. After one such block, at most one original
$I$ head remains, together with at most one new
$I$ head of age below $M$. If no outside identity
can already be selected as oldest, the remaining
original head is the next $I$ partner. After the
second block the first new head has aged at most
once, so its integer age is at most $M$; any second
new head still has age below $M$. A target in $P$
has age $M$. Thus an outside-$I$ partner is oldest
or tied, and the next owner activation gains a new
zero. This uses at most two partial clearing blocks,
two owner exchanges, and one final contact, within
the stated bound. Only the owner and $I$ initiated
before that contact, and their updates preserved
every original zero and every identity-$a$ copy.
\end{proof}

For a concrete repair of the smallest owner-only
cycle, use peers $0,1,2,3$, $c=\ell=3$, and
\[
 C_0=\{1^0,2^1\},\quad C_1=\{2^0\},\quad
 C_2=\{0^0,1^0,3^1\},\quad
 C_3=\{0^0,1^0,2^1\}.
\]
Owner-only activations alternate $0\to2,0\to3$
forever, but the cooperative sequence
\[
 2\to3,\qquad0\to2,\qquad0\to1
\]
first leaves $C_2=\{0^0,1^0\}$, then leaves
$C_0=\{1^0\}$, then creates fresh $0^0$ at peer $1$.
The original zero copies at $2,3$ are preserved,
and the endpoint has a universal zero anchor.

Theorem~\ref{thm:threefirstcontact} does not yet supply
the missing universal zero anchor. Contact with a
zero holder refreshes a positive active copy while
protecting old zeros. Contact with the owner can
instead remove one holder. An unrestricted length-three
callback that restores this holder count, gains a
zero copy, and restores all old zeros has not been
proved at arbitrary capacities. The dense result below
transfers Proposition~\ref{prop:zeroify} to length
three when $c\ge n-2$.
Proposition~\ref{prop:owneronlyobstruction} shows
that retaining the target alone cannot justify an
owner-only callback, even with full initial caches
and strong connectivity. Lemma~\ref{lem:threecooperativecallback}
supplies a cooperative callback at every shuffle
length under shared zero anchors, and
Lemma~\ref{lem:fullpositivecallback} allows a common
positive age vector at full shuffle. Neither lemma
proves that its shared-descriptor configuration is
always accessible. Mixed ages and missing shared
descriptors remain outside these conditional proofs.
Lemma~\ref{lem:largezerocallback} handles arbitrary
mixed target ages for proper shuffles when every
old zero holder has at least $\ell+1$ cache entries.
This occupancy condition is also conditional: it
does not cover arbitrary underfull states.

\subsection{An exact sufficient route to the remaining general theorem}

The conditional callback lemmas should not be mistaken
for the entire missing proof. The following statement
specifies a sufficient remaining obligation precisely.

\begin{proposition}[Induction from zero-count improvement]
\label{prop:zeroexpansioncriterion}
For every allowed $n,c$ and each fixed shuffle length
$\ell\ge3$, consider the following proof obligation:
\begin{quote}
For every connected state $x$ and fixed identity $a$
with $Z_a(x)<n-1$, there is a finite positive-probability
path to $y$ such that $Z_a(y)>Z_a(x)$.
\end{quote}
If this obligation is established at all these lengths,
the full convergence and conditional uniform-sampling
target follows for every fixed shuffle length.
The obligation remains unproved over the full
parameter range $\ell\ge3$.
\end{proposition}
\begin{proof}
We first record the inductive common-atom step.
Suppose common-atom accessibility is already known
for all smaller shuffle lengths and all their
allowed capacities and populations. At length
$\ell$, a universal zero anchor $a$ lets us reserve
$a^0$ in requests and replies and simulate the
$(c-1,\ell-1)$ kernel on residual components.
On a component of size $k$, the effective parameters
are $\min(c-1,k-1)$ and $\min(\ell-1,k-1)$.
These are covered by the inductive hypothesis;
single-vertex residual components need no exchanges.

The construction of Proposition~\ref{prop:threezeroanchor}
therefore applies unchanged. Normalize the residual
component of an inspected owner-oldest $Q$ with
root $Q$, then contact $Q$ to delete one owner head
using its sole self-$a$ reply. Fresh $a$ and at most
$\ell-1$ extras fit in capacity $c\ge\ell$, so the
passive zero anchor is preserved. Recompute residual
components and repeat until the owner is empty.
Normalize the residual components and merge their
roots with the same four exchanges; all their
involved caches have size at most two. Normalize
the now-connected residual graph once more to
reach $H(V,c)$ with root $a$. The same small-cache
root swap permits any prescribed labelled order.

The zero-count improvement obligation, applied at
most $n-1$ times, supplies a universal zero anchor
from every initial state. Intermediate counts need
not be monotone: strict improvement at the finitely
many macro endpoints suffices. The preceding
construction gives a finite path to a fixed atom.
Length one and length two supply the induction
bases by Theorems~\ref{thm:thesisone}
and~\ref{thm:thesisgeneraltwo}.

At each fixed $n,c,\ell$, the uniform finite-age
reset family is finite. Select one atom path from
each family member, maximize their finite lengths,
minimize their positive probabilities, and pad at
the empty-root self-loop. Combined with the uniform
reset, this gives a common-atom minorization.
Lemma~\ref{lem:atom} gives unique stationarity and
uniform geometric convergence. Label symmetry and
relabelled full-cache atoms give conditional
uniform samples and positive availability.
\end{proof}

There is also a counting limit to the current
shared-anchor approach. In
Lemma~\ref{lem:threecooperativecallback},
$P\cap(I\cup\{a\})=\varnothing$ and $|P|=\ell-2$,
so necessarily
\[
 Z_a=|I|\le n-\ell+1.
\]
Thus, for $\ell\ge4$, that lemma cannot handle the
last $\ell-3$ missing zero copies. The full-shuffle
positive-vector lemma has the analogous restriction
$Z_a\le n-c+1$. A general proof needs a separate
endgame or a progress argument without these shared
additional identities. Establishing their reachability
alone would not remove this counting obstruction.

\subsection{Combining the lemmas: a complete dense length-three theorem}

The next argument permits a zero copy to become positive temporarily
and restores it immediately afterwards. This additional freedom closes
length three at every capacity $c\ge n-2$. It does not impose any
new algorithmic rule or initial cache-size assumption.

\begin{lemma}[Dense cooperative owner service at length three]
\label{lem:dense3callback}
Let $\ell=3$ and $c\ge\max(3,n-2)$. Fix owner $a$ and its zero-holder
set $I$. If $a$ caches $p\notin I\cup\{a\}$, a finite
positive-probability path strictly increases $Z_a$ and does not
decrease $N_a$. Every original zero is present again at the endpoint.
At most one original zero is absent at any intermediate instant,
and the next owner exchange restores it.
If $p$ initially lacks identity $a$ altogether, the path can additionally
be chosen to increase $N_a$ by at least one.
\end{lemma}
\begin{proof}
At these capacities no surviving active resident is ever evicted:
after mandatory removal from a cache of size $d$, there are
$c-d+1\ge n-d-1$ vacancies, at least the number of possible new valid
reply identities. This is the counting argument of
Lemma~\ref{lem:densenoevict}, independent of shuffle length.

First observe a two-target service rule. If $a$ caches two distinct
peers $p,r\notin I\cup\{a\}$, let only $a$ initiate and send both
as extras before the first outside-$I$ contact.
Before that contact, every responder has $a^0$.
Include $a$ in its first reply. That responder receives and protects
the entire set $\{a,p,r\}$, and on repeats replies with exactly this set.
A first visit grows the active cache by at most one; a repeat shrinks
it because the self-$a$ reply is filtered and both targets are known.
There are finitely many first visits and the targets cannot be evicted.
An outside-$I$ contact therefore occurs and gains a zero, preserving
all earlier zeros and all identity-$a$ holders.

For a single target $p$, inspect an actual oldest owner head $Q$.
If it is outside $I$, the owner contacts it and finishes.
Otherwise freeze the owner and prepare $Q$ as follows.
If $p\notin C_Q$, or $|C_Q|\ge4$, its reply can contain $a$ and
exclude $p$; no preparation is needed.
In the other cases $C_Q$ is $\{a^0,p^m\}$ or
$\{a^0,p^m,R^b\}$.

If the residual $R$ exists and lies outside $I$, contact $Q$
from the owner with reply $\{a,p,R\}$. Both outside targets are now
cached, so the two-target rule finishes.
Otherwise, if $R\in I$ and $b\ge m$, activate $Q\to R$ with extras
$a,p$. Require $a$ in the reply, and $p,Q$ whenever they are known.
These requirements fit in three entries. This keeps $Q$'s anchor
zero, keeps its target $p$, and gives $R$ the protected identities
$a,p,Q$. Repeat this preparation until its stopping condition changes.
If the cache grows to size at least four, use the first case.
Every repeat responder already knows all three identities and
replies with exactly $\{a,p,Q\}$, leaving $C_Q=\{a^0,p^{m'}\}$.
Thus there are at most $|I|$ such internal preparation events.
If the reply produces a different outside residual, use the two-target
case instead.

In the remaining case $p$ is an actual oldest entry of $Q$.
Execute $Q\to p$. This removes $p$ from $Q$ because its owner cannot
return self-$p$. The dense no-eviction rule retains identity $a$ at $Q$,
although its age can now be positive. Immediately execute the promised
owner contact $a\to Q$; its oldest choice is unchanged because $a$
was frozen. Require $a$ in this reply, which automatically excludes $p$.
Fresh received $a^0$ restores $Q$'s zero.
The unprepared first case likewise uses an owner reply containing
$a$ and excluding $p$.

Call this preparation followed by one owner activation a macro.
Before a finishing branch, every macro restores all zeros in the
fixed set $I$, preserves all identity-$a$ holders, and increases the
owner's unrefreshed target age by exactly one. Write its initial
target age as $A$, and let $t$ count completed nonfinishing macros.
For $s\notin\{a,p\}$ let $M_s(t)$ be the global maximum stored age,
or zero if absent. Then
\[
 M_s(t+1)\le M_s(t)+1,\qquad
 M_Q(t+1)\le M_Q(t)
\]
when $Q$ is the inspected owner partner of this macro.
Indeed, every internal $Q\to R$ selects its sole residual entry:
apart from anchors $a,p$, it increments no retained or transmitted
descriptor. The optional final $Q\to p$ increments at most one
residual descriptor once. The owner was frozen and increments its
own copies before receiving these prepared copies, so no age exceeds
the preceding global maximum by more than one. Identity $Q$ is absent
from its own cache, appears freshly at age zero during preparation,
and is removed rather than forwarded in the owner activation.
Its global maximum therefore cannot grow at all.

Consequently $K_s(t)=M_s(t)-t$ is nonincreasing, and $K_Q$ drops by
at least one when $Q$ is selected. Owner-oldest selection requires
$M_Q(t)\ge A+t$, hence $K_Q(t)\ge A$. Each identity $Q$ can therefore
be selected at most $(M_Q(0)-A+1)_+$ times. Infinitely many
nonfinishing macros are impossible. This proves strict zero gain.

Finally suppose $p$ was initially a nonholder. If this path has
not yet created identity $a$ at $p$, then $p$ is still cached by
the owner, since it has not been selected and cannot be evicted.
Apply the proved zero-gain construction again with the enlarged
zero-holder set. Each application increases $Z_a$ and never loses
an identity-$a$ holder. Within at most $n-1$ applications, $p$
must receive $a$: otherwise a strictly increasing zero count
could never include $p$. Once created there, its copy cannot
be lost during these constructions because $p$ remains passive
and every subsequent request to it includes $a$.
Thus the endpoint also has at least one additional holder.
\end{proof}

\begin{corollary}[Dense length-three zeroification]
\label{cor:dense3zeroify}
At $\ell=3$ and $c\ge\max(3,n-2)$, every state has a finite path
to a state whose existing identity-$a$ copies are all zero, with
endpoint holder count at least its initial $N_a$.
Each positive-copy stage strictly increases $Z_a$.
Old zero copies are restored at the stage endpoints.
In any closed communicating class, if $k=\max Z_a$, then
$N_a\le k$ in every class state and all $k$ holders at a
maximizing state have age zero.
\end{corollary}
\begin{proof}
Freeze $a$ and all its zero holders and use
Theorem~\ref{thm:threefirstcontact} on a positive copy outside that set.
Contact with a zero holder refreshes an active positive copy and
increases $Z_a$ without decreasing $N_a$.
Contact with the owner removes one positive copy at the initiator
$p$ and puts fresh $p^0$ in the owner's cache. Apply the
nonholder version of Lemma~\ref{lem:dense3callback}: it restores the
lost holder count, gains a zero, and restores every previous zero.
Each complete stage increases the bounded integer $Z_a$.
Repeating while positive copies remain terminates within $n-1$
stages and never decreases the endpoint holder count.
The classwise bound follows by applying this construction inside
the closed class, exactly as in Corollary~\ref{cor:maxholders}.
\end{proof}

\begin{theorem}[Uniform sampling for dense three-entry shuffles]
\label{thm:dense3uniform}
For $\ell=3$, every $n\ge4$, and every allowed common capacity
$c\ge n-2$, the adopted kernel has a unique stationary law
and converges to it uniformly geometrically in total variation.
Every local request size $1\le b\le c$ has positive stationary
availability and a uniform limiting ordered sample conditional
on availability. Initially underfull caches and arbitrary finite
integer ages are allowed.
\end{theorem}
\begin{proof}
In a closed class fix $a$ and maximize $Z_a=k$.
The preceding corollary gives $N_a\le k$ throughout the class;
at a maximizing state all holders form the zero set $I$.
Also $k\ge1$, since otherwise weak connectivity forces a
nonempty owner cache and an owner activation creates a zero.
The owner cannot cache a nonholder by
Lemma~\ref{lem:dense3callback}, so $C_a\subseteq I$.

There is no edge from a holder $h\in I$ to a nonholder $x$.
Otherwise keep only $h$ active and, before a nonholder contact,
send $a,x$ as its two extras. At the terminal contact keep
$a$ as an extra and fill any remaining extra slot legally.
Before a nonholder response, require $a$ in every
reply and $h,x$ when known. First responders learn all three;
repeats reply with exactly $\{a,h,x\}$ and shrink the active
cache. First visits grow it by at most one, so this finite
first-visit/repeat argument eventually reaches a nonholder.
The owner $a$ need never be selected: its refreshed zero
anchor is no older than retained $x$, so choose a nonanchor
in a tie. Dense no-eviction keeps $a$ at $h$, and the
nonholder receives it, giving $N_a\ge k+1$, a contradiction.
Thus $I$ is outward closed except for its owner anchor.

Reserve $a^0$ and normalize residual components of $I$ using
the general length-two common-atom theorem.
For each owner-oldest head $Q$, normalize its component
with root $Q$, giving $C_Q=\{a^0\}$, then execute $a\to Q$.
Its self-$a$ reply deletes that owner head without adding
another. At $Q$, fresh $a$ and at most two extras fit in
$c\ge3$ and preserve its zero. Repeating empties the owner.
All resident identities of $I$ still lie in $I\cup\{a\}$.

Suppose $J=V\setminus(I\cup\{a\})$ is nonempty.
Weak connectivity now supplies an edge from $J$ to $I$:
the owner is empty, $I$ has no outward edge to $J$,
and no nonholder in $J$ can cache $a$.
Freeze $I$ and $a$ and use Theorem~\ref{thm:threefirstcontact}
until an outside initiator $P$ is about to contact some $Q\in I$.
No outside peer can acquire $a$ before this boundary.
Pause, freeze the outside, normalize $Q$'s residual component
with root $Q$, and resume the pending exchange.
The oldest choice of $P$ is unchanged and $C_Q=\{a^0\}$.
If $c\ge4$, its old anchor plus at most three incoming
entries fit. Both endpoints now have $a^0$, increasing $Z_a$.
For $c=3,n=4$, a nonholder can cache at most $n-2=2$
identities, so the same argument applies with a packet of
size at most two.

Only $c=3,n=5$ remains. Among maximizing states with
empty owner, choose one minimizing the number $d$ of
residual weak components on $I$. Run the same paused
boundary construction; residual normalization preserves $d$.
Packets of size at most two again give a zero gain.
Otherwise $C_P=\{Q,R,S\}$ and the exact boundary endpoint is
\[
 C_P'=\{a^0,R^u,S^v\},\qquad
 C_Q'=\{P^0,R^u,S^v\}.
\]
Incoming protection forces $Q$'s old anchor out; dense
no-eviction retains both extras at $P$. The new zero set
is $I'=(I\setminus\{Q\})\cup\{P\}$, still of size $k$,
and the owner is still empty.
If $R$ or $S$ was in $J$, new holder $P$ has an outward
edge to a nonholder, contradicting the established closure.
Thus both lie in $I\setminus\{Q\}$.
If $Q$'s normalized component had another peer, some
remaining old holder would still cache its root $Q$,
which is now a nonholder; this is the same contradiction.
Therefore that component was a singleton.
Removing it and adding $P$ attached to $R,S$ in the
remaining components strictly decreases $d$, contradicting
its minimality. This rules out $J$ also in the final case.

Every closed class consequently contains a universal zero
anchor. Proposition~\ref{prop:threezeroanchor} takes all
classes to the same labelled canonical atom.
The uniform finite-age reset and empty-root self-loop
give a uniform common-atom minorization as in
Theorem~\ref{thm:thesisgeneraltwo}.
Lemma~\ref{lem:atom}, label symmetry, and relabelled full-cache
atoms give the stated convergence, availability, and uniform samples.
\end{proof}

\begin{corollary}[The next induction step at dense length four]
\label{cor:densefourstructure}
Let $\ell=4$, $n\ge5$, and $\max(4,n-2)\le c\le n-1$.
The holder-only uniform first-contact conclusion of
Theorem~\ref{thm:threefirstcontact} holds at this length,
with finite horizon and positive probability bound independent
of initial ages. In addition, a universal zero anchor has
a finite path to any prescribed labelled canonical state.
These two conclusions do not assert full length-four sampling
from arbitrary states.
\end{corollary}
\begin{proof}
In the minimum-age-raising part of the first-contact proof,
reserve one anchor and replace the reduced length-two kernel
by length three on $n-1$ identities, at capacity $c-1$.
Its frozen-cut first-contact theorem is already valid at
every capacity. For a closed residual component of size $k$,
the effective capacity
\[
 d_k=\min(c-1,k-1)\ge k-2
\]
and length $e_k=\min(3,k-1)$ give a uniform common atom by
Theorem~\ref{thm:dense3uniform} when $k\ge4$; smaller components
use the length-one or length-two theorem.
Thus all minimum-raising phases have uniform finite horizons.
The final young-count/holder-maximum routing calculation is
unchanged: at length four there is still room for both the
old anchor and the required young difference extra.
The buffer bound becomes
$|D_q|+|Q\cap U|\le\min(4,|Q|)$, by the same all-length
counting lemma.

For a universal zero anchor, remove it and use the same dense
length-three theorem on every residual component, with the
effective parameters just checked. Owner emptying, the
four-exchange component merges, and root changes in
Proposition~\ref{prop:threezeroanchor} then apply unchanged.
\end{proof}

\begin{proposition}[The last zero at length three without density]
\label{prop:nearcomplete3}
Let $\ell=3$ and $4\le c\le n-1$, with no density restriction.
Suppose $p$ is the only nonowner without $a^0$, and the owner
caches $p$. A finite positive-probability path reaches a
universal zero anchor, restoring all old zeros.
The cache of $p$ may be empty, and the other caches and ages
are arbitrary subject to the stated zero copies.
\end{proposition}
\begin{proof}
Use the macro construction in Lemma~\ref{lem:dense3callback}.
There cannot be a second outside-$I$ residual $R\ne p$,
so its two-target branch is absent.
The remaining uses of density have the following local
replacements. An internal preparation starts from
$\{a^0,p^m,R^b\}$ and receives $a^0$ in its reply.
At most two new entries fit its at least two mandatory
vacancies at $c\ge4$, so $p$ survives.
In the preparation $Q\to p$, a two-entry $Q$ has room
for the entire reply. A three-entry $Q$ can need at
most one victim. If $a$ is received it is protected;
otherwise, whenever a victim is needed, its sole
residual $R$ is not received and can be deleted
instead of $a$. Thus $Q$ retains identity $a$.
In the owner exchange, Lemma~\ref{lem:ownertargetretain}
retains target $p$ using the self-$a$ reply.

All old zeros are restored at macro endpoints.
The same macro-age inequalities hold: internal preparations
do not increase nonanchor maxima, at most one residual
is incremented in the optional $Q\to p$, and the inspected
owner head's maximum cannot grow. The finite $K_s$ budget
therefore forces a contact outside $I$. That sole outside
peer is $p$, whose fresh zero completes the anchor.
\end{proof}

\subsection{Two further endpoint combinations at arbitrary lengths}

\begin{lemma}[An underfull zero responder can retain its anchor]
\label{lem:underfullzeroreply}
Let $2\le\ell<c$. If a nonzero-$a$ initiator $p\ne a$
contacts an underfull zero holder $q$, there is a legal
reply/victim choice that creates $a^0$ at $p$ and keeps
$a^0$ at $q$. In particular, if $p$ is the only nonowner
without $a^0$, every old zero holder is underfull, and
$C_p\ne\varnothing$, a universal zero anchor is reachable
in at most two exchanges, preserving all old zeros.
\end{lemma}
\begin{proof}
Choose a reply containing $a$. If $a$ is in the incoming
request $U$, its old zero is protected by minimum merging.
Otherwise write $b=|C_q|<c$, $h=|U\cap C_q|$, and
$r=(b+|U|-h-c)_+$ for its required victims.
There are $b-h-1$ eligible non-$a$ residents and
\[
 r\le b-h-1,\qquad r\le\min(\ell,b)-1,
\]
using $|U|\le\ell<c$ and $b<c$.
Choose these $r$ victims, put them and $a$ in the reply,
and fill to its prescribed size. This retains the old
zero and gives the active peer a protected zero.

For the second claim activate $p$. A zero-holder oldest
partner finishes by the first claim. Otherwise select
owner $a$, requiring an $I$ head in its reply whenever
one exists, where $I=V\setminus\{a,p\}$.
After mandatory removal, any remaining $p$ heads belong
to $I$, and its next exchange finishes.
If $p$ empties, the owner reply had no usable $I$ head,
so its pre-cache was empty or $\{p\}$. No usable reply
means no active victims; hence $p$'s pre-cache was
exactly $\{a\}$. The owner is now $\{p^0\}$, and
its contact with empty $p$ finishes instead.
\end{proof}

\begin{lemma}[A three-exchange fresh-copy gain at every longer length]
\label{lem:emptytargetthree}
Let $3\le\ell\le c$, let $I$ be the zero holders of $a$,
and let $p\notin I\cup\{a\}$ have empty cache.
Suppose an actual oldest owner head $Q$ belongs to $I$,
$p$ is an actual oldest head of $Q$, $|C_Q|\ge3$, and
$\operatorname{ids}(C_Q)\subseteq\{a,p\}\cup I$.
Then three exchanges increase $Z_a$ by one and restore
every old zero. No common-age or occupancy condition
is needed beyond this local configuration.
\end{lemma}
\begin{proof}
Execute $Q\to p$, choosing $a$ and some other
$R\in I\setminus\{Q\}$ as extras. All other required
extras also belong to $I$. The empty reply leaves
$a^1$ at $Q$; $p$ receives $Q^0,a^1$ and at least
one other descriptor of age at least one.
Immediately execute $a\to Q$, whose inspected oldest
choice is unchanged, restoring $a^0$ at $Q$.
The cache of $p$ has not changed since its passive receipt.
Its actual oldest can be selected among the forwarded
old zero holders $R$: their ages are at least one,
while $Q$ has age zero and $a$ has age one.
Execute $p\to R$ with its positive $a$ as extra and
$a^0$ in the reply. The received positive copy protects
$R$'s old zero by minimum merging, and $p$ acquires
a protected zero. Every other old zero is unchanged.
\end{proof}

These local configurations are conditional: their universal
accessibility has not been established. In contrast,
Theorem~\ref{thm:dense3uniform} proves its entire parameter
regime from arbitrary allowed initial states.

The singleton-mixture conclusion also extends to this adopted
convention: if each activation independently uses
$\widehat P_1$ with probability $\varepsilon>0$ and any
fixed $\widehat P_\ell$ otherwise, Proposition~\ref{prop:thesisunion}
keeps the connected state space closed. For $c\ge2$,
Theorem~\ref{thm:thesisone} gives the same common atom
with minorization $\varepsilon^T\gamma$.
Lemma~\ref{lem:atom} and Theorem~\ref{thm:stationary}
prove full convergence and conditional local uniformity.
This remains a variable-length protocol, distinct from the
fixed-length target.

\section{Established results and remaining problem}

The shuffle-length-one case is now proved for every population
size and cache capacity: Theorems~\ref{thm:mandatoryone}
and~\ref{thm:generalone} establish full-state convergence from
arbitrary finite ages, unique stationarity, positive availability,
and conditional local uniformity.
The latter proof supplies the missing common-atom construction:
source separation, a reserved callback, and induction connect
every zero-age state to the same labelled canonical overlay.
At capacity one the stationary law is explicit; at higher
capacities symmetry identifies the local conditional law without
requiring an explicit formula for the global stationary law.
Theorem~\ref{thm:three} additionally handles $n=3,c=\ell=2$.
The earlier retention theorem concerns a different convention.

Shuffle length two now has a uniform age bound for every
capacity: Theorem~\ref{thm:tworeset} makes a finite family
accessible with a uniform probability, independently of initial
ages. Lemma~\ref{lem:twoconnect} supplies connectivity, and
Proposition~\ref{prop:tworegimes} supplies stationary regimes
and continuous-time convergence to their mixture.
A common-regime or classwise symmetry argument is still needed
for general conditional uniformity. The activation-count version
also needs an aperiodicity argument.
Corollary~\ref{cor:twotails} additionally supplies exponential
stationary age tails and finite exponential moments.
At maximal capacity, Corollary~\ref{cor:twocomplete} proves
almost-sure finite completion of the underlying graph, with
an explicit uniform geometric hitting bound.

For every shuffle length, Theorem~\ref{thm:lineages} proves
age-dependent accessibility of erasing all initial descriptor
lineages. This alone cannot establish uniform sampling:
Theorem~\ref{thm:largerobstruction} disproves that conclusion
for the stated uniform-victim completion at every fixed
$3\le\ell\le c$, for sufficiently large populations and
weakly connected initial states. This includes proper shuffles,
and is a conditional-marginal result, rather than only a
pathwise disconnection example.
Corollary~\ref{cor:strongobstruction} strengthens the obstruction
to strongly connected full-network initial states, still with
full caches and zero ages. Its gateway argument controls the
probability of the isolation prefix independently of population
size; initial strong connectivity does not remove the failure.

Received-identity protection without age refresh supplies an explicit alternative
with preserved connectivity. Its singleton kernel is unchanged,
and its two-entry age-reset construction is still valid.
Theorem~\ref{thm:protectedmixture} fully proves convergence
for the protected protocol that occasionally performs singleton
shuffles. For fixed protected $\ell=2$, the remaining issue is
the stationary regime, rather than age escape. For fixed
protected $\ell\ge3$, age recurrence and a common regime remain
to be established. These are the meaningful remaining positive
proof targets; the failed completion cannot be covered by a
blanket convergence theorem.

The adopted target now additionally uses minimum-age duplicate
merging, following the thesis convention of Section~\ref{sec:thesis}.
For that kernel, Proposition~\ref{prop:thesisunion} establishes
lossless identity-union exchange and connectivity at every
shuffle length. Theorem~\ref{thm:thesisone} transfers singleton
convergence, and Theorem~\ref{thm:thesistwo} repairs the
two-entry age-reset argument using an initial age gap.
Minimum-age refresh can change future partner choices, so
the kernels and stationary age laws must not be identified
with their no-refresh counterparts.
Theorem~\ref{thm:thesisgeneralreset} now establishes uniform
finite-age recurrence for every fixed $\ell\ge2$ through a
lexicographic color potential. Corollary~\ref{cor:thesisgeneralregimes}
gives finitely many stationary regimes, uniform geometric
continuous-time convergence to their mixture, and exponential
stationary age tails.
Theorem~\ref{thm:thesisc2} now closes the full sampling
proof for $c=\ell=2$ and arbitrary $n$: an extremal
fresh-copy argument makes an anchor universal in every
closed class, which supplies a uniform common atom.
Theorem~\ref{thm:thesishalfdense} also closes shuffle
length two for every $c\ge3$ with $n\le2c$.
It uses holder-count-preserving zeroification, an
empty-owner double callback, and a count of full caches:
a transfer that fails to create an additional holder
must eventually increase this bounded integer.
Theorem~\ref{thm:thesisdense} retains the simpler
nearly-full-cache argument for $c\ge n-2$.
Theorem~\ref{thm:thesisgeneraltwo} removes the capacity
restriction completely. With only one current holder
initiating, all other original zero holders remain
passive. A global young-entry count never decreases
and strictly increases at each zero-holder contact.
On a constant-count segment, only migrations remain:
the tracked anchor ages by one each event, whereas
the normalized maximum age of every selected identity
strictly drops. Such a segment is finite. This proves
outward spread for every $c\ge3$, and completes full
uniform sampling at $\ell=2$ for every allowed capacity
and population size.
The remaining fixed-length task is a common regime or
sufficient classwise symmetry for $\ell\ge3$,
together with availability and, for the
embedded event chain, aperiodicity.
The earlier disconnection counterexamples
do not apply to this adopted target.
For $\ell=2$, Theorem~\ref{thm:firstcontact} and
Corollary~\ref{cor:repeatedcontact} now rule out permanent
communication isolation across any fixed cut, with a uniform
geometric waiting bound. Proposition~\ref{prop:rootreachability}
also makes every peer reachable from a prescribed root.
Neither statement by itself proves a common stationary
regime. Corollary~\ref{cor:anchorcriterion} identifies a
concrete sufficient remaining step: reach one identity
stored by every other peer. Its ages and the owner's
cache need not be prescribed in advance. That accessibility
criterion is now established at every cache capacity
for shuffle length two. This does not settle larger
shuffle lengths.
For $\ell=3$, Theorem~\ref{thm:threefirstcontact} now
rules out permanent cut isolation for every allowed
population and capacity, using only descriptor-holder
initiators. Proposition~\ref{prop:threezeroanchor}
reduces a common-atom proof to accessibility
of a universal zero-age anchor.
Theorem~\ref{thm:dense3uniform} now completes the
length-three sampling proof for every $c\ge n-2$.
Its cooperative callback temporarily lets an inspected
zero holder contact a target, then restores its zero
by the promised owner contact. A maximum-age budget
on these complete macros forces fresh-copy progress.
Zeroification, outward closure, and a minimum residual
component argument supply a universal zero anchor.
Smaller length-three capacities remain open.
For arbitrary larger lengths, the corresponding callback
and saturation construction remain missing.
Corollary~\ref{cor:densefourstructure} carries the
first-contact and universal-anchor reduction to length
four at $c\ge n-2$; it leaves universal-anchor accessibility
as a separate obligation.
The owner-only alternative is explicitly impossible
in the family of Proposition~\ref{prop:owneronlyobstruction};
the original all-peer scheduler can leave that
restricted family, as the three-event repair shows.
The cooperative callback lemmas establish progress
when old zero holders share suitable additional
descriptors: zero ages at arbitrary shuffle length,
or common integer ages at full shuffle. These are
conditions on a state reached during a proof, not
additional assumptions on the protocol or initial
network. They do not establish generic access to
the required shared descriptors.

The scheduling and sampling target cannot be omitted.
Theorem~\ref{thm:periodic} rules out a theorem covering every fixed
periodic phase schedule at $c=\ell=1$.
The retention-specific frozen states in Theorem~\ref{thm:twins}
do not refute the mandatory-removal target.
Proposition~\ref{prop:availability} concerns unconditional
full-capacity requests and does not refute conditional uniformity.
The present note does not establish a fixed-length convergence
theorem for every connectivity-preserving CYCLON completion,
nor an efficient or polylogarithmic mixing bound.

\section{Finite verification}
The accompanying \texttt{verify.py} checks the periodic trace including
ages, every possible initiating/subset choice in a small frozen-twin
example, the zero-age absorption certificate, the frozen construction's connectivity, the availability
argument on small cache states, and the capacity-one identity chain
exhaustively for $2\le n\le6$. It checks inverse moves, communicating
classes, stationary local counts, and Hamiltonian adjacent swaps.
It covers 14,831 connected identity states, 87,858 marked transitions,
and 564 Hamiltonian swap words.

The additional checker \texttt{verify\_reset.py}, using
\texttt{cyclon3.py}, both in \texttt{mandatory-removal/},
follows the constructive
three-peer reset without truncating the ages reached along its
paths. It exhaustively checks all 46,550 weakly connected initial
states whose stored ages belong to $\{0,1,2,3,4\}$.
Every constructed transition is checked against the exact
oldest-first kernel; every path reaches $Z(1,2,3)$ within the
proved nineteen-exchange limit. The longest observed path in this
finite check has fourteen exchanges. The symbolic proof above,
rather than this finite enumeration, establishes the theorem for
arbitrarily large finite integer ages. All these checks corroborate
the proofs; none replaces a missing general convergence argument.

The checker \texttt{verify\_capacity\_one.py} in the same directory
checks all weakly connected capacity-one identity states through
$n=6$, including empty caches, cycle reduction, tree closure,
inverse moves, tree irreducibility, uniform local counts, and
the Hamiltonian-path swap gadget, including swaps next to the
empty root.

The checker \texttt{verify\_singleton\_general.py} checks the
entire arbitrary-capacity induction with exact ages, including
every oldest choice, clearing bound, protected cache entry,
unknown-source invariant, induced connectivity requirement, and
zero-age block boundary. It passes 1,102 exhaustive small
forward states and 264 random aged initializations through
$n=10$. The largest constructed trajectory in those checks has
8,472 events. An independent implementation additionally passed
6,030 exhaustive connected zero-age states through $n=4$ and
1,000 random arbitrary-age initializations through $n=9$,
with ages up to 9,999.

The checker \texttt{verify\_lineage\_erasure.py} checks 3,300
initial states covering every parameter combination through
$n=10$, including empty and underfull caches and disconnected
states. It checks legal buffers, replacement victims, numerical
routing bounds, immune-peer preservation, and complete initial
lineage extinction. Its runs exercise 307 forced migrations
and have maximum observed length 2,237 events.
These finite and random checks support the symbolic proofs;
the all-parameter lineage result remains an accessibility
theorem, not a general mixing theorem.

The checker \texttt{verify\_capacity\_two\_reset.py} checks the
sharper two-entry-cache reset on 2,240 initial states through
$n=30$, with ages up to $10^6$ and 16,238 forced migrations.
The checker \texttt{verify\_shuffle\_two\_reset.py} checks the
arbitrary-capacity $\ell=2$ construction on 6,840 states
through $n=20$, every supported capacity, and ages up to $10^6$.
It checks the auxiliary fresh-count inequality at every
nonterminal routing event, pair retention, immunity, migration
and phase bounds, and final extinction. It exercises 27,395
forced migrations; the largest observed trajectory has
1,015 exchanges.

The checker \texttt{verify\_large\_shuffle\_obstruction.py}
checks 39 parameter choices, including proper shuffles.
It verifies the exact cache and buffer sizes, legal ages and
victims, core connectivity and outside closure, the permanent
isolated set, exact rational event probabilities, and the
strict conditional-total-variation threshold.
It also checks the gateway modification's strong connectivity,
full-cache zero-age initialization, the same isolation
certificate, and the stronger theorem's exact probability
lower bound and population threshold. The no-touch estimate
itself follows from the rotational-symmetry and Poisson
counting argument in Corollary~\ref{cor:strongobstruction}.
The general negative statement follows symbolically from
Theorem~\ref{thm:largerobstruction}; these finite checks do not
simulate the enormous sufficient population thresholds.

The checker \texttt{verify\_thesis\_two\_reset.py} checks the
minimum-age age-gap construction on 6,840 deliberately
nonvacuous initial states through $n=20$, every supported
capacity, and four age patterns. Initial ages scale with the
reset horizon and can exceed $10^{12}L$, so the tested old
color is never vacuous. It checks received-identity-protected
victims, exact minimum-age merging, the gap separation,
immunity, young-count inequalities, and the final age bound.
Its runs exercise 12,508 forced migrations, 12,056 younger
old--old refreshes, and 33,222 early tracked conversions.
The longest observed path has 188 exchanges. These tests
support Theorem~\ref{thm:thesistwo}; they do not replace its
symbolic proof or establish a common stationary regime.

The checker \texttt{verify\_thesis\_general\_local.py}
exhaustively checks the all-length color controller on
1,099,391 admissible local colored states through $n=7$,
including every one of their 3,887,740 permitted controller
completions. It verifies feasibility, protected victims,
and strict lexicographic decrease. Another 190 cases cover
larger transfers through size nineteen and capacities
through twenty.
The independent checker
\texttt{verify\_thesis\_general\_reset.py} checks 11,400
aged initializations through $n=20$, every supported
capacity and fixed shuffle length, and 373,413 controlled
exchanges. It verifies actual ages, oldest-first choices,
minimum merging, the gap separation, and the final bound;
the longest observed construction takes 89 exchanges.
These checks support Theorem~\ref{thm:thesisgeneralreset}.
They establish no additional common-class or sampling claim.

The checker \texttt{verify\_thesis\_anchor.py} checks the
anchor reduction and canonical construction on 576 states
through $n=10$, including arbitrary anchor labels,
nonempty owner caches, disconnected residual components,
and residual ages zero or $10^6$.
It checks 1,829,148 actual protected minimum-age exchanges,
including 1,822,679 independently compared reduced
singleton exchanges and 5,140 component-merging exchanges.
It covers excess capacity on smaller residual components
and the residual capacity-one case.
These checks support Proposition~\ref{prop:anchorcanonical};
they do not establish that an anchor is always accessible.

The checker \texttt{verify\_thesis\_first\_contact.py}
checks the minimum-raising controller exhaustively on
28,036 small outside-cache states through $n=4$, every
capacity and frozen set, and outside ages in $\{0,1\}$.
It checks 35,374 controlled activations, reaching the
frozen set in 16,002 cases and raising the outside minimum
in the other 12,034. Another 1,365 random states through
$n=15$, with ages up to $10^6$, pass the same check.
An additional 455 cases start with the tracked minimum
above the earlier bounded-age reset level and check that the outside color-reset
construction contacts the frozen set instead of finishing
a bounded-age reset. The minimum-raising checks support
Theorem~\ref{thm:firstcontact}; its current polynomial
route replaces that earlier reset-or-hit proof.

The checker \texttt{verify\_thesis\_first\_contact\_fast.py}
checks the strengthened polynomial route, including
proof-only young-OR duplicate markings, on 9,100 cases
through $n=15$, all supported capacities and frozen-set
sizes. It checks 21,507 exchanges, including 108 forced
migrations and 2,296 completed minimum raises; the longest
observed witness has 41 events.
Its \texttt{--zeroify} run checks another 9,100 cases
and 32,008 exchanges, including the terminal merge and
owner callback. Every block preserves old zero copies,
strictly increases their count, and has endpoint holder
count at least its starting value. These checks support
Proposition~\ref{prop:zeroify}; the longest observed block
has 56 events.

The checker \texttt{verify\_thesis\_c2\_saturation.py}
checks the new saturation proof on 52,964 exhaustive
initial states through $n=5$ with ages in $\{0,1\}$,
covering 118,498 legal oldest-partner boundary witnesses.
Another 684 random boundary witnesses and 684 outward
infection gadgets cover populations through $n=20$,
arbitrary anchor labels, and ages up to $10^6$.
All five extremal branches are exercised, including
2,065 strict residual-component decreases. Normalization
executes the actual minimum-age singleton embedding,
rather than assigning its endpoint caches.
Separate checks cover 775 source-only witnesses, 579
positive-copy witnesses, and 864 empty-owner, initially
zero-free witnesses. In total 788,171 exact exchanges
pass oldest-first, buffer, protection, age, zero-count,
and component-count checks. The symbolic argument in
Theorem~\ref{thm:thesisc2}, rather than enumeration of
network sizes, establishes the arbitrary-population result.

The checker \texttt{verify\_thesis\_c2\_quantitative.py}
executes the complete quadratic controlled path from 1,008
connected arbitrary-age initializations with $3\le n\le30$,
including ages up to $10^6$. All 351,100 exchanges pass the
kernel checks, the per-event probability bound $1/(4n)$,
and the total horizon $16n^2$; the longest observed path
has 1,154 exchanges at $n=30$. Branch coverage includes
727 positive-copy blocks, 950 owner-boundary repairs,
all incoming-boundary gain cases, and ten isolated-component
reductions, each costing exactly one event. A separate
exhaustive check covers all 1,728 connected capacity-one
identity configurations on three through five peers for
the two-event full-cache availability construction.
These checks support Theorem~\ref{thm:c2explicitmix}
and Proposition~\ref{prop:c2availabilityquant}; they do
not establish a polynomial mixing time. The exact
three-peer benchmark above is independently reproduced by
\texttt{verify\_c2\_recurrent\_mixing.py}, using rational
transition probabilities and integer matrix powers.
The separate checker
\texttt{verify\_c2\_refresh\_counterexample.py} enumerates
all eight terminal paths of the three-peer random-permutation
example. It verifies both complete refreshing and the
nonuniform conditional law $(3/7,4/7)$ exactly, together
with the successful-source-cover conditional law $(3/5,2/5)$.
It also checks a reachable state whose expected old-copy drift
is positive. These tests distinguish temporal connectivity
and refreshing from a sufficient sampling-mixing argument.

The checker \texttt{verify\_c2\_ready\_drift.py} covers
14,260 exact endpoint cases on three through six identity
labels with ages zero through three. It checks preservation
of target-holder count, the one-unit readiness-loss bound,
the exceptional-clock charging rule, 394 readiness losses,
3,800 readiness creations, and 1,880 full-reply readiness
transfers. These checks support the structural case split
in Theorem~\ref{thm:c2readycontact}. The displayed
potential drift supplies the arbitrary-size expectation
and tail bounds analytically.

The exact checker
\texttt{verify\_c2\_three\_peer\_antisymmetric.py}
enumerates all $487$ connected positivity-projection states
and certifies closure, equivariance, the exact four-event
availability lower bound $5/9$, and the exact eight-event
antisymmetric contraction $34447/69984<1/2$.
Its $974$ large unequal-positive-age lifts additionally
check the analytic projection. All certificate inequalities
use integer or rational arithmetic. The reduction from
arbitrary numerical ages and the iteration of the certified
operator bound are proved in
Propositions~\ref{prop:c2threepositivity}
and~\ref{prop:c2threeordinarysample}; the calculation is
not a finite-age truncation or an arbitrary-$n$ claim.

The checker \texttt{verify\_c2\_age\_projection\_obstruction.py}
verifies the cap-plus-rank obstruction by exact two-event
query laws at $K=1,100,10^6$. The symbolic comparison in
Proposition~\ref{prop:c2caprankobstruction} applies above
every fixed absolute-age cap. It excludes only that
particular proposed projection.

The checker \texttt{verify\_thesis\_dense.py} checks
43,338 exhaustive local nearly-full exchanges with no
initiator-extra eviction and 468 contacts with a singleton
anchor responder. Another 1,060 random outward-contact
witnesses and 1,060 random singleton-responder contacts
cover populations through $n=30$ and ages up to $10^6$.
All 3,145 random controlled exchanges pass; the longest
boundary-contact prelude has 27 events. These checks
support the local steps of Theorem~\ref{thm:thesisdense};
the extremal closed-class proof establishes its generality.

The checker \texttt{verify\_thesis\_half\_dense.py}
executes the outward-gain construction of
Lemma~\ref{lem:halfdensegain} in 10,777 initializations
through $n=30$, including all capacities
$c\ge\max(3,\lceil n/2\rceil)$, random states, and
forced disjoint replies with ages up to $10^6$.
All 36,204 numerical-age exchanges pass. The run
includes 1,203 full-cache migrations, 6,920 underfull
migrations, 2,364 owner double callbacks, and an
explicit $n=6,c=3$ path with three consecutive
full-cache migrations before the gain. These checks
exercise the exceptional branches; the bounded
full-cache potential proves termination for arbitrary
allowed population and cache sizes.

The checker \texttt{verify\_thesis\_general\_two.py}
tests Lemma~\ref{lem:generaloutwardgain} in 20,136
initializations through $n=30$ and every capacity
$3\le c\le n-1$, including the formerly uncovered
$n>2c$ cases. All 74,801 exact numerical-age exchanges
pass the buffer, victim, oldest, color-OR and
young-count checks. The run includes 273 young-partner
migrations, 3,106 strict zero-holder resets, 16,443
old-partner migrations, 8,227 owner finishes, 6,905
underfull finishes and 5,004 direct gains. Global
maximum-age checks also verify that each identity's
maximum grows by at most one per controlled event,
and cannot grow at all when that identity is the
selected partner. These finite checks corroborate
the symbolic termination argument; they do not
replace it with population-size enumeration.

The checker \texttt{verify\_thesis\_holder\_routing.py}
tests the high-age final routing phase of
Theorem~\ref{thm:threefirstcontact} on 20,475 cases
through $n=15$, including all $3\le\ell\le c$.
Its 54,280 events include 33,804 strict young-count
increases and one strict holder-maximum increase on
a young-count plateau. These checks concern the
routing phase, not the preceding minimum-age raising
or full higher-shuffle convergence.

The checker \texttt{verify\_thesis\_owner\_cycle.py}
checks the exact obstruction in
Proposition~\ref{prop:owneronlyobstruction} for
$m=3,\ldots,30$: 8,120 invariant-family oldest
branches and all 434 initially full first branches.
At the proper-shuffle capacity $c=m+1$ it checks
another 8,120 invariant branches and 434 first branches,
with every passive cache remaining underfull.
It also checks the four-peer cycle and its cooperative
three-event repair. The symbolic family proves the
auxiliary obstruction at every $m\ge3$.

The checker \texttt{verify\_thesis\_cooperative\_owner.py}
checks the three-entry shared-zero callback through
$n=30$ at every $c\ge3$. Its arbitrary-residual suite
has 15,120 starts and 186,403 exchanges, including
27,535 outside-boundary blocks and 13,223 empty blocks.
It independently compares 130,525 residual singleton
exchanges with the full kernel. The earlier closed-residual
suite has another 15,120 starts and 647,065 exchanges.
The general-shuffle checker
\texttt{verify\_thesis\_cooperative\_general.py}
tests 660 cases through $n=12$, all $3\le\ell\le c$,
and 4,774 exact exchanges. These support
Lemma~\ref{lem:threecooperativecallback}, including
its strict decrease of owner-cached zero-holder heads.

The checker \texttt{verify\_thesis\_cooperative\_vector.py}
tests Lemma~\ref{lem:fullpositivecallback} in 540
full-shuffle cases through $n=12$, common ages up to
$10^6$, and 1,421 exact exchanges. Its 72 low-inside
stopping blocks exercise the integer-age two-head
argument. The checks preserve all shared ages and
old zeros and enforce the bound of two clearing
blocks and three owner activations. Higher-shuffle
uniform sampling is not inferred from these
conditional verification results.

The checker \texttt{verify\_thesis\_proper\_owner.py}
tests Lemma~\ref{lem:largezerocallback} in 1,716 cases
through $n=15$, at every proper length
$3\le\ell<c\le n-1$, with initial target ages zero
and $10^6$. All 7,734 exact exchanges pass, including
6,018 internal contacts and 1,470 active replacement
choices protecting the designated target. The run
has 640 target terminals and 1,076 other terminals;
its longest path has 11 exchanges. It also checks
the unrefreshed target age, every identity's global
maximum-age budget, passive cache occupancy, old-zero
preservation, and holder counts. It tests cached-target
service under the stated occupancy condition, rather
than arbitrary reachability of that condition.

The checker \texttt{verify\_thesis\_dense\_three\_owner.py}
supports Lemma~\ref{lem:dense3callback} with 18,400 cases
through $n=15$ at both dense capacities, mixed ages up
to $10^6$, and 117,946 exact exchanges.
Its 30,950 internal preparations, 27,952 temporary
zero-loss/restoration events, 40,510 nonfinishing
macros, and 94 two-target exits exercise every
macro branch. The longest tested trajectory has
39 exchanges. Every macro checks the exact target age,
the global maximum-age inequalities, its selected-head
budget, and restored zeros; every event checks
nondecrease of the identity-holder count.
The same checker tests Proposition~\ref{prop:nearcomplete3}
without density in 27,200 cases through $n=20$
at every $4\le c\le n-1$: 350,417 exchanges,
120,506 internal preparations, 94,201 temporary
zero-loss/restoration events, and 108,510 macros.
The longest such trajectory has 52 exchanges.

The checker \texttt{verify\_thesis\_dense\_three\_boundary.py}
exhausts 144 normalized boundary configurations for
the $n=5,c=3$ endgame in
Theorem~\ref{thm:dense3uniform}, with ages zero,
one, and $10^6$. The exact protected merges give
108 zero gains, 27 outward-edge contradictions, and
nine strict component-count decreases.
These finite checks support the symbolic arbitrary-size
proof; they do not replace its extremal argument.

The checker \texttt{verify\_thesis\_underfull\_bootstrap.py}
supports Lemma~\ref{lem:underfullzeroreply} with
33,520 exhaustive local transitions through $n=8$
and 5,720 proper-shuffle near-complete starts
through $n=15$. Its 6,300 exchanges cover
5,140 direct gains, 572 owner-then-zero contacts,
and eight singleton-owner callbacks. It also
checks the excluded full-shuffle failure, where
a full novel request forces the old zero out.

The complete mathematical text is contained in this single
\texttt{main.tex}; it has no external TeX inputs or bibliography file.

\begin{thebibliography}{9}
\bibitem{peerswap}
R.~Guerraoui, A.-M.~Kermarrec, A.~Kucherenko, R.~Pinot, and M.~de Vos.
\emph{PEERSWAP: A Peer-Sampler with Randomness Guarantees}.
arXiv:2408.03829, 2024.
\url{https://arxiv.org/pdf/2408.03829}.
\bibitem{cyclon}
S.~Voulgaris, D.~Gavidia, and M.~van Steen.
\emph{CYCLON: Inexpensive Membership Management for Unstructured P2P
Overlays}. Journal of Network and Systems Management, 13:197--217,
2005. Published article, author-hosted copy:
\url{https://acropolis.aueb.gr/~spyros/www/papers/Cyclon.pdf}.
\bibitem{cyclonperformance}
F.~Bonnet, F.~Tronel, and S.~Voulgaris.
\emph{Brief Announcement: Performance Analysis of Cyclon,
an Inexpensive Membership Management for Unstructured P2P Overlays}.
Distributed Computing (DISC 2006), LNCS 4167, pp.~560--562, 2006.
\url{https://doi.org/10.1007/11864219_45}.
\bibitem{cyclonthesis}
S.~Voulgaris.
\emph{Epidemic-Based Self-Organization in Peer-to-Peer Systems}.
PhD thesis, Vrije Universiteit Amsterdam, 2006.
Chapter~2, especially Figure~2.3 on p.~19.
\url{https://acropolis.aueb.gr/~spyros/www/papers/Thesis-Voulgaris.pdf}.
\end{thebibliography}
\end{document}
